\documentclass[11pt,a4paper]{article}
\usepackage[T1]{fontenc}
\usepackage[utf8]{inputenc}
\usepackage{lmodern}
\usepackage[margin=27mm,headheight=15pt]{geometry}
\usepackage{amsmath,amssymb,amsthm,mathtools}
\usepackage{microtype}
\usepackage{booktabs,longtable,array}
\usepackage[dvipsnames]{xcolor}
\usepackage{xurl}
\usepackage{fancyhdr}
\usepackage{hyperref}
\hypersetup{colorlinks=true,linkcolor=MidnightBlue,urlcolor=MidnightBlue,
  citecolor=MidnightBlue,pdftitle={A computer-assisted construction of a c=8 Haagerup CFT},
  pdfauthor={Research manuscript},pdfsubject={Proof manuscript with explicit review obligations}}
\urlstyle{tt}
\setlength{\emergencystretch}{2em}
\setlength{\parskip}{0.35em}
\pagestyle{fancy}
\fancyhf{}
\fancyhead[L]{\small Central-charge-eight Haagerup construction}
\fancyhead[R]{\small Review draft}
\fancyfoot[C]{\thepage}
\newtheorem{theorem}{Theorem}[section]
\newtheorem{lemma}[theorem]{Lemma}
\newtheorem{proposition}[theorem]{Proposition}
\newtheorem{corollary}[theorem]{Corollary}
\newtheorem{claim}[theorem]{Claim}
\theoremstyle{definition}
\newtheorem{definition}[theorem]{Definition}
\newtheorem{remark}[theorem]{Remark}
\newtheorem{certificate}[theorem]{Certificate input}
\DeclareMathOperator{\Rep}{Rep}
\DeclareMathOperator{\Com}{Com}
\DeclareMathOperator{\FPdim}{FPdim}
\DeclareMathOperator{\Hom}{Hom}
\DeclareMathOperator{\End}{End}
\DeclareMathOperator{\Tr}{Tr}
\DeclareMathOperator{\wt}{wt}
\DeclareMathOperator{\Span}{span}
\DeclareMathOperator{\qdim}{qdim}
\DeclareMathOperator{\ch}{ch}
\DeclareMathOperator{\Aut}{Aut}
\renewcommand{\Vec}{\mathrm{Vec}}
\newcommand{\HH}{\mathcal H}
\newcommand{\CC}{\mathbb C}
\newcommand{\QQ}{\mathbb Q}
\newcommand{\ZZ}{\mathbb Z}
\newcommand{\one}{\mathbf1}
\newcommand{\artifact}[1]{\textnormal{\nolinkurl{#1}}}
% Breakable review markers require only standard LaTeX and xcolor.
% Their textual labels remain meaningful in a monochrome printout.
\newenvironment{reviewnote}[1]{\par\addvspace{8pt}\begingroup
  \color{Sepia}\noindent\rule{\linewidth}{0.7pt}\par\nobreak
  \noindent\textbf{REVIEW --- #1}\par\nobreak\small\noindent\ignorespaces
}{\par\noindent\rule{\linewidth}{0.7pt}\par\endgroup\addvspace{8pt}}
\newenvironment{criticalnote}[1]{\par\addvspace{8pt}\begingroup
  \color{BrickRed}\noindent\rule{\linewidth}{1pt}\par\nobreak
  \noindent\textbf{HIGH PRIORITY REVIEW --- #1}\par\nobreak\small\noindent\ignorespaces
}{\par\noindent\rule{\linewidth}{1pt}\par\endgroup\addvspace{8pt}}
\title{A computer-assisted construction of a\\[3pt]
  \texorpdfstring{\(c=8\)}{c=8} Haagerup CFT\\[5pt]
  \large Proof manuscript with explicit review obligations}
\author{Research manuscript}
\date{27 September 2026}

\begin{document}
\maketitle
\begin{abstract}
We describe an explicit conformal vertex subalgebra of the lattice VOA
\(V_{E_6\oplus A_2}\), generated by an \(A_{1,1}\otimes A_{1,39}\)
current algebra, the ambient stress tensor, and one three-term lattice field.
Exact finite certificates, together with a recursive strong-generation
argument, prove that its current space has dimension six. From the
stated embedded-field and finite-action inputs we deduce strong
rationality, actual lattice branching, and the twelve-component
Evans--Gannon character vector with parameter zero. Sector counting and
reverse reconstruction give strong unitarity and a completely rational
conformal net whose representation category is unitarily ribbon
equivalent to the Haagerup center. This yields a diagonal unitary full
net theory at $c_L=c_R=8$. Complete unitarity of the VOA and a chosen
$\HH_3$ holomorphic VOA extension remain conditional on one explicit
geometric annulus comparison. We include fresh exact certificate
replays and limited Lean formalizations, with their scopes distinguished
from the mathematical arguments and the remaining construction inputs.
\end{abstract}

\begin{criticalnote}{The remaining analytic hypothesis}
The net-sector comparison and linear module reconstruction, previously
labelled A2 and A3, are proved below from the stated algebraic inputs
and cited reconstruction theorems. The annulus comparison A1 remains
unproved: the specific geometric operators used by the fusion-positivity
theorem must be matched, including their domains and central lifts.
The conformal-net realization does not require A1. Complete VOA
unitarity and the unitary ribbon identification of $\Rep(U)$ do require
it. Exact field identities remain certificate inputs; the finite
replays and the narrowly scoped Lean proofs do not verify the entire
construction. This is a revised research draft, without external
expert validation.
\end{criticalnote}

\clearpage
\tableofcontents
\clearpage

\section{Introduction and statement of results}
\label{sec:statement}

The Haagerup fusion categories provide a test case for the realization of
finite tensor categories by chiral conformal field theory. Their centers
are modular, but modular data alone do not construct a vertex operator
algebra: the operator products, positivity, and representation theory
must be supplied separately. Evans and Gannon computed the relevant
modular data and the possible character vectors at central charge eight
\cite[Sections 2.4 and 5]{EG}. In particular, the character with six
weight-one states suggests a small current algebra inside a much larger
lattice algebra. This observation motivates the construction studied here.

We start with actual vectors in a positive-definite lattice VOA. The
ambient vertex operators fix all products and their associativity from
the outset. The first substantive problem is then to prove that these
vectors generate a proper subalgebra. Generating additional fields in
weights two and three does not settle that question: singular products
of higher fields could still create new currents. We address this issue
by a finite strong-generation criterion. Its input consists of exact
normal-product identities and the recursive provenance of the higher
generators. Its output controls every conformal weight.

There are three distinct stages. The first constructs the subalgebra and
proves properness from explicit finite data. The second uses positive
conformal vectors, branching, and fusion-category arguments to analyze
its representation theory. The third compares VOA modules with conformal
net sectors, proves the net realization, and isolates the remaining
condition for complete VOA unitarity. In particular,
the existence of a positive vacuum inner product is not substituted for
positivity of all fusion pairings.

Our presentation uses the normal-product conventions of De Sole--Kac
\cite[Section 2]{DSK}, separates character-theoretic information as in
Evans--Gannon \cite{EG}, and states analytic transport hypotheses in the
operator-theoretic setting of Carpi--Kawahigashi--Longo--Weiner \cite{CKLW}.
These works also guide the level of detail: finite input is recorded as
such, and each passage to an infinite-dimensional conclusion is made by
a separate lemma. The exact arithmetic certificates are accompanying
data, not substitutes for these implications.

\subsection{The candidate and its finite construction}

Let
\[
 B=V_{E_6\oplus A_2},\qquad H=V_{E_8},\qquad B\subset H
\]
be the standard index-three lattice extension. The notation \(A_{1,k}\)
means the simple affine VOA \(L_k(\mathfrak{sl}_2)\). Fix the embedding
\(\mathcal K=A_{1,1}\otimes A_{1,39}\subset B\) described in
Section~\ref{sec:construction}. Let \(u\in B_2\) be the field in
\eqref{eq:u}, and define the \emph{actual generated} vertex algebra
\begin{equation}\label{eq:U}
 U=\langle\mathcal K,\omega_B,u\rangle\subset B.
\end{equation}
All products are computed in this lattice VOA. Thus associativity is
inherited from the ambient algebra, rather than assumed for an abstract
presentation inferred from character coefficients.
Throughout the normal-space and mode formulas, \(K=\mathcal K_1\)
denotes only the six-dimensional current space, not the full affine VOA.

\begin{proposition}[Computational properness conclusion]\label{prop:mainproper}
Accept the exact lattice identities and provenance recorded in
Section~\ref{sec:ledger}. The algebra \(U\) is strongly generated by
actual fields of weights at most four. Its weight-one space is
\[
 U_1=\mathfrak{sl}_2\oplus\mathfrak{sl}_2,\qquad \dim U_1=6.
\]
Consequently \(U\subsetneq B\), since \(\dim B_1=86\).
\end{proposition}
The finite certificates and the all-weight implication are presented in
Sections~\ref{sec:finite}--\ref{sec:properness}. This is the most directly
checkable part of the construction.

\begin{corollary}[Initial graded pieces]\label{cor:initial-grades}
Under the same finite-certificate hypotheses, $U$ is a simple unitary
VOA of CFT type with the ambient conformal vector and central charge
eight. Moreover,
\[
 U_2=S_2,\qquad U_3=S_3,\qquad
 \ch_U(\tau)=q^{-1/3}(1+6q+120q^2+956q^3+O(q^4)).
\]
Here $S_2,S_3$ are the concrete spaces in
Certificate~\ref{cert:initial}. This assertion does not require a modular
character calculation.
\end{corollary}
The proof is given after the properness argument. In particular the
coefficient $6010$ in the proposed character below has a different
dependency from the first four coefficients.

We write $\HH_3$ for the unitary Haagerup fusion category with six
simple objects $\one,\alpha,\alpha^2,\rho,\alpha\rho,\alpha^2\rho$
and fusion rules
\[
 \alpha^3=\one,\qquad \alpha\rho=\rho\alpha^{-1},\qquad
 \rho^2=\one+\rho+\alpha\rho+\alpha^2\rho.
\]
Here $\operatorname{FPdim}(\rho)=(3+\sqrt{13})/2$.
The other unitary category with this fusion ring, denoted $\HH_2$,
is Morita equivalent to $\HH_3$. Consequently their Drinfeld centers
are equivalent even though the fusion categories are not. In the
statement below $Z$ denotes this center, not merely its modular matrices.

\begin{theorem}[Conformal-net realization]\label{thm:main-net}
Assume the exact finite construction inputs recorded in
Section~\ref{sec:ledger} and Certificate~\ref{cert:structural-inputs}.
The VOA $U$ is strongly rational and strongly unitary, with central
charge eight. Its completely rational smeared-field conformal net
$\mathcal N_U$ satisfies
\[
 \Rep(\mathcal N_U)\simeq Z(\HH_3)
\]
as unitary ribbon categories. The twelve-component character vector
is the $\gamma=0$ Evans--Gannon vector, and in particular
\begin{equation}
 \ch_U(\tau)=q^{-1/3}\bigl(1+6q+120q^2+956q^3+6010q^4+\cdots\bigr).
 \label{eq:vacchar}
\end{equation}
There is a diagonal unitary full conformal net with $c_L=c_R=8$
and partition function $\sum_{i=1}^{12}|\chi_i(\tau)|^2$.
\end{theorem}
The proof is assembled in Theorem~\ref{thm:positive-stress-rationality},
Proposition~\ref{prop:analytic-strong-unitarity},
Corollary~\ref{cor:analytic-net-haagerup}, and
Proposition~\ref{prop:analytic-diagonal}. The character argument is
given before the analytic comparison and does not presume it.

\begin{theorem}[Conditional VOA identification]\label{claim:main}
Under the same construction inputs, assume additionally the geometric
annulus comparison in Claim~\ref{clm:analytic-annulus}. Then $U$ is
completely unitary and
\begin{equation}\label{eq:maincat}
 \Rep(U)\simeq Z(\HH_3)
\end{equation}
as unitary ribbon categories. There is a unitary holomorphic extension
$H'\cong V_{E_8}$ whose category of all algebra modules is $\HH_3$;
its algebra-bimodule category is
$\HH_3\boxtimes\HH_3^{\mathrm{op}}$.
\end{theorem}

Here \(H'\) is obtained by transporting the canonical \(\HH_3\)
Lagrangian algebra. It need not coincide with the originally embedded
extension \(H/U\); the proposed argument identifies that original
all-module category only up to the alternatives \(\HH_2,\HH_3\).
The full-CFT defects in Theorem~\ref{claim:main} preserve \(U\), and need
not commute with the enlarged \(E_8\) chiral algebra.

\subsection{How the qualifications are used}
Amber \textbf{REVIEW} notes identify the boundary of exact construction
inputs and their replay coverage. Red \textbf{HIGH PRIORITY REVIEW}
notes isolate the unresolved geometric comparison A1 and explain its
consequences. A proved lemma is not made into an additional hypothesis
merely because its argument is new; its proof and external dependencies
are supplied for scrutiny. Conversely, the unproved comparison is
explicitly retained as a hypothesis wherever it is used.

\begin{center}
\begin{tabular}{@{}p{0.27\linewidth}p{0.65\linewidth}@{}}
\toprule
Part of the argument & Nature of its evidence \\
\midrule
Explicit fields and finite closure & Exact arithmetic over
\(\QQ(\sqrt{13},\sqrt{-3})\), saved recipes and full residuals.\\
Properness in all weights & Recursive algebraic theorem plus finite
certificates; no character upper-bound extrapolation.\\
Rationality and branching & Positive-stress and category arguments,
embedded-field and finite-top inputs, and cited external theorems.\\
Haagerup modular data & Rank exclusions, indicator and character
arguments printed below, with exact finite arithmetic checked separately.\\
Net realization and strong unitarity & Actual branching, linear
reconstruction, sector counting, and unitary classification.\\
Complete VOA unitarity and extension & Conditional on A1, then fusion
positivity, unitary classification, and extension theorems.\\
\bottomrule
\end{tabular}
\end{center}

\section{The explicit lattice construction}\label{sec:construction}
\subsection{Lattice, cocycle, and currents}
Use simple lattice vectors indexed \(0,\ldots,7\). The Gram matrix
has diagonal entries two and off-diagonal entries minus one on
\[
 (0,1),(1,2),(2,3),(3,4),(2,5),(6,7),
\]
with all remaining off-diagonal entries zero. Fix the cocycle
\begin{equation}\label{eq:cocycle}
 \epsilon(a,b)=(-1)^{\sum_i a_i b_i+
                      \sum_{i>j}a_i b_jG_{ij}}.
\end{equation}
Write $L=E_6\oplus A_2$ and $\mathfrak h=L\otimes_{\ZZ}\CC$.
The underlying vector space is
\[
 B=S\bigl(\mathfrak h\otimes t^{-1}\CC[t^{-1}]\bigr)
       \otimes\CC_\epsilon[L].
\]
A monomial $h_1(-n_1)\cdots h_r(-n_r)e^\alpha$ has weight
$(\alpha,\alpha)/2+\sum_i n_i$. In particular its weight is a
nonnegative integer. The exponent in \eqref{eq:cocycle} is bilinear
modulo two, and hence defines a normalized two-cocycle. Directly,
\[
 \frac{\epsilon(\alpha,\beta)}{\epsilon(\beta,\alpha)}
       =(-1)^{(\alpha,\beta)},\qquad
 \epsilon(\alpha,-\alpha)=(-1)^{(\alpha,\alpha)/2}.
\]
These identities fix the signs of all root brackets. For example,
$e^\alpha_0e^{-\alpha}=-\alpha(-1)\one$ for a root $\alpha$.
The exponential fields obey
\begin{equation}\label{eq:lattice-product}
 Y(e^\alpha,z)e^\beta
 =\epsilon(\alpha,\beta)z^{(\alpha,\beta)}
   \exp\!\left(\sum_{n>0}\frac{\alpha(-n)}{n}z^n\right)e^{\alpha+\beta}.
\end{equation}
Together with the Heisenberg commutation relations, this formula
specifies the exact mode engine used by the certificates.

Put $\alpha_i=(0,\ldots,1,\ldots,0)$ and
$\alpha_L=(-1,-2,-3,-2,-1,-2,0,0)$. With
\[
 (c_0,\ldots,c_7)=(5,8,9,8,5,0,2,2),\qquad
 h_R=\sum_i c_i\alpha_i,
\]
the two current triples, in the order $(H,E,F)$, are
\begin{align}
 (H_L,E_L,F_L)&=
   (\alpha_L(-1)\one,e^{\alpha_L},-e^{-\alpha_L}),\label{eq:leftcurrents}\\
 (H_R,E_R,F_R)&=
   \left(h_R(-1)\one,\sum_{c_i\ne0}e^{\alpha_i},
                    -\sum_i c_i e^{-\alpha_i}\right).
                    \label{eq:rightcurrents}
\end{align}
Current letters also denote their zero-mode actions when applied to
a state. Thus, for example, $F_RE_Rv$ means
$(F_R)_0(E_R)_0v$.

\begin{lemma}[The two affine levels]\label{lem:levels}
The triples \eqref{eq:leftcurrents}--\eqref{eq:rightcurrents} commute
and generate $A_{1,1}\otimes A_{1,39}$.
\end{lemma}
\begin{proof}
We have $(\alpha_L,\alpha_L)=2$ and
$(\alpha_L,\alpha_i)=0$ whenever $c_i\ne0$.
Also $(h_R,\alpha_i)=2$ for every $c_i\ne0$, and
$(h_R,h_R)=78$. Formula~\eqref{eq:lattice-product} gives
$[H_0,E_0]=2E_0$, $[H_0,F_0]=-2F_0$ and $[E_0,F_0]=H_0$
for both triples. The double poles are
\[
 (H_L)_1H_L=2\one,\quad (E_L)_1F_L=\one,\quad
 (H_R)_1H_R=78\one,\quad (E_R)_1F_R=39\one.
\]
Orthogonality gives all mixed singular products equal to zero.
These are the affine relations at levels one and thirty-nine.
The positive-definite lattice representation is integrable; equivalently
one may use the principal $\mathfrak{sl}_2$ embeddings in $A_5$ and
$A_2$, whose levels $35$ and $4$ add in the diagonal. Thus the
generated factors are the simple positive-integral-level affine VOAs.
Their tensor product is simple and embeds in $B$.
\end{proof}

\subsection{The defining field and the positive real form}
The exact coefficient convention uses unnormalized simple-root
Heisenberg oscillators and a torus gauge in which the principal raising
current is the sum of the corresponding simple-root vertices. The
lowering current is \(-\sum_i c_i e^{-\alpha_i}\), with
\[
 (c_0,\ldots,c_7)=(5,8,9,8,5,0,2,2).
\]
The torus rescales simple-root vertices by
\( (1/\sqrt5,1/\sqrt8,1/3,1/\sqrt8,1/\sqrt5,1,1/\sqrt2,1/\sqrt2)\).
In this convention define
\begin{align}
 u={}&e^{(0,0,0,0,0,-1,1,1)}\notag\\
 &+\frac{\sqrt{13}+3\sqrt{-3}}{20}
                   e^{(0,0,0,1,1,-1,0,0)}\notag\\
 &+\frac{\sqrt{13}-3\sqrt{-3}}{20}
                   e^{(1,1,0,0,0,-1,0,0)}.\label{eq:u}
\end{align}
The compact real form and PCT operation must be transported through
this gauge. Ordinary coefficient conjugation in the displayed gauge
is not a replacement for the ambient PCT map.

Here is the transport explicitly. Set
$\widetilde c=(5,8,9,8,5,1,2,2)$ and let $g$ be the lattice torus
automorphism which fixes the oscillators and multiplies $e^\alpha$ by
$\prod_i\widetilde c_i^{-\alpha_i/2}$. If $\theta$ is the ordinary
lattice PCT involution, the transported involution is
$\theta_g=g\theta g^{-1}$. For real lattice oscillators,
\begin{equation}\label{eq:pct}
 \theta_g\bigl(h_1(-n_1)\cdots h_r(-n_r)e^\alpha\bigr)
 =(-1)^r\Bigl(\prod_i\widetilde c_i^{\alpha_i}\Bigr)
       h_1(-n_1)\cdots h_r(-n_r)e^{-\alpha}.
\end{equation}
It acts antilinearly on coefficients, so fixes $\sqrt{13}$ and
negates $\sqrt{-3}$. The transported Hermitian form is
$\langle v,w\rangle_g=\langle g^{-1}v,g^{-1}w\rangle$ and is
positive definite. Formula~\eqref{eq:pct} squares to the identity.
The same torus element acts on $V_{E_8}$: the positive real character
on $L$ extends by its real linear logarithm to the common Cartan space.
We transport the form and PCT on $H$ by that element as well, so the
entire inclusion $U\subset B\subset H$ uses one unitary convention.
It preserves the current algebra, including the root factor even
though its two root vectors have different norms in this gauge.

All three momenta in \eqref{eq:u} have squared norm four and Cartan
weight $(1,13)$. Substitution in the lattice formula verifies that
$u$ is a highest vector of current type $(2,14)$, where the entries
now denote representation dimensions. In the original positive unitary
basis, $g^{-1}u$ is, up to a common scalar,
\[
 e^a+\zeta e^b+\overline\zeta e^c,\qquad
 \zeta=\frac{\sqrt{13}+3i\sqrt3}{\sqrt{40}},\qquad |\zeta|=1,
\]
with $a,b,c$ the three displayed momenta. Thus the exact construction
fixes a relative phase; there is no remaining continuous parameter.

The ambient and affine conformal vectors are fixed by these conventions:
\begin{align}
 \omega_B&=\frac12\sum_{i,j}(G^{-1})_{ij}
               \alpha_i(-1)\alpha_j(-1)\one,\label{eq:ambient-stress}\\
 \omega_{A_{1,k}}&=\frac{1}{2(k+2)}
       \left(\tfrac12 H_{-1}H+E_{-1}F+F_{-1}E\right).
       \label{eq:sugawara}
\end{align}
The conformal vector of $U$ is $\omega_B$, not the sum of the two
affine conformal vectors. Their central charges sum to
$1+117/41<8$.

\subsection{The constructed spaces in weights two and three}

\begin{certificate}[Actual initial spaces]\label{cert:initial}
The saved exact construction gives a PCT-stable 120-dimensional
space \(S_2\subset U_2\), containing \(u,\omega_B\), the affine
vacuum descendants, two neutral primitive fields, and eight charged
current multiplets. The current multiplet dimensions are
\[
 (1,5),(1,9),(1,9),(1,13),(2,2),(2,6),(2,8),(2,14).
\]
These contribute 96 vectors, in addition to the 22 affine vacuum
vectors and two neutral fields. The exact generated grade-three
space \(S_3\subset U_3\) has dimension 956. Its normal subspace
\(N_3=K_{-1}S_2+TS_2\) has dimension 634, and
\[
 S_3=N_3+\Span\{a_0b:a,b\in S_2\},\qquad \dim(S_3/N_3)=322.
\]
The spaces are current-stable. These are dimensions of constructed
finite spaces, not initial assertions that they exhaust \(U_2,U_3\).
\end{certificate}
The source records are
\artifact{hi_cft_20260923_c8_exact_fields.json},
\artifact{hi_cft_20260923_c8_exact_unitary_certificate.json}, and
\artifact{hi_cft_20260923_c8_minimal_generators.json}, together with
the exact grade-three checkpoint bound into the final ledger.

The minimal-generator record is logically prior to the closure records.
It starts with the six specified currents, $\omega_B$, and $u$, then
closes their weight-two span under current zero modes and $u_1$.
Every intermediate vector is therefore an actual $U$-vector. The final
span has dimension $120$ and contains each independently specified
highest vector and neutral conformal vector. The equality with the saved
space $S_2$ is checked only after this construction. This is the step
that relates the large certificate presentation to the short definition
\eqref{eq:U}; an ambient weight-two dimension count alone would not do so.

For reference the two neutral fields, denoted $e,f$, satisfy
\begin{equation}\label{eq:two-coset-stresses}
 e+f=\omega_B-\omega_{A_{1,1}}-\omega_{A_{1,39}},\qquad
 c_e=\frac{2400}{1271},\quad c_f=\frac{70}{31}.
\end{equation}
They commute as Virasoro fields. The exact current and stress
eigenvalues on the eight primitive highest vectors are as follows.
The two occurrences of $(1,9)$ are distinct vectors.
\begin{center}
\begin{tabular}{@{}ccc@{}}
\toprule
Current dimensions $(n,m)$ & $e_1$ eigenvalue & $f_1$ eigenvalue\\
\midrule
$(1,5)$ & $798/1271$ & $38/31$\\
$(1,9)$ & $200/1271$ & $42/31$\\
$(1,9)$ & $1512/1271$ & $10/31$\\
$(1,13)$ & $420/1271$ & $20/31$\\
$(2,2)$ & $192/1271$ & $49/31$\\
$(2,6)$ & $1830/1271$ & $3/31$\\
$(2,8)$ & $752/1271$ & $24/31$\\
$(2,14)$ & $180/1271$ & $13/31$\\
\bottomrule
\end{tabular}
\end{center}
Each row obeys
\[
 \frac{n^2-1}{12}+\frac{m^2-1}{164}+h_e+h_f=2.
\]
The affine vacuum space contributes $22$ dimensions in weight two:
$4$ from $A_{1,1}$, $9$ from $A_{1,39}$, and $3\cdot3$ mixed
current products. The level-one null multiplet removes five of the
nine universal weight-two states of its affine factor. Together with
the two singlets and the $96$ charged states, this gives $120$.
The exact normal-space and rank calculations certify these spaces
inside the lattice algebra, including their intersections.

\section{Finite closure and the all-weight criterion}\label{sec:finite}
Write \(Y(a,z)=\sum_{n\in\ZZ}a_nz^{-n-1}\) and \(T=L(-1)\).
A singular product has mode index \(n\geq0\). The zero mode \(a_0\)
is a derivation and commutes with \(T\).

\begin{definition}[Differential normal span]\label{def:normal-span}
For a homogeneous positive-weight subspace $S\subset V$, let
$\mathcal N(S)$ be the linear span of the vacuum and every expression
obtained from vectors of $S$ by translation and arbitrary parenthesized
$(-1)$ products. Let $\mathcal N_d(S)=\mathcal N(S)\cap V_d$.
This definition makes translation and normal multiplication closed
operations on $\mathcal N(S)$; closure under singular products is a
separate assertion. We use the usual convention that a set strongly
generates when its differential normal span is the whole algebra.
\end{definition}

\begin{lemma}[Filtered Wick reduction]\label{lem:filtered-wick}
Let $V$ be a nonnegatively integer-graded vertex algebra with
$V_0=\CC\one$, and let $S$ have strictly positive weights. Fix $w$.
Suppose every singular product of two homogeneous generators of total
weight less than $w$ belongs to $\mathcal N(S)$. Then the same is
true for two normal expressions of total weight less than $w$.
It is also true at total weight $w$ if at least one expression is a
derivative of a positive-weight expression or a normal product of
two positive-weight expressions.
\end{lemma}
\begin{proof}
Use induction on total input weight. A vacuum factor has zero singular
bracket. Sesquilinearity removes a derivative and lowers the total input
weight; the resulting powers of the formal variable and translations
preserve coefficientwise membership in $\mathcal N(S)$. Skew symmetry
allows a composite input to be placed on the right. The noncommutative
Wick identity is
\begin{equation}\label{eq:wick}
 [x_\lambda:y z:]=:[x_\lambda y]z:+:y[x_\lambda z]:
             +\int_0^\lambda [[x_\lambda y]_\mu z]\,d\mu.
\end{equation}
The first two brackets have smaller total input weight, since an
omitted factor has positive weight. Their coefficients therefore admit
homogeneous expressions in $\mathcal N(S)$. A coefficient $x_jy$ in
the integral has weight $\wt(x)+\wt(y)-j-1$. Its pairing with $z$
has total weight at most $\wt(x)+\wt(y)+\wt(z)-1$, so induction
applies again. If the coefficient is scalar, its bracket is zero.
Normal multiplication preserves the first two terms by definition.
This proves both assertions, including the case in which both inputs
are composite. Neither a preferred parenthesization nor independence
of normal words was used.
\end{proof}

\begin{theorem}[Recursive strong-generation criterion]\label{thm:recursive}
Let \(V\) be a nonnegatively graded vertex algebra with \(V_0=\CC\one\).
Choose finite homogeneous subspaces \(S_d\subset V_d\),
\(1\leq d\leq h\), where $h\geq2$, and let \(\mathcal N\) be their differential
normally ordered span, allowing every parenthesization. Put
\(K=S_1\) and \(A=S_2\). Assume
\begin{enumerate}
\item \(k_ns,a_ns\in\mathcal N\) for every \(k\in K\), \(a\in A\),
\(s\in S_d\), and \(n\geq0\);
\item for \(d\geq3\),
\[
 S_d\subset A_0S_{d-1}+\mathcal N_d(S_1+\cdots+S_{d-1}).
\]
\end{enumerate}
Then \(\mathcal N\) is a vertex subalgebra. In particular, the chosen
spaces strongly generate the vertex algebra that they generate.
\end{theorem}

\begin{proof}
The first assumption, the derivation property and \([T,a_0]=0\)
give \(a_0\mathcal N\subset\mathcal N\) for \(a\in A\).
Explicitly, $a_0(x_{-1}y)=(a_0x)_{-1}y+x_{-1}(a_0y)$, so
induction on the construction of a normal expression proves this
stability before any singular-product closure is known.

Now induct lexicographically on
\((\wt(s)+\wt(t),\wt(s))\) for generator pairs \(s,t\).
Left weight one or two is the first assumption. Otherwise the second
assumption reduces the left input to \(s=a_0b\), with
\(\wt(b)=\wt(s)-1\), plus normal remainders. The latter are covered
by Lemma~\ref{lem:filtered-wick}: a homogeneous remainder of weight
$\wt(s)$ in lower-weight generators has no single undifferentiated
generator of that weight. Use
\begin{equation}\label{eq:recursive-jacobi}
 (a_0b)_nt=a_0(b_nt)-b_n(a_0t).
\end{equation}
The first term has smaller total pair weight before applying \(a_0\),
which preserves \(\mathcal N\). The known normal expression for
\(a_0t\) has weight \(\wt(t)+1\). Every derivative or composite term
reduces total pair weight in the second term. The only remaining
possibility is a single undifferentiated generator of weight
\(\wt(t)+1\); its pairing with \(b\) has the original total weight
and smaller left weight. It is therefore an earlier induction case.
If \(\wt(t)=h\), this single-generator possibility is absent.

All generator singular products close. Wick induction then gives closure
for all normal expressions. Translation and normal multiplication
preserve \(\mathcal N\) by definition. For $r\geq1$,
\[
 x_{-r}y=\frac{1}{(r-1)!}(T^{r-1}x)_{-1}y,
\]
so all negative modes close too.
\end{proof}
For the Wick identities see \cite[Section 2]{DSK}.
The full local derivation is
\artifact{hi_cft_20260923_next_structure_recursive_strong_certificate.md}.

\begin{remark}[Finite scope of the input]\label{rem:finite-scope}
Every membership in the first hypothesis has output weight at most
$h+1$. Thus the criterion avoids checking all generator pairs through
weight $2h-1$. The second hypothesis is essential: it asserts actual
zero-product identities, rather than agreement of a dimension or a
character coefficient. The proof takes place in the given ambient
vertex algebra and makes no assertion that its normal words are free.
This resolves the all-weight implication formerly marked F2.
\end{remark}

\subsection{The finite spaces and prerequisite controls}
Choose actual quotient lifts \(P_3\subset S_3\) and an actual triplet
\(P_4\subset U_4\), with joint current highest weight \((0,2)\), and set
\begin{align}
 N_4&=K_{-1}S_3+TS_3+(S_2)_{-1}S_2,&W_4&=N_4+P_4,\label{eq:normal4}\\
 N_5&=K_{-1}W_4+TW_4+(S_2)_{-1}S_3.&&\label{eq:normal5}
\end{align}
These denote full differential normal spaces, using the lower closure
identities to reorder expressions. A computational solver may use a
proper subset of actual words in these spaces.

The triplet is obtained from an actual \(a_0p\) by current-zero-mode
projectors and lowering. It satisfies
\(P_4\subset (S_2)_0S_3+N_4\).

\begin{lemma}[Bounds for the current projectors]\label{lem:current-bounds}
On ambient weight $d$, a Cartan element represented by $h\in L\otimes\mathbb R$
has eigenvalues $m$ satisfying $m^2\leq2d(h,h)$. At weight four
the two current highest weights are therefore at most $4$ and $24$.
\end{lemma}
\begin{proof}
Cartan zero modes act on a Fock monomial of momentum $\alpha$ by
$(h,\alpha)$, and commute with its oscillators. The grading gives
$(\alpha,\alpha)\leq2d$, so the estimate is Cauchy--Schwarz in the
positive lattice form. Every eigenvalue is integral. Apply the estimate
to $(h,h)=2,78$ and $d=4$, and then to the highest vectors in the
finite-dimensional semisimple current decomposition.
\end{proof}

Consequently finite polynomials in \(FE\) give exact projectors.
For an input of weight \((0,2)\), one may use
\[
 \prod_{\ell=1}^{2}\left(1-\frac{F_LE_L}{\ell(\ell+1)}\right),
 \qquad
 \prod_{\ell=2}^{12}\left(1-\frac{F_RE_R}{\ell(\ell+1)-2}\right).
\]
The raising operator acts first in every \(FE\) factor.
Indeed, on the weight-$m$ line in an $\mathfrak{sl}_2$ representation
of highest weight $j$, $FE$ acts by $(j-m)(j+m+2)/4$. The first
polynomial retains left highest weight zero. At right weight two the
second retains highest weight two and kills every possible larger
highest weight. Thus all denominators are nonzero and the projector
has a finite, justified range.

\begin{certificate}[Lower identities and the actual triplet]\label{cert:lower}
The lower-closure artifact contains 3,970 completed exact identities.
The actual \(P_4\) vectors satisfy 18 current-zero-mode identities and
six independently checked inclusions
\[
 L(1)P_4\subset S_3,\qquad L(2)P_4\subset S_2.
\]
These direct Virasoro checks are available before the weight-four gate.
\end{certificate}
The two gates left for the all-weight argument are
\begin{align}
 (S_2)_0S_3&\subset W_4 &&\text{(Gate 4)},\label{eq:gate4}\\
 (S_2)_0P_4&\subset N_5 &&\text{(Gate 5)}.\label{eq:gate5}
\end{align}

\section{Gate 4: exact seeds, raising, and singlet recovery}\label{sec:gate4}
Put \(P=S_3/N_3\), of independently known dimension 322. For a fixed
joint-highest primitive field \(a\in S_2\), the map
\[
 f_a:P\longrightarrow B_4/W_4,\qquad [p]\longmapsto[a_0p]
\]
is well defined because the lower relations give \(a_0N_3\subset N_4\).
Current covariance implies that its kernel is stable under both commuting
raising operators \(E_L,E_R\).

\begin{lemma}[Raising saturation]\label{lem:raising}
If \(C_a\subset\ker f_a\) and
\(P=E_LP+E_RP+C_a\), then \(f_a=0\).
\end{lemma}
\begin{proof}
On \(M=P/\ker f_a\) the hypothesis says
\(M=E_LM+E_RM\). The commuting ideal \((E_L,E_R)\) acts nilpotently
because all current weights are bounded. Iterating gives
\(M=(E_L,E_R)^rM=0\) for sufficiently large \(r\).
\end{proof}

In this application nilpotence is entirely finite-dimensional. At weight
three Lemma~\ref{lem:current-bounds} bounds the two Cartan eigenvalues
by $|m_L|\leq3$ and $|m_R|\leq21$. Either raising operator increases
$m_L+m_R$ by two. Every word of length $25$ in the two raising
operators therefore vanishes on $P$. The quotient by $\ker f_a$
inherits these operators because the kernel is stable. This gives an
explicit termination bound for the argument; there is no limiting step.

The exact raising coordinates comprise 644 images. Their rank is 294.
For each of nine independent highest left fields, the completed seed
sets and these images have rank 322. A nonzero minor at the good split
prime \(p=1{,}000{,}003\) proves a nonzero characteristic-zero minor.
This proves spanning because the exact ambient quotient dimension is
already 322. It does not give an upper rank for an unknown weight-four
normal space.

The original exact run supplies 2,145 seed identities. The saturation
calculation asks for 130 further inputs; 118 have direct exact residual
certificates. The remaining twelve outputs have weight \((0,0)\).
An exact replay of the known kernel and twelve additional identities
at output weight \((0,2)\) shows that both raising currents kill each
of these twelve remaining classes modulo \(W_4\).

\begin{lemma}[Recovery from a joint-singlet projection]\label{lem:singlet}
Let \(t\in B_4\) have weight \((0,0)\), with
\(E_Lt,E_Rt\in W_4\). Define
\begin{equation}\label{eq:projector}
 \Pi=\prod_{j\in\{2,4\}}
       \left(1-\frac{4F_LE_L}{j(j+2)}\right)
       \prod_{j\in\{2,4,\ldots,24\}}
       \left(1-\frac{4F_RE_R}{j(j+2)}\right).
\end{equation}
If \(\Pi t\) is a linear combination of \(\Pi n_i\), with
\(n_i\in N_4\), then \(t\in W_4\).
\end{lemma}
\begin{proof}
Both \(W_4\) and \(N_4\) are stable under current zero modes.
Every \(FE\) factor kills \([t]\in B_4/W_4\), so
\([\Pi t]=[t]\). Each \(\Pi n_i\) remains in \(N_4\subset W_4\).
Thus \([t]=0\). Equivalently, the weight-zero highest class is a singlet
in the finite semisimple current module. No dimension assertion about
the selected normal columns is needed.
\end{proof}

The prime search selects 21 independent projected normal columns.
All twelve projected identities have complete residual zero over
\(\QQ(\sqrt{13},\sqrt{-3})\). The final verifier repeats those
residuals by independent integer denominator clearing, authenticates
the raising premise, and repeats all nine saturation ranks. By
Lemmas~\ref{lem:raising} and~\ref{lem:singlet}, all \(9\cdot322=2898\)
required highest-left identities hold. Current covariance gives the
charged descendants. The second neutral field follows from
\(\omega_B-\omega_{K_L}-\omega_{K_R}\) minus the retained neutral field.
This proves \eqref{eq:gate4} with the certified lower controls.

\subsection{Positive modes available after Gate 4}
\label{sec:bootstrap}
The normal-product argument at the next gate uses positive modes on
$S_3$. We establish them here, before their first use. This also explains
why no additional unverified positive-mode gate is required.

\begin{lemma}[Virasoro and current bootstrap]\label{lem:bootstrap}
The lower identities, the direct $P_4$ controls, and Gate~4 imply
\begin{align}
 (S_2)_1S_3&\subset S_3,& (S_2)_2S_3&\subset S_2,&
 (S_2)_3S_3&\subset K,\label{eq:positive-S3}\\
 K_1W_4&\subset S_3,&K_2W_4&\subset S_2,&K_3W_4&\subset K,
 \label{eq:positive-W4}\\
 L(1)N_5&\subset W_4,&L(2)N_5&\subset S_3.&&
 \label{eq:positive-N5}
\end{align}
The next mode in each of the first two lines is scalar, and all further
modes vanish.
\end{lemma}
\begin{proof}
The exact lower inputs give all singular products of $K$ and $S_2$
with $S_2$ in the expected lower spaces. The provenance
$S_3=N_3+(S_2)_0S_2$ and the current commutator give
\[
 K_0S_3\subset S_3,\qquad K_1S_3\subset S_2,\qquad
 K_2S_3\subset K.
\]
For instance,
\begin{equation}\label{eq:k1-provenance}
 k_1(a_0b)=(k_0a)_1b+(k_1a)_0b+a_0(k_1b).
\end{equation}
The Virasoro commutators for a weight-two vector are
\begin{align}
 [L(1),a_n]&=(2-n)a_{n+1}+(L(1)a)_n,\label{eq:L1-weight2}\\
 [L(2),a_n]&=(3-n)a_{n+2}+3(L(1)a)_{n+1}+(L(2)a)_n.
 \label{eq:L2-weight2}
\end{align}
They show $L(1)S_3\subset S_2$ and $L(2)S_3\subset K$.
Applied to the three summands of $N_4$, they give
$L(1)N_4\subset S_3$ and $L(2)N_4\subset S_2$.
For a normal input $a_{-1}b$ the latter formula includes the scalar
term $(L(2)a)_{-1}b$; it must not be omitted. It lies in $S_2$.
The six direct controls of Certificate~\ref{cert:lower} extend these
inclusions to $W_4$ without using Gate~5.

Each of the eight charged primitive multiplets and the two neutral
singlets in $S_2$ is quasiprimary. Indeed, $L(1)$ is equivariant for
current zero modes, and none of their current types is $(3,1)$ or
$(1,3)$, the two types in $K$. For such a vector $a$ and $p\in S_3$,
Gate~4 and the preceding controls give
\begin{align*}
 2a_1p&=L(1)(a_0p)-a_0L(1)p\in S_3,\\
 a_2p&=L(1)(a_1p)-a_1L(1)p\in S_2,\\
 2a_3p&=L(2)(a_1p)-a_1L(2)p\in K.
\end{align*}
In the last line the possible $(L(2)a)_1$ term is zero since $L(2)a$
is scalar. The affine normal and derivative part of $S_2$ is handled
by the current iterate identities. This proves
\eqref{eq:positive-S3} on all of $S_2$.

The actual construction gives $P_4\subset(S_2)_0S_3+N_4$.
Equation~\eqref{eq:k1-provenance}, now with $b\in S_3$, gives
$K_1P_4\subset S_3$; the same inclusion for $N_4$ follows on its
normal generators. Since currents are primary of weight one,
\[
 [L(1),k_1]=-k_2,\qquad [L(2),k_1]=-k_3.
\]
The controls on $L(1)W_4,L(2)W_4$ consequently give
\eqref{eq:positive-W4}. The remaining modes follow from grading.

Finally apply $L(1),L(2)$ to
$N_5=K_{-1}W_4+TW_4+(S_2)_{-1}S_3$. The current and translation
summands follow from the established inclusions and the Virasoro
relations. For the third summand, the required formulas are
\begin{align*}
 L(1)(a_{-1}p)
 &=3a_0p+(L(1)a)_{-1}p+a_{-1}L(1)p,\\
 L(2)(a_{-1}p)
 &=4a_1p+3(L(1)a)_0p+(L(2)a)_{-1}p+a_{-1}L(2)p.
\end{align*}
Every term belongs respectively to $W_4$ or $S_3$ by Gate~4,
\eqref{eq:positive-S3}, and the lower identities. Reordering the last
normal term uses only already established lower products. This proves
\eqref{eq:positive-N5}.
\end{proof}

\section{Gate 5 from Jacobi identities}\label{sec:gate5}
The proof uses the full normal spaces in
\eqref{eq:normal4}--\eqref{eq:normal5}, not only the selected solver
columns. To see that the displayed formula for $N_5$ includes every
lower normal word of weight five, remove a leading derivative, or
split a normal product into weights $1+4$ or $2+3$. Reversed splits
are reordered by skew symmetry for normal products. Its correction
terms are derivatives of singular products, already controlled by
Lemma~\ref{lem:bootstrap} and the lower identities. Quasi-associativity
changes the parenthesization with correction terms of the same type.
Thus the resulting words lie in
$K_{-1}W_4+TW_4+(S_2)_{-1}S_3$.

\begin{lemma}[Normal-space controls]\label{lem:normal-controls}
Before using Gate 5, the lower controls and Gate 4 give
\begin{equation}\label{eq:normal-lemmas}
 (S_2)_0N_4\subset N_5,\qquad (N_3)_0S_3\subset N_5.
\end{equation}
\end{lemma}
\begin{proof}
The first follows by zero-mode derivation and skew reordering. For the
second, \((Ta)_0=0\) and, for a current \(x\),
\begin{equation}\label{eq:iterate}
 (x_{-1}a)_0p
 =\sum_{i\geq0}x_{-1-i}(a_ip)
  +\sum_{i\geq0}a_{-1-i}(x_ip).
\end{equation}
All inner products belong to the required lower spaces. To make the
translation reduction explicit, repeatedly use
\[
 x_{-r}v=\frac{T(x_{-(r-1)}v)-x_{-(r-1)}Tv}{r-1},\qquad r\geq2.
\]
The base terms are precisely those in \eqref{eq:normal5}; the same
identity applies to \(a_{-r}v\) with \(a\in S_2\). Thus no
unproved \(P_3\)-by-\(P_3\) closure enters \eqref{eq:normal-lemmas}.
All positive inner modes in \eqref{eq:iterate} are supplied by
\eqref{eq:positive-S3} and the current controls established before
this section. More explicitly the first sum has inner vectors in
$W_4,S_3,S_2,K,\CC\one$, in that order; the second has inner
vectors in $S_3,S_2,K,\CC\one$. Both sums are finite by grading.
\end{proof}

If \(P_4\subset N_4\), Gate 5 already follows. Otherwise
\(Q=W_4/N_4\) is the irreducible current triplet of highest weight
\((0,2)\), since the actual \(P_4\) is irreducible. Let
\(q_-,q_0,q_+\) be its weight basis, and set
\[
 \beta(a,p)=[a_0p]\in Q,\qquad
 F(a,q)=[a_0\widetilde q]\in B_5/N_5.
\]
The second map does not depend on the representative
\(\widetilde q\in W_4\), by \eqref{eq:normal-lemmas}.

\begin{lemma}[Jacobi reduction]\label{lem:jacobi}
If \(a,b\in S_2\) satisfy \(a_0b\in N_3\), then for \(p\in S_3\),
\begin{equation}\label{eq:jacobi-reduction}
 F(a,\beta(b,p))=F(b,\beta(a,p)).
\end{equation}
\end{lemma}
\begin{proof}
The difference is the class of \((a_0b)_0p\), which vanishes by
\eqref{eq:normal-lemmas} and \([a_0,b_0]=(a_0b)_0\).
\end{proof}

\subsection{The charged families}
Let \(a_0,\ldots,a_7\) denote the eight charged highest fields, with
current weights
\[
 (0,4),(0,8),(0,8),(0,12),(1,1),(1,5),(1,7),(1,13).
\]
In a product such as \((a_i)_0p\), the parenthesized subscript denotes
the field label and the next subscript denotes its mode.
The saved exact relations give
\[
 \beta(a_3,p_3)=c_3q_-,\qquad \beta(a_6,p_6)=c_6q_-,
 \qquad c_3c_6\neq0,
\]
where \(p_3=\texttt{r3\_000910}\) has weight \((0,-14)\), and
\(p_6=\texttt{r3\_000899}\) has weight \((-1,-9)\).
The coefficients are authenticated in the final Jacobi certificate.
Exact combinations of the lower relations establish \((a_i)_0a_j\in N_3\)
for the following pairs:
\begin{center}
\begin{tabular}{@{}cccl@{}}
\toprule
\(i\)&\(j\)&Input&Weight of \((a_i)_0p_j\)\\
\midrule
0&3&\(p_3\)&\((0,-10)\)\\
1,2&3&\(p_3\)&\((0,-6)\)\\
3&6&\(p_6\)&\((-1,3)\)\\
4&6&\(p_6\)&\((0,-8)\)\\
5&3&\(p_3\)&\((1,-9)\)\\
6&3&\(p_3\)&\((1,-7)\)\\
7&3&\(p_3\)&\((1,-1)\)\\
\bottomrule
\end{tabular}
\end{center}
Every listed output weight lies outside \(Q\), so
\(\beta(a_i,p_j)=0\). Equation~\eqref{eq:jacobi-reduction} and
\(c_j\neq0\) give \(F(a_i,q_-)=0\). Raise this resulting zero
with \(a_i\) held highest to obtain the other two right inputs.
Current covariance gives all left descendants.

\begin{remark}
The weight exclusion is used only at the displayed bottom inputs.
It need not hold after raising \(p_j\); using it there would be an
incorrect stronger assertion. We raise the proved zero identity instead.
\end{remark}

\subsection{The neutral family}
Let \(C\) be the full current span, inside \(P=S_3/N_3\), of
brackets of charged weight-two fields. The charged-span certificate
reconstructs 28 actual charged seeds from the exact lower relations.
Together with the 644 lowering images, of rank 294, these have rank 322.
The independent verifier also checks both \(\mathfrak{sl}_2\) commutators
and all cross-axis commutators against the saved exact raising matrices.
Since \(C\) is a current submodule, the nilpotent-ideal argument gives
\[
 P=F_LP+F_RP+C\quad\Longrightarrow\quad P/C=0.
\]
Once again the finite-prime full rank is a positive spanning witness
inside an independently exact 322-dimensional space.

For the retained neutral singlet \(e\), write
\[
 e_0a_3=\sum_i t_i(x_i)_0y_i+n,
 \qquad n\in N_3,
\]
with charged \(x_i,y_i\). Apply zero modes to \(p_3\). By Jacobi,
the charged conclusion and \eqref{eq:normal-lemmas},
\[
 (e_0a_3)_0p_3
 =\sum_i t_i[(x_i)_0,(y_i)_0]p_3+n_0p_3\in N_5.
\]
Indeed Gate 4 puts each inner product in \(W_4\), and the next charged
product is controlled on both \(N_4\) and \(P_4\). A second Jacobi
identity gives
\[
 e_0((a_3)_0p_3)=(a_3)_0(e_0p_3)+(e_0a_3)_0p_3\in N_5.
\]
The nonzero coefficient of \(q_-\) proves the neutral bottom identity;
covariance gives the triplet. The other neutral field is obtained from
the total stress minus the two affine Sugawara stresses and \(e\).
This proves Gate 5. Its 24 independent highest channels comprise
23 charged channels and one neutral channel; they are consequences
of the displayed identities, not 24 separately recomputed grade-five
lattice products.

\section{From the finite gates to properness}\label{sec:properness}
Lemma~\ref{lem:bootstrap} supplies the positive-current products and
$L(1)N_5\subset W_4$, $L(2)N_5\subset S_3$ before Gate~5.
For a quasiprimary primitive $a\in S_2$,
Gate 5 now yields, for \(p\in P_4\),
\begin{align*}
 2a_1p&=L(1)(a_0p)-a_0L(1)p\in W_4,\\
 a_2p&=L(1)(a_1p)-a_1L(1)p\in S_3,\\
 2a_3p&=L(2)(a_1p)-a_1L(2)p\in S_2,\\
 a_4p&=L(2)(a_2p)-a_2L(2)p\in K.
\end{align*}
The next mode is scalar. The 30 earlier prime-only \(a_1P_4\)
tests are therefore derived exactly here; they are not an omitted gate.

All hypotheses of Theorem~\ref{thm:recursive} hold for
\(K,S_2,S_3,P_4\). Their normal span is a vertex subalgebra of \(U\)
containing \(K,u,\omega_B\), hence equals \(U\). A normal word in
positive-weight generators cannot acquire weight one unless it is a
weight-one generator. Thus \(U_1=K\), proving
Proposition~\ref{prop:mainproper}.

\begin{proof}[Proof of Corollary~\ref{cor:initial-grades}]
A differential normal word of weight two uses only $K$ and $S_2$;
its current descendants are already in $S_2$. A word of weight three
uses only $K,S_2,S_3$, and the lower normal words lie in $N_3$.
Thus the strong-generation result proves $U_2=S_2$ and $U_3=S_3$.
The exact PCT certificate gives $\theta_gS_2=S_2$, while
$\theta_gK=K$. Since $S_2\subset U$ contains the original generators,
$\theta_gU=U$. The inherited positive form makes $U$ a unitary
subalgebra of $B$, hence simple by \cite[Proposition 5.29]{CKLW}.
It contains $\omega_B$, so has central charge eight, the same
nonnegative grading, and $U_0=\CC\one$. The character assertion is
now the definition of the graded trace.
\end{proof}

\begin{reviewnote}{F1: what the finite certificate establishes}
This proof of properness has the strongest computational support in the
manuscript. It still relies on correctness of the lattice mode engine,
the saved actual-vector recipes, and the stated algebraic lemmas.
Hashes authenticate the records; hashes alone do not prove those lemmas
or independently reconstruct every earlier mode calculation. The Lean
checks in Appendix~\ref{sec:formalization} cover selected abstract
implications, not the VOA construction or the complete strong-generation
argument. No second implementation of the entire construction is
claimed. In particular, \(\dim U_4=6010\) is not used as an input to
this proof.
\end{reviewnote}

% Fragment: regularity, branching, fusion-ring reduction, and characters.
% The main document supplies theorem environments, reviewnote, and artifact.

\section{Regularity from a positive conformal-vector identity}
\label{sec:regularity}

This section concerns the actual PCT-stable subalgebra $U\subset B$ defined
in the construction. Its conclusions do not require properness. We use
the convention $Y(a,z)=\sum_{n\in\mathbb Z}a_nz^{-n-1}$ and write
$L^e(n)=e_{n+1}$ for the Virasoro operators of a conformal vector $e$.
The subalgebras used below need not contain the ambient conformal vector.
We call an inclusion $A\subset V$ \emph{grading compatible} if the
$A$-weight of each homogeneous vector of $A$ equals its $V$-weight.
Thus a mode $a_n$ changes both gradings by $\operatorname{wt}(a)-n-1$;
the equality of the two energy operators is not required on all of $V$.
All subalgebras in the positive identity below are grading compatible
and stable under the ambient PCT involution.

\begin{certificate}[Embedded structural data]
\label{cert:structural-inputs}
The following are assertions about actual fields and their exact modes
in the lattice construction, in addition to the finite closure ledger.
\begin{enumerate}
\item The Potts fields lie in $U$, are PCT stable, and generate the
regular unitary algebra $K(\mathfrak{sl}_2,3)$, with its translation
and intrinsic grading inherited from the ambient VOA. The exact
records are \artifact{hi_cft_20260923_proof_potts.json} and
\artifact{hi_cft_20260923_proof_potts_exact.json}.
\item The actual fields $z,a$ of weights two and three lie in $U$,
are PCT stable, have central charge $8/5$, and satisfy the stress,
primary, and translation identities in
\artifact{hi_cft_20260923_proof_pf13_exact.json},
\artifact{hi_cft_20260923_proof_pf13_membership.json}, and
\artifact{hi_cft_20260923_proof_remaining_coset.json}.
Their full singular first rows agree with the parafermion rows at
$c=8/5$, $\lambda=14/473$, as polynomials in $z,a,a_1a,a_1(a_1a)$.
The polynomial records are
\artifact{hi_cft_20260923_proof_pf13_upper_poles.json},
\artifact{hi_cft_20260923_proof_pf13_nonlinear.json},
\artifact{hi_cft_20260923_proof_pf13_truncation.json},
\artifact{hi_cft_20260923_proof_pf13_simple_pole_algebra.json}, and
\artifact{hi_cft_20260923_proof_pf13_dly.json}.
The membership record also gives the positive identity
\eqref{eq:positive-stress}.
\item The lattice--Potts core $T$ below is embedded with its stated
charge realization. The six joint spectral vectors used to construct
its order-thirteen extension are actual vectors of $U_2$, with the
eigenvalues printed below. Their exact source is
\artifact{hi_cft_20260923_next_structure_core_extension.json}; the
construction of $T$ is described in
\artifact{hi_cft_20260923_proof_report.md}, Section~4.
\item The zero-mode action on the complementary eighty currents is
irreducible, and the four mixed lattice tops of dimension eighty-one
are irreducible as $U$-tops. These are exact
field-action assertions, as recorded in
\artifact{hi_cft_20260923_proof_current_module_v2.json} and
\artifact{hi_cft_20260923_next_81top_all.json}, respectively.
\end{enumerate}
All polynomial equalities here mean exact vector equalities in the
fixed lattice realization, not projected or floating-point matches.
\end{certificate}

The input from the local construction is the identity
\begin{equation}
 \omega_U=\omega_{A_{1,1}}
   +\frac{3403}{4563}\omega_{A_{1,39}}
   +\frac{25}{6}\overline p+\frac{25}{26}\overline z.
 \label{eq:positive-stress}
\end{equation}
Here $p$ and $z$ are the conformal vectors of actual embedded copies of
the three-state Potts algebra $K(\mathfrak{sl}_2,3)$ and the parafermion
algebra $K(\mathfrak{sl}_2,13)$; a bar denotes normalized averaging over
the compact current group. Each subalgebra is regular and unitary, and
its intrinsic grading agrees with the ambient grading on its states.
The averages take place in the finite-dimensional space $U_2$.
More concretely, $p,z$ lie in the already constructed current-stable
space $S_2$, whose invariant subspace is spanned by the four mutually
orthogonal stresses $\omega_{A_{1,1}},\omega_{A_{1,39}},e,f$.
For either vector $v$, its Haar average is therefore
\[
 \overline v=\sum_{w\in\{\omega_{A_{1,1}},\omega_{A_{1,39}},e,f\}}
       \frac{\langle v,w\rangle_g}{\langle w,w\rangle_g}\,w.
\]
Thus the exact coefficients in \eqref{eq:positive-stress} can be
verified inside $S_2$ without assuming $S_2=U_2$ or properness.
The compact orbit of a conformal vector has compact convex hull in the
real space of PCT-fixed weight-two vectors. Its Haar average belongs
to that hull. Carath\'eodory's theorem therefore replaces each average
by a finite convex combination of orbit points. In particular,
\begin{equation}
 \omega_U=\sum_{i=1}^N c_i e_i,\qquad c_i>0.
 \label{eq:finite-positive-stress}
\end{equation}
Zero coefficients, if any, are omitted. All simple modules of each
regular subalgebra $A_i$ have nonnegative lowest weights, with zero
occurring only for its vacuum module.

\begin{reviewnote}{R1: the embedded regular factors are local inputs}
Certificate~\ref{cert:structural-inputs} states the model-specific
field inputs explicitly. Proposition~\ref{prop:pf13-finite-reconstruction}
below proves the transport from the complete checked parafermion rows
to the actual factor identification. It does not replace verification
of the listed lattice identities, Potts embedding, or positive stress
identity. These inputs are separate from the finite closure theorem.
\end{reviewnote}

\begin{proposition}[Finite reconstruction of the parafermion factor]
\label{prop:pf13-finite-reconstruction}
Let $V,V'$ be simple vertex algebras of CFT type generated by a
conformal vector $L$ and a weight-three primary $a$. Set
$b=a_1a$ and $d=a_1b$, of weights four and five. Suppose the full
singular products $a_na,a_nb,a_nd$, $n\geq0$, agree as homogeneous
differential normal polynomials in $L,a,b,d$, with the same normal
ordering. Suppose also that their self-row is
\[
 a_5a=\frac c3\mathbf1,\quad a_4a=0,\quad a_3a=2L,\quad
 a_2a=TL,\quad a_1a=b,\quad a_0a=\frac12Tb-\frac1{12}T^3L,
\]
and that $a_nb=0$ for $n\geq4$, $a_3b\in\mathbb Ca$, and
$a_2b\in\mathbb CTa$. Then there is an isomorphism $V\cong V'$
identifying $L,a,b,d$. In particular $a_1d$ and both simple poles
must be specified polynomials; they cannot be left as independent
new generators.
\end{proposition}
\begin{proof}
We first recover the remaining generator OPEs without a circular
normal-product reduction. The Virasoro rows follow by commuting
Virasoro modes through $b=a_1a$ and $d=a_1b$:
\[
 L(m)b=(2m+1)a_{m+1}a,\qquad
 L(m)d=(2m+1)a_{m+1}b+a_1L(m)b\quad(m\geq1).
\]
These are linear in generators and derivatives under the displayed
hypotheses. Thus Virasoro creation operators can be moved to the
outside of any normal word using already known commutators.
Sesquilinearity removes derivatives. Since three non-Virasoro
factors have weight at least nine, the possible two-factor cores
through weight eight are exactly
\[
\begin{array}{c|c}
\text{weight bound}&\text{cores besides individual fields}\\\hline
6&:aa:\\
7&:aa:,\ :ab:\\
8&:aa:,\ :ab:,\ :bb:,\ :ad:.
\end{array}
\]
Reversing the order of a core uses skew symmetry and its already
known pair bracket. Commuting an external $a_n$ or $b_n$ through
outer Virasoro creation operators reduces their number; any resulting
negative mode is itself a normal product.

Write $A_r$ for the evaluation of $a_nP$, $n\geq0$, on every normal
polynomial $P$ of weight at most $r$, and $B_r$ for the corresponding
assertion with $b_n$. We use the noncommutative Wick identity \cite[Section~2]{DSK}
\[
 [x_\lambda:yz:]=:[x_\lambda y]z:+:y[x_\lambda z]:
              +\int_0^\lambda[[x_\lambda y]_\mu z]\,d\mu.
\]
The following order is acyclic. First $A_6$ is known: the only
two-factor core is $:aa:$, and $[a_\lambda a]$ is linear in the
vacuum and derivatives of $L,b$, whose brackets with $a$ are given.
The iterate identity
\begin{equation}
 b_ns=a_1(a_ns)-a_0(a_{n+1}s)
          +a_{n+1}(a_0s)-a_n(a_1s)
 \label{eq:pf13-b-iterate}
\end{equation}
with $s=b$ now determines $b_nb$, because each inner $a_jb$ has
weight at most six. Next $A_7$ follows: its new core $:ab:$ uses
only the just-computed $b,b$ row in the Wick integral.
Equation~\eqref{eq:pf13-b-iterate} with $s=d$ then determines
$b_nd$, since every inner $a_jd$ has weight at most seven.

At this point $B_6$ is known. For its core $:aa:$ the Wick integral
brackets a coefficient of $[b_\lambda a]$, of weight at most six,
with $a$, which is $A_6$ after skew symmetry. Single-field inputs
use the known $b,b$ and $b,d$ rows. We obtain $A_8$ next: its core
$:bb:$ uses $B_6$ on coefficients of $[a_\lambda b]$, and its core
$:ad:$ uses the known $L,d$ and $b,d$ rows. Finally $B_7$ follows
from $B_6$ on the coefficients of $[b_\lambda a]$ in the new core
$:ab:$. With $s=d$, the second iterate identity
\[
 d_ns=a_1(b_ns)-a_0(b_{n+1}s)
          +b_{n+1}(a_0s)-b_n(a_1s)
\]
computes $d_nd$: the first two terms use $A_8$, and the last two
use $B_7$. All generator singular products have therefore been
expressed by the same polynomials in the two algebras.

For the uniqueness step, take the vertex subalgebra $D\subset V\times V'$
generated by the diagonal pairs of the four named fields. Common
polynomial OPE closure and filtered Wick induction show that $D$ is
their differential normal span. Every generator has positive weight,
so $D_0=\mathbb C(\mathbf1,\mathbf1)$ and $D_n=0$ for $n<0$.
It follows that $D$ has a unique maximal proper graded ideal: every
proper graded ideal has zero degree-zero part, and so does their sum.
The coordinate projections are surjective because $L,a$ generate
each target. They preserve the conformal vector and hence the grading.
Their kernels are maximal proper graded ideals because $V,V'$ are
simple. The kernels coincide, which gives the asserted isomorphism.
No freeness or independence of normal monomials has been used.
\end{proof}

The exact parafermion records instantiate this proposition at
$c=8/5$ and $\lambda=14/473$. The upper-pole record supplies the
self-row and the remaining positive poles, the nonlinear record
supplies $a_0b$, and the truncation record supplies $a_1d$.
The algebraic simple-pole derivation supplies $a_0d$ from the
independently checked upper poles and skew symmetry. The
Dong--Lam--Yamada change-of-generators record matches the seven
weight-six decoupling coefficients with those of
$K(\mathfrak{sl}_2,13)$, using \cite[Appendix B]{DLY}. The common first rows are the recursive
formulas (5.4), (5.7), and (A.2) of Linshaw \cite{Linshaw}. Simplicity of the actual
generated subalgebra follows from its PCT stability; its translation
operator agrees with the ambient one by the complementary positive
energy argument. Thus the proposition transports the checked
polynomial identities to an identification with the parafermion
algebra. The exact membership, PCT, translation, and polynomial
identities remain the explicit finite inputs to this identification.

\begin{lemma}[Orthogonal energy bound]\label{lem:positive-c2}
Let $V$ be a unitary VOA of CFT type admitting an identity of the form
\eqref{eq:finite-positive-stress}, with the properties just stated.
Then $V$ is $C_2$-cofinite.
\end{lemma}
\begin{proof}
For an ordinary simple $A_i$-module $M$, put
\[
 C_{2,A_i}(M)=\operatorname{span}\{a_{-2}m:a\in A_i,\ m\in M\}.
\]
These quotients are finite-dimensional by \cite[Proposition~5.2]{ABD}.
They are graded for the intrinsic operator $L^{e_i}(0)$.
There are finitely many simple $A_i$-modules. Hence
there is a finite number $D_i$ bounding the intrinsic energy in all
their $C_{2,A_i}$ quotients; this includes their lowest weights.
Regularity decomposes the weak $A_i$-module $V$ into ordinary simples.
Consequently every intrinsic eigenvector of energy above $D_i$ lies
in $C_{2,A_i}(V)$. The intrinsic energy preserves each finite-dimensional
ambient grade $V_n$ and is self-adjoint for its positive form. We may
therefore diagonalize it on $V_n$. Since
\[
 C_{2,A_i}(V)\subset C_2(V),
\]
every vector in $V_n\cap C_2(V)^\perp$ has zero projection to the
intrinsic eigenvectors of energy above $D_i$. It follows that
\[
 \langle L^{e_i}(0)v,v\rangle\le D_i\|v\|^2
 \qquad (v\in V_n\cap C_2(V)^\perp).
\]
For $v\in V_n\cap C_2(V)^\perp$, the conformal-vector identity gives
\[
 n\|v\|^2
 =\sum_i c_i\langle L^{e_i}(0)v,v\rangle
 \le\Bigl(\sum_i c_iD_i\Bigr)\|v\|^2.
\]
The subspace $C_2(V)$ is ambient graded. Its orthogonal complement in
each finite-dimensional $V_n$ represents $(V/C_2(V))_n$, so the last
inequality makes the quotient vanish for
$n>\sum_i c_iD_i$. Finiteness of each $V_n$ proves the claim.
No simultaneous diagonalization of the different energy operators is
required. The same proof applies directly to averages, since every
conjugate factor has the same bound $D_i$.
\end{proof}

\begin{lemma}[Highest spaces for a compatible subalgebra]
\label{lem:compatible-top}
Let $A\subset V$ be grading compatible and regular, and let $M$ be a
weak $V$-module. Suppose a finite-dimensional subspace $E\subset M$ is
stable under every ambient degree-zero mode from $A$ and is killed by
every mode from $A$ of negative ambient degree. Then $E$ is a direct
sum of irreducible Zhu-algebra modules belonging to the tops of simple
ordinary $A$-modules. In particular $L^A(0)|_E$ is diagonalizable,
with eigenvalues among their lowest weights.
\end{lemma}
\begin{proof}
Write
\[
 \Omega_A(M)=\{v\in M:a_nv=0\text{ for all homogeneous }a\in A,
                   \ n\ge\operatorname{wt}_A(a)\}.
\]
Compatibility puts $E$ inside this space. Regularity decomposes
$M\downarrow_A$ as an algebraic direct sum of ordinary simples, and
$\Omega_A$ commutes with that sum. The highest space of an ordinary
simple is its top. Zero modes endow $\Omega_A(M)$ with the corresponding
semisimple Zhu-algebra action; its invariant finite-dimensional
subspace $E$ is therefore a sum of simple tops. This also explains why
one must impose the lowering condition for all states of $A$, rather
than only for a chosen generating set.
\end{proof}

\begin{lemma}[Vacuum-like vectors]\label{lem:vacuum-like}
Let $M$ be a weak module over a VOA $V$. If $v\in M$ satisfies
$L_M(-1)v=0$, then $a_nv=0$ for all $a\in V$ and $n\ge0$, and
$a\mapsto a_{-1}v$ is a $V$-module homomorphism $V\to M$ sending
$\mathbf1$ to $v$.
\end{lemma}
\begin{proof}
Translation covariance gives
$L_M(-1)a_{n+1}v=-(n+1)a_nv$. Lower truncation and downward induction
give $a_nv=0$ for every $n\ge0$. The module Jacobi identity then gives
$b_m(a_{-1}v)=(b_ma)_{-1}v$; equivalently this follows by taking the
constant coefficient in weak associativity, since the fields applied
to $v$ have no negative powers. This is the asserted homomorphism.
\end{proof}

\begin{lemma}[Lowest weights and projectivity of the vacuum]
\label{lem:vacuum-projective}
Let $V$ satisfy Lemma~\ref{lem:positive-c2}, and suppose every component
algebra has nonnegative simple lowest weights and its vacuum is the
unique simple of weight zero. Every simple ordinary $V$-module has
nonnegative lowest conformal weight, the only weight-zero simple is
$V$, and $V$ is projective among grading-restricted generalized modules.
\end{lemma}
\begin{proof}
Let $E=M_h$ be the lowest space of a simple ordinary module.
For a homogeneous state $a$,
$[L_V(0),a_n]=(\operatorname{wt}_V(a)-n-1)a_n$ on every module.
Thus degree-zero modes preserve $E$, and lowering modes kill $E$.
Lemma~\ref{lem:compatible-top} implies that every
$L^{e_i}(0)|_E$ is diagonalizable with nonnegative real spectrum.
Taking traces, rather than assuming the operators commute, gives
\begin{equation}
 h\dim E=\sum_i c_i\operatorname{Tr}_E L^{e_i}(0)\ge0.
 \label{eq:top-trace}
\end{equation}
In particular $h$ is real and nonnegative. If $h=0$, every summand vanishes.
Hence every partial energy is zero
on $E$, which lies in sums of vacuum tops for every $A_i$. Consequently
$L^{e_i}(-1)E=0$, and \eqref{eq:finite-positive-stress} gives
$L_V(-1)E=0$. By Lemma~\ref{lem:vacuum-like}, any nonzero $v\in E$
gives a nonzero module map $V\to M$, hence an
isomorphism.

Now let $\pi:N\twoheadrightarrow V$ be an epimorphism in the generalized
module category. $C_2$-cofiniteness gives finite length, and the preceding
argument excludes negative weights in all its composition factors.
The finite generalized weight-zero space $N_{[0]}$ therefore has no
lower-weight target for a lowering mode. It is again in every
$\Omega_{A_i}(N)$ and is invariant under degree-zero modes. On this
space $L_V(0)$ is nilpotent, so
\[
 0=\operatorname{Tr}_{N_{[0]}}L_V(0)
  =\sum_i c_i\operatorname{Tr}_{N_{[0]}}L^{e_i}(0).
\]
Diagonalizability and nonnegativity of each partial energy imply it is
zero. The vacuum-top argument gives $L_V(-1)N_{[0]}=0$. Generalized
weight decompositions are preserved by $\pi$, so one can choose
$v\in N_{[0]}$ with $\pi(v)=\mathbf1$. Then
\[
 s:V\longrightarrow N,\qquad s(a)=a_{-1}v,
 \qquad \pi s=\mathrm{id}_V,
\]
is a section. Every epimorphism onto the vacuum splits.
\end{proof}

\begin{lemma}[A projective unit in a finite braided $r$-category]
\label{lem:projective-unit}
Let $\mathcal C$ be a finite complex linear abelian braided monoidal
category with simple unit, right-exact tensor product, and
Grothendieck--Verdier duality with dualizing object $\mathbf1$.
If $\mathbf1$ is projective, then $\mathcal C$ is rigid and semisimple.
\end{lemma}
\begin{proof}
Realize $\mathcal C$ as
finite-dimensional modules for a finite-dimensional algebra, denote
underlying vector-space dimension by $d_0$, and let $P$ correspond to
its regular module. Every $Z$ is a quotient of $P^{\oplus d_0(Z)}$.
For fixed $X$, right exactness gives
\[
 d_0(Z\otimes X)\le C_Xd_0(Z),\qquad
 \dim\operatorname{End}(X^{\otimes n})
 \le d_0(\mathbf1)^2C_X^{2n}<n!
\]
for large $n$, where $C_X=\max\{1,d_0(P\otimes X)\}$.
This uses only right exactness. For nonzero indecomposable $X$,
Grothendieck--Verdier evaluation
$DX\otimes X\to\mathbf1$ is nonzero, hence surjective onto the simple
unit, and splits by projectivity. A finite abelian category is
idempotent complete, and its Hom spaces are finite-dimensional.
Thus $X$ is non-negligible and has moderate growth in the precise
sense of \cite[Theorem~1.1]{EP}; that theorem proves $X$ rigid.
Finite direct sums give
rigidity of all objects. Tensoring is now exact, and
\[
 \operatorname{Hom}(X,-)
 \cong\operatorname{Hom}(\mathbf1,-\otimes X^*)
\]
is exact. Every object is projective, so $\mathcal C$ is semisimple.
\end{proof}

\begin{theorem}[Positive stresses imply strong rationality]
\label{thm:positive-stress-rationality}
A simple unitary VOA of CFT type satisfying the hypotheses of
Lemma~\ref{lem:vacuum-projective} is strongly rational. In particular,
the actual generated algebra $U$ is strongly rational, independently
of its properness in $B$.
\end{theorem}
\begin{proof}
Unitarity gives self-contragredience. Lemmas~\ref{lem:positive-c2} and
\ref{lem:vacuum-projective} give $C_2$-cofiniteness and a projective
vacuum. The category of grading-restricted generalized modules is
finite abelian and braided, its tensor product is right exact, and
contragredient duality is Grothendieck--Verdier duality with dualizing
object $V'=V$; see \cite[Sections~2.3 and~3]{McRae}.
Lemma~\ref{lem:projective-unit} proves semisimplicity, hence rationality
in the $C_2$-cofinite setting. This argument does not assume that the
nonvacuum modules are unitarizable.
The local argument is recorded in
\artifact{hi_cft_20260923_next_structure_rationality.md}.
\end{proof}

\section{Branching and the boundary fusion ring}
\label{sec:branching}

From this point we use the properness conclusion and $\dim U_1=6$ of
the finite closure argument. Before that conclusion is supplied, every
statement in this section is conditional on the proper branch.

\subsection{A universal top bound and actual lattice branching}

The sharper bound uses an actual regular subalgebra $A_L\subset U$.
In the notation of the local construction, a rational lattice--Potts
core $\mathcal T$ of central charge $14/5$ and $K(\mathfrak{sl}_2,13)$
of central charge $8/5$ have pointed charges
$(j,k)\in(\mathbb Z/13)^2$ with quadratic form
$-(j^2+k^2)/13$. Put $A=\mathcal T\otimes K(\mathfrak{sl}_2,13)$.
The field certificate supplies nonzero weight-two vectors $v$ with
\[
 (\omega_U-\omega_A)_1v=0,\qquad
 z_1v=\frac{12}{13}v,\qquad
 (\omega_{\mathcal T})_1v=\frac{14}{13}v,
 \qquad p_1v=\frac23v,
\]
at Cartan weights $(\pm1,\pm5)$ and $(0,\pm8)$.
These are the top characters of the simple-current charges
$(5,1)$ and its inverse, up to reversing the second charge.
The label is forced by actual, unreduced lowest weights: the
$\tau$-labelled $\mathcal T$ module of charge $5$ starts at $31/65$,
whereas $14/13-31/65$ is not integral, and the only parafermion modules
starting at $12/13$ are the two inverse simple currents.

Here is why the witness gives an embedded extension. The kernel of
$(\omega_U-\omega_A)_1$ in $U$ is a connected unitary extension of
$A$ with conformal vector $\omega_A$. Regularity of $A$ and finite
ambient homogeneous spaces make this extension finite. In a connected
unitary extension, invertible constituents have multiplicity at most
one and their support is a subgroup: the invariant pairing gives a
nonzero product of a constituent with its inverse into the vacuum
line; associativity then makes multiplication by it invertible on
each supported invertible line. Integer grading makes that support
isotropic. The only isotropic lines for the displayed quadratic form
are $k=5j$ and $k=-5j$, since $5^2=-1\pmod{13}$. The witnessed
nonzero charge generates one of these lines, and both cannot occur
because together they generate nonintegral-spin charges. Consequently,
after the harmless choice of parafermion charge orientation,
\[
 A_L=\bigoplus_{k=-5j}\mathcal T_{1,j}\otimes\mathrm{PF}_{0,k}
 \subset U,\qquad
 \operatorname{Rep}(A_L)\simeq
 \mathrm{Fib}\boxtimes\mathrm{SO}(3)_{13}.
\]
The field witness and extension argument are in
\artifact{hi_cft_20260923_next_structure_order13_core.md} and
\artifact{hi_cft_20260923_next_structure_core_extension.json}.
There are fourteen simple modules, labelled $x_a=(1,a)$ and
$y_a=(\tau,a)$ for $0\le a\le6$. Indeed the free order-thirteen
extension removes the hyperbolic pointed factor; no orbit has a
stabilizer, so there is no fixed-point resolution. The two remaining
labels are the Fibonacci label and the even affine label $2a$.

We give enough charge data to reconstruct their tops. For
$\epsilon\in\{1,\tau\}$ put
\[
 p_1(m)=\begin{cases}0,&3\mid m,\\2/3,&3\nmid m,\end{cases}
 \qquad
 p_\tau(m)=\begin{cases}2/5,&3\mid m,\\1/15,&3\nmid m.\end{cases}
\]
The actual lowest weights of the lattice--Potts modules are
\begin{equation}
 h_{\mathcal T}(\epsilon,j)
 =\min_{m\equiv j\ (13)}
 \left(\frac{m^2}{156}+\frac{\mathbf1_{m\text{ odd}}}{4}
                   +p_\epsilon(m)\right).
 \label{eq:core-charge-minimum}
\end{equation}
An even minimizing $m$ contributes a left affine singlet and an odd
one a left doublet; their right Cartan charge is $m$. Choosing
$|m|\le6$ bounds the minimum by $14/13$, while $|m|\ge13$ gives
$m^2/156>14/13$. Thus $-12\le m\le12$ contains every minimizer.
For the representative $-6\le k\le6$, the parafermion formula is
\begin{equation}
 h_{\mathrm{PF}}(a,k)
 =\frac{a(a+1)}{15}-\frac{k^2}{13}+\max(0,|k|-a).
 \label{eq:pf-charge-minimum}
\end{equation}
The last term is the least affine grade needed for charge $2k$
outside the finite $\mathfrak{sl}_2$ top. This is equivalently the
actual-top formula in \cite[Section~2.4]{ALY}; the parafermion top is
one-dimensional. Finally minimize
$h_{\mathcal T}(\epsilon,j)+h_{\mathrm{PF}}(a,-5j)$ over
$j\in\mathbb Z/13$, using the representative in
\eqref{eq:pf-charge-minimum}.

Let $E$ be any simple $U$-module top and write $x_a,y_a$ also for the
nonnegative integer multiplicities of these $A_L$ tops in $E$.
Lemma~\ref{lem:compatible-top} justifies taking complete $A_L$ tops
here: the finite space $E$ is invariant under all their Zhu zero modes,
not merely contained in the sum of their highest spaces.
Set $D=\dim E$. For each top, let $L$ be the trace of the left affine
energy, $Q=\operatorname{Tr}H_R(0)^2$, and $P,Z$ the Potts and
parafermion energy traces. Its contribution to \eqref{eq:top-trace} is
\begin{equation}
 W=L+\frac{83}{6084}Q+\frac{25}{6}P+\frac{25}{26}Z.
 \label{eq:top-W}
\end{equation}
This is a trace formula on a complete right $\mathfrak{sl}_2$
representation, not an assertion about a single partial-energy
eigenvector. In fact, for an irreducible right $\mathfrak{sl}_2$
representation of highest weight $r$,
\[
 \frac{\operatorname{Tr}H_R(0)^2}{\dim}=\frac{r(r+2)}3,
 \qquad L^{A_{1,39}}(0)=\frac{r(r+2)}{164}.
\]
Thus the Sugawara trace is $3Q/164$, and multiplying by $3403/4563$
gives $83Q/6084$. Haar conjugation preserves traces on $E$.
The exact dimensions and trace contributions are
\[
\begin{array}{c|r|r@{\qquad}c|r|r}
 \text{top}&d&4680W&\text{top}&d&4680W\\\hline
 x_0&1&0&y_0&1&7800\\
 x_1&1&600&y_1&7&36960\\
 x_2&5&10440&y_2&15&110880\\
 x_3&9&56640&y_3&2&4080\\
 x_4&2&5520&y_4&14&81480\\
 x_5&14&148560&y_5&6&18240\\
 x_6&6&54240&y_6&2&4920
\end{array}
\]
For clarity, the lowest core weights and all minimizing absolute
charges used in this table are
\[
\begin{array}{c|c|l@{\qquad}c|c|l}
 &h_{A_L}&|j|&&h_{A_L}&|j|\\\hline
 x_0&0&0&y_0&2/5&0\\
 x_1&2/15&0&y_1&8/15&0,5\\
 x_2&2/5&0,3&y_2&4/5&0,2,3,5\\
 x_3&4/5&0,2,3,6&y_3&1/5&2\\
 x_4&1/3&6&y_4&11/15&1,2,4,6\\
 x_5&1&1,4,6&y_5&2/5&1,4\\
 x_6&4/5&4&y_6&1/5&4
\end{array}
\]
For each positive $j$ both signs occur. A convenient independent
check on every row is
\[
 Z=h_{A_L}d-L-Q/156-P,\qquad
 W=\frac L{26}+\frac{7Q}{936}+\frac{125P}{39}
                              +\frac{25}{26}h_{A_L}d.
\]
Complete right $\mathfrak{sl}_2$ characters imply
$y_2\ge x_5$ and $y_4\ge x_5+x_6$: these compare charges $10,12$
and charges $7,9$, respectively, with the common left doublet removed
in the second comparison. Direct comparison of the displayed table gives
\begin{align}
 4680\bigl(\tfrac{17}{10}D-W_{\rm total}\bigr)
 ={}&8460(y_2-x_5)+28716(y_4-x_5-x_6)\notag\\
 &+7956x_0+7356x_1+29340x_2+14964x_3+10392x_4\notag\\
 &+22212x_6+156y_0+18732y_1+11832y_3\notag\\
 &+1188y_4+29496y_5+10992y_6\ge0.
 \label{eq:sharp-top-bound}
\end{align}
Since $W_{\rm total}=hD$, this proves $0\le h\le17/10<2$ for
\emph{every} simple $U$-module, without assuming it is unitarizable.

\begin{reviewnote}{R2: scope of the core and finite-top certificates}
The charge minimizations and the nonnegative identity
\eqref{eq:sharp-top-bound} are specified above; their finite exact
check is
\artifact{hi_cft_20260923_next_structure_order13_top_check.json}.
The embedded lattice--Potts core and the spectral field witness remain
construction inputs, supplied by the exact core-extension certificate.
The finite-top irreducibility below is likewise imported from exact
field actions, particularly
\artifact{hi_cft_20260923_next_81top_all.json}; it should not be replaced
by numerical cyclicity alone.
\end{reviewnote}

The strongly rational restriction $B\downarrow_U$ is semisimple. Its
simple summands have integer lowest weights, hence only weights $0,1$
are possible. The vacuum occurs once, while the complementary current
space has dimension $86-6=80$. Its certified irreducible zero-mode
action therefore implies
\begin{equation}
 B\downarrow_U=U\oplus X,\qquad h_X=1,\qquad\dim X_{\rm top}=80.
 \label{eq:B-branching}
\end{equation}
The four mixed nontrivial $E_6\oplus A_2$ lattice sectors start at
weight $1$, with $81$-dimensional tops. Their next possible top weight
would be $2$, excluded by \eqref{eq:sharp-top-bound}. Their exact
top-action certificates use a one-dimensional Cartan eigenline,
positive-metric adjoint identities, and a cyclic span of dimension
$81$. The adjoint identities make the generated operator algebra a
star algebra; orthogonal complements are therefore invariant, and
the unique eigenline plus cyclicity proves irreducibility.

Write the actual lattice extension as $H=B\oplus J\oplus J^2$. Its
restriction is consequently
\begin{equation}
 H\downarrow_U=\mathbf1\oplus X\oplus Y\oplus Y^*,\qquad
 h_Y=1,\quad\dim Y_{\rm top}=81.
 \label{eq:H-branching-fork}
\end{equation}
At this stage $Y\cong Y^*$ has not been assumed. The three types
$\mathbf1,X,Y$ are distinguished by their lowest spaces.

\subsection{The rank-four or rank-six alternative}

Put $\mathcal C=\operatorname{Rep}(U)$ and
$\mathcal F=\mathcal C_H$, the category of \emph{all} right modules
over the extension algebra $H$. The strict positivity of nonvacuum
lowest weights gives the positive spherical structure through character
asymptotics. Since $H$ is holomorphic, it is Lagrangian and the
center correspondence \cite[Section~4]{DMNO}
gives $\mathcal C\simeq Z(\mathcal F)$.
Write $s=\operatorname{FPdim}_{\mathcal C}B$. Extension dimension
identities give
\[
 d_X=s-1,\qquad d_Y=d_{Y^*}=s,\qquad
 \operatorname{Dim}\mathcal C=9s^2,\quad
 \operatorname{Dim}\mathcal F=3s.
\]
Pseudo-unitarity of $\mathcal F$ follows from that of its center and
the squared global-dimension and Frobenius--Perron-dimension identities.
No $C^*$ structure is needed in this section. All categorical dimensions
below use the positive spherical structure and therefore equal
Frobenius--Perron dimensions.

For center induction $I$ and forgetful functor $F$, one has $H=I(\mathbf1)$
and $FI(\mathbf1)=\bigoplus_{a\in\operatorname{Irr}\mathcal F}a\otimes a^*$.
Adjunction therefore gives
\begin{equation}
 \operatorname{rank}\mathcal F
 =\dim\operatorname{End}_{\mathcal C}(H)
 =\begin{cases}4,&Y\not\cong Y^*,\\6,&Y\cong Y^*.
 \end{cases}
 \label{eq:rank-fork}
\end{equation}
\begin{lemma}[Automorphisms give invertible boundary objects]
\label{lem:automorphisms-invertibles}
For a pseudo-unitary fusion category $\mathcal F$, algebra
automorphisms of its canonical Lagrangian algebra $I(\mathbf1)$ inject
into $\operatorname{Inv}(\mathcal F)$.
\end{lemma}
\begin{proof}
Normalize the special Frobenius structure so that
$m\Delta=\mathrm{id}$ and set $f*g=m(f\otimes g)\Delta$.
If $f$ is an algebra automorphism, then $f*f=f$, and the same holds
for its composition inverse. Under the two-product description
\cite[Lemma~3.1 and Corollary~3.4]{BJ}, composition is fusion-ring
multiplication and the primitive convolution idempotents are
$d_a[a]$. Thus both $f$ and $f^{-1}$ have expansions with coefficients
$0$ or $d_a>0$. Their fusion product is $[\mathbf1]$.
Nonnegative fusion coefficients rule out cancellation. Every pair of
objects in the two supports must therefore multiply to copies of the
unit; rigidity then forces each pair to consist of inverse invertible
objects. The two supports are singletons and their dimensions are
one. Composition preserves this assignment, proving the injection.
\end{proof}

The actual order-three automorphisms of $H/B$ consequently inject into
$\operatorname{Inv}(\mathcal F)$. This injection is already spherical:
the proof used positivity of dimensions and fusion coefficients,
without using adjoints of morphisms.

\begin{lemma}[No sign extension]\label{lem:no-sign-extension}
The nonvacuum summand $X$ of $B\downarrow_U=U\oplus X$ is not
invertible.
\end{lemma}
\begin{proof}
If $X$ were invertible, self-duality of $B$ and connectedness would
give $X^2=\mathbf1$. The extension would have its $\mathbb Z/2$
grading, and its sign automorphism $g$ would satisfy $B^g=U$.
On currents this would give
$(E_6\oplus A_2)^g=U_1\cong\mathfrak{sl}_2\oplus\mathfrak{sl}_2$,
of rank two.

Every finite-order automorphism of a complex simple Lie algebra is
conjugate to a toral automorphism times a diagram automorphism;
see \cite[Section~5.2]{RLYG}. It therefore fixes a toral subalgebra
whose dimension is the number of orbits of the diagram automorphism
on simple roots. In type $E_6$ this number is at least four, and in
type $A_2$ at least one. Since the nonisomorphic ideals $E_6,A_2$
are preserved separately, $(E_6\oplus A_2)^g$ contains a toral
subalgebra of dimension at least five. Such a subalgebra cannot lie
in the rank-two Lie algebra $U_1$, a contradiction.
\end{proof}

If rank six had six invertibles, then $3s=6$ and
$d_X=s-1=1$, contradicting Lemma~\ref{lem:no-sign-extension}.
Thus there are
exactly three invertibles in either rank. Their associator is also
determined: the local $B$-module category has metric group
\[
 (\mathbb Z/3)^2,\qquad q(a,b)=\frac{2a^2+b^2}{3}\pmod1.
\]
The actual extension $H/B$ selects the diagonal isotropic line
$L_+=\{(a,a)\}$. Its complement
$L_-=\{(a,-a)\}$ is also isotropic, and
$(\mathbb Z/3)^2=L_+\oplus L_-$.
The pointed subcategory on $L_-$ is symmetric with trivial twists,
hence Tannakian; as a tensor category it is the untwisted
$\operatorname{Vec}_{\mathbb Z/3}$. Free $H/B$-module induction is
tensor on this subcategory. It is fully faithful since
\[
 \operatorname{Hom}_{H/B}(x\otimes(H/B),y\otimes(H/B))
 \cong\operatorname{Hom}_{\operatorname{Rep}(B)}
       (x,y\otimes(H/B))
\]
has dimension $\delta_{x,y}$ for $x,y\in L_-$: distinct elements of
$L_-$ represent distinct cosets of $L_+$. Every free module is a
direct sum of these three simple modules, and separability makes
every module a summand of a free module. Thus induction identifies the entire
all-module category for $H/B$ with this untwisted pointed category.
Viewing these modules in $\operatorname{Rep}(U)$ forgets only the
locality condition for the intermediate algebra $B$; it preserves
$H$-module morphisms and relative tensor products. This gives the
required full tensor embedding into $\mathcal F$.
Merely knowing the group of invertibles would not determine this
associator.

\subsection{Excluding rank four}

At rank four the fusion ring is near-group:
\[
 a^3=1,\qquad a\rho=\rho a=\rho,\qquad
 \rho^2=1+a+a^2+m\rho.
\]
There are exactly two self-dual simples, $\mathbf1$ and $\rho$.
Consequently \cite[Theorem~1.1]{Larson}, which assumes
pseudo-unitarity, restricts $m$ to $0,2,3,6$.
For $m=0$, $s=2$ contradicts Lemma~\ref{lem:no-sign-extension}.
For $m=2$, the categorical reduction
\cite[Theorem~7.1 and Corollary~7.4]{EGO} supplies the twisted $A_4$ double.
Its electric regular algebra gives a holomorphic $c=8$ extension
$H'$ with $U=(H')^{A_4}$ by the extension correspondence
\cite{HKL}. The nonvacuum lowest weights are positive, so $H'$ is
of CFT type, and \cite[Theorem~1]{DM} identifies it with $V_{E_8}$.
For completeness, the nontrivial order-two fixed Lie algebras are
$D_8$ and $E_7+A_1$, with dimensions $120,136$ and adjoint traces
$-8,24$. At order three the possibilities are
$A_8$, $E_6+A_2$, $D_7+\mathbb C$, and $E_7+\mathbb C$,
with dimensions $80,86,92,134$ and traces $-4,5,14,77$.
These lists follow by selecting Kac coordinates of total mark two
or three in the affine $E_8$ diagram, whose marks are
$1,2,2,3,3,4,4,5,6$. For order three the trace is
$(3\dim\mathfrak g^g-248)/2$, since the two nontrivial eigenspaces
are dual. Thus the two relevant traces are bounded below by $-8$
and $-4$. Averaging over $A_4$ gives
\[
 \dim(E_8^{A_4})\ge\frac{248+3(-8)+8(-4)}{12}=16,
\]
contradicting $\dim U_1=6$.

The remaining cases require a spherical induction argument before any
comparison with affine characters. We give this reduction explicitly.

\begin{lemma}[Near-group induction]\label{lem:near-group-induction}
Suppose $m=3k$, and put $d^2=3+3kd$, $s=2+kd$.
There are integers $p,q\ge0$, $p+q=k$, and a primitive cube root
$\zeta$ such that the known center objects have forgetful images
\[
\begin{array}{c|c|c}
 \text{object}&F(\text{object})&\text{twist}\\\hline
 X&1+k\rho&1\\
 Y,Y'&1+a+k\rho,\ 1+a^2+k\rho&1,1\\
 C&a+a^2+k\rho&\zeta\\
 D,D'&a+p\rho,\ a^2+p\rho&\zeta^2,\zeta^2\\
 E,E'&a+q\rho,\ a^2+q\rho&\zeta^2,\zeta^2
\end{array}
\]
Every further simple $Z_j$ forgets to $n_j\rho$, $n_j>0$. Its twist
$\theta_j$ satisfies
\begin{align}
 \sum_jn_j^2&=6+3k^2+4pq,\label{eq:near-group-length}\\
 S_1:=\sum_jn_j^2\theta_j&=k^2-kd-(p-q)^2\zeta^2,\notag\\
 \sum_jn_j^2\theta_j^2&=\overline{S_1}
                 +\varepsilon\sqrt{9k^2+12},
 \qquad\varepsilon=\nu_2(\rho)\in\{1,-1\}.
 \label{eq:near-group-moments}
\end{align}
\end{lemma}
\begin{proof}
The four one-dimensional representations of the fusion algebra have
formal codegrees $3s,9s/d^2,3,3$. Thus induction of the unit is
$1+X+Y+Y'$, with dimensions $1,d^2/3,s,s$ and trivial twists.
Since $FI(1)=4\cdot1+a+a^2+3k\rho$, adjunction and comparison of
rational and irrational dimension coefficients give the first two
rows. Also
\[
 FI(a)=1+4a+a^2+3k\rho,\qquad
 FI(a^2)=1+a+4a^2+3k\rho.
\]
The two inductions have norm four, share one constituent, and contain
$Y,Y'$ respectively. Their remaining norm-three parts are therefore
multiplicity-free, with a single common constituent $C$.
Write its $\rho$ multiplicity as $c$ and the other two as $p,q$.
Then $c+p+q=2k$.

The untwisted pointed subcategory gives $\nu_3(a)=1$, so positivity
in the twist-trace formula forces every constituent of $I(a)$ to
have cube-root twist. Its first indicator, separated over the
disjoint fields $\mathbb Q(d)$ and $\mathbb Q(\zeta)$, gives
\[
 2+2\theta_C+\theta_D+\theta_E=0,\qquad
 k+c\theta_C+p\theta_D+q\theta_E=0.
\]
The first equality forces $\theta_C=\zeta$ and
$\theta_D=\theta_E=\zeta^2$, up to simultaneous conjugation; the
second gives $c=k$ and $p+q=k$. Ribbon duality preserves the primed
twists. Subtracting the nine known squared dimensions from $9s^2$
gives \eqref{eq:near-group-length}. The first two twist traces of
$I(\rho)$ give \eqref{eq:near-group-moments}, using
$3/d=d-3k$ and $3s/d=\sqrt{9k^2+12}$.
\end{proof}

\begin{proposition}[The multiplicity-three modular representation]
\label{prop:m3-modular}
For $k=1$, the center has the modular data of
$G_{2,3}\otimes A_{2,1}$ or their simultaneous reversal.
\end{proposition}
\begin{proof}
One of $p,q$ is zero. The corresponding $D,D'$ are invertible and
generate a nondegenerate pointed category $\mathcal P$ with twists
$1,\zeta^2,\zeta^2$. Write
$Z(\mathcal F)=\mathcal M\boxtimes\mathcal P$ using its modular
centralizer. The known objects form three $\mathcal P$-orbits;
representatives $1,X,C$ in $\mathcal M$ have dimensions
$1,1+d,s$ and twists $1,1,\zeta$. Balancing verifies their trivial
monodromy with $D$ from $D X=E$ and $DC\in\{Y,Y'\}$.
The remaining free orbits, together with
\eqref{eq:near-group-length}, leave three simples in $\mathcal M$,
all of dimension $d$.

Take $\zeta=e^{2\pi i/3}$, the other choice conjugating the answer.
For their twists $z_1,z_2,z_3$, the two moments become
\[
 \sum_i z_i=\alpha:=\frac{-1-i\sqrt7}{2},\qquad
 \sum_i z_i^2=-\frac12-i(\varepsilon-\tfrac12)\sqrt7.
\]
The second sum would have absolute value four for $\varepsilon=-1$,
so $\varepsilon=1$. Since $\alpha^2+\alpha+2=0$, Newton's identity
and the unit modulus of all $z_i$ give elementary symmetric functions
$e_2=\overline\alpha$, $e_3=1$. Thus the twists are the roots of
$t^3-\alpha t^2+\overline\alpha t-1$, namely
$\zeta_7^3,\zeta_7^5,\zeta_7^6$.

The actual extension algebra $1+X$ lies in $\mathcal M$. Its local
module category has dimension three and is the pointed category with
nontrivial twist $\zeta$, fixed by its multiplicative central charge.
Its two nonvacuum simples must both restrict to $C$: each underlying
object has dimension $s$, cannot contain the unit by induction
reciprocity, and no other combination of the six known dimensions
has these properties. Hence the branching columns are $1+X,C,C$.
Modular intertwining for this extension
\cite[Theorems~4.1 and~4.5]{KO} fixes the first three rows
of $\widehat S=\sqrt3sS_{\mathcal M}$:
\[
 \begin{pmatrix}
 1&1+d&s&d&d&d\\
 1+d&1&s&-d&-d&-d\\
 s&s&-s&0&0&0
 \end{pmatrix}.
\]
All six objects are self-dual, the last three by their distinct
twists. Their fusion matrices are real symmetric, so Verlinde
diagonalization and the positive vacuum row make $S_{\mathcal M}$
real. For the remaining block $K$, orthogonality and the modular
relation yield
\[
 \sum_jK_{ij}=d,\qquad
 \sum_jK_{ij}z_j=d+i\sqrt3s\,z_i^{-1}.
\]
These determine each real row: three distinct points on a circle are
not collinear. The affine $G_{2,3}$ data
\cite[Table~4 and equation~(3.8)]{AGJS} realize these rows and twists
in the reversed orientation; conjugating gives the other possibility.
Thus the entire representation, rather than merely its twist list,
is forced.
\end{proof}

\begin{lemma}[The balanced multiplicity-six pattern is impossible]
\label{lem:m6-balanced}
In Lemma~\ref{lem:near-group-induction} with $k=2$, one has
$\{p,q\}=\{0,2\}$.
\end{lemma}
\begin{proof}
Suppose $p=q=1$, and put $d=3+2\sqrt3$, $r=1+d$.
The known dimensions are $1$, $1+2d$, three copies of $2r$, and four
copies of $r$. The remaining dimensions are $n_jd$, with
$\sum_jn_j^2=22$. Write
$u=(2-\sqrt3)/12$, $v=(2+\sqrt3)/12$.
The normalized vacuum row is
\[
 (u,v,1/3,1/3,1/3,1/6,1/6,1/6,1/6,n_j\sqrt3/12).
\]
Signed Galois symmetry \cite[Theorem~II(iii)]{DLN}, for an
automorphism reversing $\sqrt3$, sends this row to the $X$ row:
its first entry becomes $v>0$, uniquely identifying $X$ and a
positive sign. The forced row is
\[
 (v,u,1/3,1/3,1/3,1/6,1/6,1/6,1/6,-n_j\sqrt3/12).
\]
The object $X$ is self-dual because its dimension is unique.
Verlinde's formula consequently gives
\[
 N_{XX}^{X}=\frac{v^3}{u}+\frac{u^3}{v}
              +\frac39+\frac4{36}-\frac{22}{48}
            =\frac{97}{72}+\frac{32}{72}-\frac{33}{72}
            =\frac43,
\]
contradicting integrality. Here $u+v=1/3$, $uv=1/144$ verify the
first two terms directly.
\end{proof}

\begin{lemma}[A universal multiplicity-six character projection]
\label{lem:m6-projection}
The unbalanced case has a nondegenerate pointed factor $\mathcal P$
with twists $1,\zeta^2,\zeta^2$. In its modular centralizer there is
a character projection with matrices
\begin{equation}
 S_v=-\frac12\begin{pmatrix}1&2&1\\1&0&-1\\1&-2&1\end{pmatrix},
 \qquad T_v=\operatorname{diag}(1,-i,-1),
 \label{eq:m6-projection}
\end{equation}
before the common central-charge scalar, or the conjugate matrices.
\end{lemma}
\begin{proof}
The zero multiplicity gives the pointed factor just as for $k=1$.
In the centralizer the known simples have dimensions
$1,2d+1,2d+2$ and twists $1,1,\zeta$; all others have dimensions
$n_jd$, with $\sum n_j^2=6$.
The Gauss sum and the second indicator give
\[
 \sum n_j^2z_j=-2-4i,\qquad
 \sum n_j^2z_j^2=-2+4i(1-\varepsilon).
\]
The second equality excludes $\varepsilon=-1$ by absolute values.
Since $(1+\Re z)(1+\Im z)\ge0$ on the unit circle, summing gives
$6-2-4+\tfrac12\Im(-2)=0$. Thus each $z_j$ is $-1$ or $-i$,
with total squared multiplicity two or four, respectively.

Put
\[
 v_0=\chi_1-\chi_X,\qquad
 v_A=-\frac12\sum_{z_j=-i}n_j\chi_j,\qquad
 v_B=-\sum_{z_j=-1}n_j\chi_j.
\]
The vacuum row is $(u,v,1/\sqrt3,n_j/4)$, with
$u=(2\sqrt3-3)/12$, $v=(2\sqrt3+3)/12$; Galois symmetry gives
the $X$ row $(v,u,1/\sqrt3,-n_j/4)$.
Subtracting these rows yields
$Sv_0=-v_0/2-v_A-v_B/2$. All three combinations are fixed by
duality. Applying $S^2=1$ and
$STS=iT^{-1}ST^{-1}$ to $v_0$ then determines
\[
 Sv_A=-v_0/2+v_B/2,\qquad
 Sv_B=-v_0/2+v_A-v_B/2,
\]
which is \eqref{eq:m6-projection}. This proof permits either one
dimension-$2d$ object or four dimension-$d$ objects at twist $-i$;
no unproved rank-nine classification has been used.
\end{proof}

\begin{lemma}[Pole cancellation]\label{lem:pole-cancellation}
Let $\sigma$ be a finite-image representation of
$\mathrm{SL}_2(\mathbb Z)$ with no invariant vector. Two weight-zero
vector-valued modular functions, holomorphic on the upper half-plane,
with representation $\sigma$ and equal principal parts at infinity
are equal.
\end{lemma}
\begin{proof}
Their difference has no pole at infinity. Covariance makes every
component holomorphic at every cusp of the finite-index kernel of
$\sigma$. It descends to a holomorphic function on the compactified
modular curve and is constant. The resulting constant vector must be
invariant under $\sigma$, hence zero.
\end{proof}

\begin{proposition}[Current obstructions for the remaining near groups]
\label{prop:near-group-current-obstructions}
Under the hypotheses above, $m=3$ forces
\[
 \dim U_1=15+A+10B\quad\text{or}\quad14+42A,
 \qquad A,B\in\mathbb Z_{\ge0},
\]
whereas $m=6$ forces
\[
 \dim U_1=8+A\quad\text{or}\quad5+12P,
 \qquad A,P\in\mathbb Z_{\ge0}.
\]
In particular none allows $\dim U_1=6$.
\end{proposition}
\begin{proof}
We give the theta normalizations, since a factor of two in the vacuum
coordinate would change these obstructions. For $m=3$, use the affine
$G_{2,3}$ labels $1,7,27,77,14,64$, and for each $A_{2,1}$ label put
\[
 v=(\chi_1-\chi_{64},-\chi_{77},-\chi_{14},-\chi_7).
\]
The complement of the extension character is the even Fourier block
on $\mathbb Z/7$, with entries
\[
 S_{00}=-1/\sqrt7,\quad S_{0r}=-2/\sqrt7,\quad
 S_{r0}=-1/\sqrt7,\quad
 S_{rs}=-2\cos(4\pi rs/7)/\sqrt7,
\]
and twist exponents $r^2/7$, $0\le r\le3$.
Together with the even $A_2$ block, multiplication by $\eta^8$
identifies it with the Weil representation of a positive rank-four
lattice with discriminant form
$\mathbb Z_7(r^2/7)\oplus\mathbb Z_3(s^2/3)$.
The actual extension gives $\chi_1+\chi_{64}=\chi_B$. Thus a
difference of two such character vectors has vacuum projection
$\delta v_0=2\delta\chi_U$.

Take $L=L_7\oplus A_2$, where $L_7$ has Gram matrix
$\left(\begin{smallmatrix}2&1\\1&4\end{smallmatrix}\right)$.
For a basis $u,v$ of $L_7$, write $w=v-u/2$ and $x=pu+tw$.
Set $a=p^2$, $b=7t^2/4$, and $c=(y,y)/2$ on the $A_2$ factor.
The quadratics $H_1=a-b$ and $H_2=a+b-c$ are harmonic.
Poisson summation gives vector-valued theta series of weight four.
Their required coefficients are
\[
\begin{array}{c|c|r|r}
 (r,s)&q\text{-power}&\Theta_{H_1}&\Theta_{H_2}\\\hline
 (1,0)&1/7&-1/7&1/7\\
 (3,0)&2/7&3/7&4/7\\
 (0,0)&1&2&-4
\end{array}
\]
The first row comes from $2w/7$, the second from the two shortest
vectors with $a=1/4,b=1/28$, and the last from two $L_7$ roots and
six $A_2$ roots. The only nonvacuum polar top dimensions are $A$ at
$2/7$ and $B$ at $1/7$, with zero allowed for an integer-shifted top.
Subtract the affine reference, which has $22$ currents and $A=7,B=0$.
Pole matching gives
\[
 \eta^8\delta v=\lambda\Theta_{H_1}+\mu\Theta_{H_2},\qquad
 \lambda=7-A+4B,\quad\mu=7-A-3B.
\]
This identity follows from Lemma~\ref{lem:pole-cancellation}: after
division by $\eta^8$, the remainder has no poles; its only possible
$T$-fixed line is $(r,s)=(0,1)$, which the Fourier $S$ matrix does
not preserve. The vacuum coefficient therefore gives
$2(\dim U_1-22)=2\lambda-4\mu$, the first formula.

For the reverse orientation, realize the negative discriminant form
by the positive lattice with Gram matrix
\[
 G_7^-=
 \begin{pmatrix}2&-1&0&0\\-1&2&-1&0\\
 0&-1&2&-1\\0&0&-1&6\end{pmatrix}.
\]
Its dual generator $\ell=(16,11,6,1)/21$ has energy $8/21$;
$(r,s)\mapsto(9r+7s)\ell$ identifies the discriminant form.
Project the theta series to functions even separately in $r,s$.
Let $H=\|x_{A_3}\|^2-3\|x_{A_3^\perp}\|^2$ and
$F=-6D_2\Theta$, where $D_2=q\,d/dq-E_2/6$ is the modular derivative.
The necessary coefficients after projection are
\[
\begin{array}{c|c|r|r}
 (r,s)&q\text{-power}&F&\Theta_H\\\hline
 (0,0)&0&1&0\\
 (0,0)&1&-84&24\\
 (2,1)&2/21&3/14&-2/7
\end{array}
\]
There are twelve roots, all in $A_3$. The polar coset has shortest
vector $(1,2,3,4)/21$, perpendicular to $A_3$; the separately even
projection halves its coefficient, since the paired coset has no
vector of the same norm. If $A$ is its actual polar top dimension,
matching the coefficient $-A$ gives
$\eta^8v=F+(3+14A)\Theta_H/4$. Here $T$ has no fixed vector, so
pole cancellation again proves equality. Since
$\eta^8\chi_B=1+78q+\cdots$, its vacuum coefficient is
$2\dim U_1-94$. Thus
$2\dim U_1-94=-84+6(3+14A)$, giving $14+42A$.

For $m=6$, tensor the projection in
Lemma~\ref{lem:m6-projection} with its pointed factor and first take
the orientation with twists $1,i,-1$. After multiplication by
$\eta^8$ this is the Weil representation for
$L=A_1\oplus A_1\oplus A_2$, restricted by interchange of the two
$A_1$ factors and conjugation of the $A_2$ charge.
The harmonic quadratic
$H=\|x_{A_1^2}\|^2/2-\|y_{A_2}\|^2/2$ has coefficients
\[
 [q^{1/4}]\Theta_{H,(10,0)}=\frac12,
 \qquad [q]\Theta_{H,(00,0)}=-2.
\]
The first counts the two half-roots in a nontrivial $A_1$ coset;
the second counts four $A_1$ roots and six $A_2$ roots.
Let $A=\sum n_j\dim(M_j)_{\rm top}$ over the quarter-twist polar
modules. This is a nonnegative integer, and its projected polar
coefficient is $-A/2$. Use as reference the order-three simple-current
extension of $A_{2,9}$ tensored with $A_{2,1}$: it has $16$ currents
and $A=8$. Its projected representation agrees with
\eqref{eq:m6-projection}; its matrices are also given in
\cite[Lemma~3.4 and Proposition~3.5]{Yu}. Thus cancellation gives
\[
 \eta^8\delta v=(8-A)\Theta_H,
 \qquad 2(\dim U_1-16)=-2(8-A).
\]
As in the first orientation above, the sole $T$-fixed coordinate is
not $S$-invariant. This gives $\dim U_1=8+A$.

For the reverse orientation use the positive lattice with Gram matrix
\[
 G_6^-=
 \begin{pmatrix}2&-1&-1&-1\\-1&2&0&0\\
 -1&0&2&0\\-1&0&0&4\end{pmatrix}.
\]
Its first three basis vectors span $A_3$; if the basis is $a,b,c,e$,
then $w=e+a+(b+c)/2$ is perpendicular to $A_3$ and has norm three.
The discriminant generators $(b+e)/2,(c+e)/2,2w/3$ have quadratic
values $3/4,3/4,2/3$ and mutually integral pairings. The polar coset
is represented by $-w/3$, with energy $1/6$. This shortest vector
is unique, and the projection does not halve its coefficient.
With $H=\|x_{A_3}\|^2-3\|x_{A_3^\perp}\|^2$ and $F=-6D_2\Theta$,
\[
\begin{array}{c|c|r|r}
 \text{coordinate}&q\text{-power}&F&\Theta_H\\\hline
 \text{vacuum}&0&1&0\\
 \text{vacuum}&1&-84&24\\
 \text{polar}&1/6&0&-1
\end{array}
\]
Let $P$ be the weighted top dimension in the half-twist coordinate
with nonvacuum pointed charge; the weighting is by the integers $n_j$.
Contragredience equates the two pointed-charge sums, so averaging
does not change this integer. Pole cancellation gives
$\eta^8v=F+P\Theta_H$: there is no $T$-fixed coordinate in this
orientation. Comparing the vacuum coefficient yields
$2\dim U_1-94=-84+24P$, hence $\dim U_1=5+12P$.
\end{proof}

The representations used in the last proof have finite image because
they are Weil representations, with the finite-order eta multiplier
removed. Thus the proof does not assume a dimension formula for a
space of vector-valued modular forms. The coefficients $A,B,P$ are
weighted dimensions of actual polar tops, and distinct modules need
not have linearly independent characters.

\begin{remark}[Scope of the rank-four exclusion]
The universal induction, balanced-pattern obstruction, weighted
projection, and theta coefficient normalizations are printed above.
Their source calculations are
\artifact{hi_cft_20260923_next_m3_universal_data.md},
\artifact{hi_cft_20260923_next_structure_m6_universal_projection.md},
\artifact{hi_cft_20260923_next_structure_m3_character_obstruction.md},
and \artifact{hi_cft_20260923_next_structure_m6_character_obstruction.md}.
Lemma~\ref{lem:no-sign-extension} removes the formerly used pointwise
stabilizer calculation $\operatorname{Aut}(B/U)=1$ from this proof.
No classification restricted to $C^*$ categories is used here.
\end{remark}

\subsection{The rank-six ring and exclusion of higher multiplicity}

Rank four being excluded, \eqref{eq:rank-fork} gives
$H=\mathbf1\oplus X\oplus2Y$. The composition algebra of
$\operatorname{End}_{\mathcal C}(H)$, identified with the complex
fusion algebra by \cite[Corollary~3.4]{BJ}, is
\[
 K_0(\mathcal F)\otimes\mathbb C
 \cong\mathbb C\oplus\mathbb C\oplus M_2(\mathbb C),
\]
In particular it is noncommutative. Let $a$ generate the three
invertibles. A permutation of order three on three noninvertibles
is either transitive or trivial. If the left action were trivial,
duality would make the right action trivial as well. The two nontrivial
character idempotents of $\mathbb C[\mathbb Z/3]$ would then be
one-dimensional direct ideals. Their four-dimensional complementary
ideal would have to be $M_2(\mathbb C)$, but the dimension character
is nonzero on that ideal, a contradiction.

Choose a noninvertible $\rho$. The right action commutes with the free
left cyclic action, so it is either the same cycle or its inverse.
The same cycle would make the fusion ring generated by two commuting
elements $a,\rho$, which is impossible. Thus $\rho a=a^2\rho$.
If $\rho^*=a^k\rho$, then
$(a^i\rho)^*=a^{i+k}\rho$; involutivity gives $2k=0$ in
$\mathbb Z/3$, hence $k=0$. All three noninvertibles are self-dual.
Frobenius reciprocity gives one copy of the unit and no other
invertible in $\rho^2$. Write its remaining coefficients as
$m_0,m_1,m_2$. The identity $\rho^2a=a\rho^2$ cyclically equates
these coefficients, so
\begin{equation}
 a^3=1,\qquad\rho a=a^2\rho,\qquad
 \rho^2=1+m(\rho+a\rho+a^2\rho),\qquad m\ge1.
 \label{eq:rank6-ring}
\end{equation}

The following induction calculation determines the multiplicity.
Put
\[
 d=\frac{3m+\sqrt{9m^2+4}}2,\quad s=1+d^2,\quad
 \Delta=9m^2+4,\quad R=\rho+a\rho+a^2\rho.
\]
\begin{lemma}[Center induction and two moments]
\label{lem:rank6-induction}
For a pseudo-unitary categorification of \eqref{eq:rank6-ring} whose
pointed subcategory is untwisted, center induction satisfies
\[
 I(1)=1+X+2Y,\quad d_X=d^2,\quad d_Y=s,
 \qquad I(a)=I(a^2)=C_0+C_1+C_2,
\]
where $\theta_X=\theta_Y=1$, $d_{C_j}=s$, and the $C_j$ twists are
$1,\zeta_3,\zeta_3^2$. Every other central simple $Z_j$ satisfies
$F(Z_j)=n_jR$, with $n_j\in\mathbb Z_{>0}$, and
\begin{align}
 \sum_jn_j^2&=4m^2+2,\notag\\
 \sum_jn_j^2\theta_j&=-\frac{m^2+m\sqrt\Delta}{2},
 \label{eq:indicator-first}\\
 \sum_jn_j^2\theta_j^2
 &=\frac{-m^2+(2\varepsilon-m)\sqrt\Delta}{2},
 \qquad\varepsilon=\nu_2(\rho)\in\{1,-1\}.
 \label{eq:indicator-second}
\end{align}
\end{lemma}
\begin{proof}
The complex fusion algebra has two characters
$a\mapsto1$, $\rho\mapsto d,-d^{-1}$, and the two-dimensional
representation
\[
 a\longmapsto\begin{pmatrix}\zeta_3&0\\0&\zeta_3^2\end{pmatrix},
 \qquad \rho\longmapsto\begin{pmatrix}0&1\\1&0\end{pmatrix}.
\]
Their formal codegrees are respectively $3s,3s/d^2,3$.
The induction--formal-codegree correspondence
\cite[Theorems~2.5 and~2.13]{Ostrik13} gives
$I(1)=1+X+2Y$, with the stated dimensions and trivial twists.
Conjugation by $a$ cycles the three noninvertibles, so their
multiplicities in the forgetful image of a central object coincide.
Since
\[
 FI(1)=6\cdot1+3mR,
\]
adjunction and dimensions give $F(X)=1+mR$ and $F(Y)=2\cdot1+mR$.

Similarly,
\[
 FI(a)=FI(a^2)=3a+3a^2+3mR.
\]
Thus the two induced objects have self-inner-products and mutual
inner-product equal to three. They are equal, multiplicity-free sums
of three simples $C_j$, disjoint from $I(1)$. Adjunction gives
$F(C_j)=a+a^2+b_jR$, where $b_j\ge0$ and $\sum_jb_j=3m$.
Untwistedness gives $\nu_3(a)=1$. The center-twist formula
\cite[Theorem~4.1]{NS} therefore reads
\[
 \sum_jd_{C_j}\theta_{C_j}^3=\operatorname{Dim}\mathcal F
                         =\sum_jd_{C_j}.
\]
Positive dimensions and equality in the triangle inequality imply
$\theta_{C_j}^3=1$. The first indicator of $a$ is zero, so
\[
 \sum_j(2+3b_jd)\theta_{C_j}=0.
\]
The real quadratic field $\mathbb Q(d)$ is disjoint from
$\mathbb Q(\zeta_3)$. Equating its two coefficients shows that the
three twists are $1,\zeta_3,\zeta_3^2$ once each and $b_j=m$ for all
$j$. All invertible constituents have now been exhausted by induction
of $1,a,a^2$; every remaining forgetful image is $n_jR$.

The center dimension identity is
\[
 9s^2=1+d^4+4s^2+9d^2\sum_j n_j^2,
\]
which simplifies to the first asserted equation. Finally adjunction
gives
\[
 I(\rho)=mX+mY+m(C_0+C_1+C_2)+\sum_jn_jZ_j.
\]
Its first and second twist traces are
$\operatorname{Dim}\mathcal F\,\nu_1(\rho)=0$ and
$\operatorname{Dim}\mathcal F\,\nu_2(\rho)=3s\varepsilon$.
Using $d^{-1}=d-3m$ gives \eqref{eq:indicator-first} and
\eqref{eq:indicator-second}.
\end{proof}

The roots in these equations form a multiset of exactly $4m^2+2$
elements, with each $\theta_j$ repeated $n_j^2$ times. No distinctness
of the twists has been assumed.

Write $\Delta=t^2D$ with $D$ squarefree. Since
$(3m)^2<\Delta<(3m+1)^2$, we have $D>1$. Reduction modulo three and
seven gives $D\equiv1\pmod3$ and $7\nmid D$: the latter uses that
$m^2\equiv5\pmod7$ has no solution.
Except when $D=10$,
\begin{equation}
 \frac{\varphi(2D)}{\sqrt D}\ge\frac{12}{\sqrt{13}}.
 \label{eq:squareclass-bound}
\end{equation}
Indeed for odd $D$ the left side is the product of
$(p-1)/\sqrt p$ over its prime factors. A single prime must be at
least $13$, and two distinct primes contribute at least
$40/\sqrt{55}>12/\sqrt{13}$; all further factors exceed one.
For even $D$ there is an additional factor $\sqrt2$.
A single odd prime is $5$ (the exceptional $D=10$) or at least $11$,
and two odd primes again give a larger product.
Apply \cite[Theorem~1.3]{Larson} to two copies of the root multiset in
\eqref{eq:indicator-first}, to remove any half-integral coefficients.
Together with \eqref{eq:squareclass-bound} this gives
\[
 4m^2+2\ge\frac{mt}{2}\varphi(2D)
 \ge\frac{6m\sqrt{9m^2+4}}{\sqrt{13}}.
\]
For $m>1$ this is impossible, since
\[
 36m^2(9m^2+4)-13(4m^2+2)^2
 =4(m^2-1)(29m^2+13)>0.
\]
The exceptional squareclass requires the following equality statement.

\begin{lemma}[The conductor-forty bound]\label{lem:conductor40}
If a multiset of $N$ roots of unity has sum $-a-b\sqrt{10}$, where
$a,b\in\mathbb Z_{\ge0}$, then $N\ge a+8b$. If equality holds and
the sum of the squared roots lies in $\mathbb Q(\sqrt{10})$, that
second sum is $a$.
\end{lemma}
\begin{proof}
Choose a cyclotomic field containing all roots and
$\mathbb Q(\sqrt{10})$. Take the union of the multiset over its Galois
group $G$ fixing $\sqrt{10}$, and put $M=|G|$. This gives a multiset
of $MN$ roots, with sum $-Ma-Mb\sqrt{10}$, whose multiplicities are
constant on $G$-orbits. Multiplicities remain integral throughout.

If an orbit consists of primitive $n$th roots and $40\nmid n$, its
sum is the Ramanujan sum $\mu(n)\in\{0,1,-1\}$ and its size is
$\varphi(n)$. Its rational contribution therefore costs at least one
root per unit. A negative contribution realizes equality only for
the singleton orbit $\{-1\}$.

If $40\mid n$, the primitive roots split into two equal orbits,
distinguished by the primitive quadratic character of conductor $40$.
Their sums are
\[
 \frac{\mu(n)\pm\tau_n(\chi_{40})}{2}.
\]
Here $\mu(n)=0$, and the induced Gauss sum vanishes unless
$n=40u$ with $u$ squarefree and coprime to $40$. In that case
\[
 \tau_n(\chi_{40})=\mu(u)\chi_{40}(u)\sqrt{40}.
\]
This formula follows by separating the factors coprime to $40$ in the
Chinese remainder theorem; an extra power of a prime dividing the
conductor gives zero. Thus an orbit contributes zero or
$\pm\sqrt{10}$, at cost at least eight roots. Equality for a nonzero
contribution requires $n=40$.

The rational and irrational contributions come from disjoint types
of orbit, so $MN\ge Ma+8Mb$. Equality leaves only $Ma$ copies of
$\{-1\}$ and $Mb$ copies of the negative quadratic orbit of primitive
fortieth roots. The latter orbit has squared sum zero: reduction of
its exponents modulo $20$ gives every unit once, since
$\chi_{40}(x+20)=-\chi_{40}(x)$, and the resulting sum is
$\mu(20)=0$. The averaged squared sum is therefore $Ma$.
By the additional hypothesis, Galois averaging multiplied the original
second sum by $M$. Dividing by $M$ gives the assertion.
\end{proof}

In our application $m,t$ are even and $a=m^2/2$, $b=mt/2$.
Lemma~\ref{lem:conductor40} gives
\[
 4m^2+2\ge\frac{m^2}{2}+4mt
\]
For $m\ge4$ one has $t/m>\sqrt{9/10}>15/16$, making the right side
strictly larger than $17m^2/4>4m^2+2$. The remaining $m=t=2$ is the
equality case,
so the squared sum must be $2$, whereas \eqref{eq:indicator-second}
gives $-2$ or $-2-4\sqrt{10}$. Thus $m=1$.
The same argument, including the indispensable second-moment
hypothesis, is recorded in
\artifact{hi_cft_20260923_next_structure_rank6_indicators.md} and
\artifact{hi_cft_20260923_next_rank6_root_audit.md}.

Consequently
\begin{equation}
 d=\frac{3+\sqrt{13}}2,\qquad
 s=\frac{13+3\sqrt{13}}2,\qquad
 d_X=\frac{11+3\sqrt{13}}2.
 \label{eq:haagerup-dimensions}
\end{equation}
\begin{lemma}[The six residual twists]\label{lem:six-residual-twists}
For $m=1$, there are six residual central simples, each with $n_j=1$.
Their twists are the six primitive thirteenth roots with quadratic
nonresidue exponents, and $\nu_2(\rho)=1$.
\end{lemma}
\begin{proof}
Equation~\eqref{eq:indicator-first} gives a multiset of six roots
with sum $(-1-\sqrt{13})/2$. Average it over a cyclotomic Galois group
fixing $\sqrt{13}$, as in Lemma~\ref{lem:conductor40}. An orbit with
nonzero $\sqrt{13}$ coefficient has coefficient of absolute value
$1/2$ and at least six members. To see the equality condition, the
induced quadratic Gauss-sum formula is nonzero only at primitive
orders $13u$, with $u$ squarefree and coprime to $13$. The orbit size
is $6\varphi(u)$, so equality permits only $u=1,2$.

The entire available length is already needed for the irrational
coefficient. Hence no rational-only orbit or cancellation can occur.
The two possible negative quadratic orbits, of orders $13$ and $26$,
have sums $(-1-\sqrt{13})/2$ and $(1-\sqrt{13})/2$, respectively.
The required rational coefficient forces every orbit to be the former.
Consequently every original root is one of the six nonresidue
thirteenth roots. These six roots are linearly independent over
$\mathbb Q$: a rational relation would give a polynomial of degree
at most twelve, with zero constant term, divisible by
$\Phi_{13}(t)=1+t+\cdots+t^{12}$, and hence zero.
Comparison with the quadratic period shows that each root occurs
once. As each object was repeated $n_j^2$ times, all six objects have
$n_j=1$ and distinct twists.

Multiplication by two sends nonresidues modulo thirteen to residues,
so the second sum is $(-1+\sqrt{13})/2$.
Equation~\eqref{eq:indicator-second} now gives $\varepsilon=1$.
\end{proof}

The center thus has twelve simples. Its dimensions and twist exponents are
\[
 \begin{array}{c|cccc}
 \text{number}&1&1&4&6\\\hline
 \text{dimension}&1&d^2&s&3d
 \end{array}
 \qquad
 0,0,0,0,\tfrac13,\tfrac23,
 \tfrac2{13},\tfrac5{13},\tfrac6{13},
 \tfrac7{13},\tfrac8{13},\tfrac{11}{13}.
\]
The complete $S,T$ data follow from
\cite[Theorem~A(2), equivalently Theorem~5.1]{GSY}: this statement assumes
a complex spherical untwisted HI category. The auxiliary metric group
has prime order $13$ and its negative quadratic Gauss sum fixes its
isometry class. This identifies the standard Haagerup modular data.
It does not by itself identify a braided category; that stronger
conclusion is treated separately below and in the analytic section.

\subsection{An optional algebraic classification route}

There is a second possible route from this fusion ring to an abstract
ribbon identification. We isolate its extra input rather than folding
an unreplayed classification into the preceding modular-data argument.
Vercleyen \cite[Section~4.4 and Table~9.1, p.~171]{Vercleyen} reports an
exhaustive classification of complex pivotal multiplicity-free fusion
categories of rank at most seven. Its list for $\operatorname{HI}(\mathbb
Z_3)$ has four entries. Gauge equivalence there permits complex nonzero
scalars and ring automorphisms; no unitarity assumption is imposed.
The solution method uses exact polynomial and algebraic-number
operations, with a separate numerical redundancy check. The complete
notebooks and decision logs are available on request rather than
included in the thesis \cite[Sections~4.1.7--4.1.8 and~4.4.1]{Vercleyen}.
We have not independently replayed that exhaustion calculation.

\begin{proposition}[Ribbon identification with an additional classification input]
\label{prop:optional-algebraic-identification}
Assume, in addition to the hypotheses used above, that the four-entry
list in \cite[Table~9.1]{Vercleyen} exhausts complex pivotal fusion
categories with the $\operatorname{HI}(\mathbb Z_3)$ fusion ring. Then
$\mathcal F$ is tensor equivalent to $\mathcal H_2$ or $\mathcal H_3$,
and there is an abstract ribbon equivalence
\[
 \operatorname{Rep}(U)\simeq Z(\mathcal H_3).
\]
In particular this modular category is abstractly unitarizable.
\end{proposition}
\begin{proof}
Evans--Gannon \cite[Theorem~3(ii)]{EG15} construct four inequivalent
pivotal categories with this ring. Two have noninvertible dimension
$(3+\sqrt{13})/2$ and are the unitary categories $\mathcal H_2$ and
$\mathcal H_3$; the other two have dimension $(3-\sqrt{13})/2$.
Their theorem classifies the endomorphism realizations considered
there; the additional exhaustion hypothesis is what makes this a
complete list for the present abstract $\mathcal F$.

A pivotal structure on any categorification of this ring is unique
when it exists. Indeed, the ratio of two such structures is a tensor
automorphism of the identity. Write its scalar on a simple $Z$ as
$t_Z$. Since $\rho^2$ contains both $\mathbf1$ and $\rho$,
$t_\rho^2=1=t_\rho$. It also contains $a\rho$ and $a^2\rho$,
so their scalars are one, and $a\otimes\rho=a\rho$ gives $t_a=1$.
Thus all six scalars are one. The negative-dimensional examples
cannot acquire a positive spherical structure by repivotalization.
The positive spherical category $\mathcal F$ is therefore one of the
two positive examples.

The categories $\mathcal H_2$ and $\mathcal H_3$ are Morita equivalent
by \cite[Theorem~1.4]{GS}, so their centers are braided equivalent.
Combining this with $\operatorname{Rep}(U)\simeq Z(\mathcal F)$ gives
the claimed braided equivalence. Both sides have the positive
spherical structures, which are unique; the induced pivotal and
twist structures are consequently preserved. This makes the
equivalence ribbon. Transporting the unitary structure from the
Haagerup center proves abstract unitarizability.
\end{proof}

This proposition does not establish positivity of the canonical VOA
fusion pairings or the analytic realization of the transported adjoint.
Those are separate complete-unitarity assertions addressed later.
Neither this classification input nor the proposition is needed for
the character computation that follows.

\section{The complete character vector}
\label{sec:characters}

Let $q=e^{2\pi i\tau}$. Character exponents are $h-c/24=h-1/3$.
Using the preceding twist list and strict positivity of nonvacuum
lowest weights, the only allowed nonvacuum pole is
$\gamma q^{-7/39}$ in the twist-$2/13$ coordinate. Here
$\gamma\in\mathbb Z_{\ge0}$ is its top dimension if that module starts
at $2/13$, and is zero if its lowest weight is shifted upward. The
vacuum pole is $q^{-1/3}$ with coefficient one.

\begin{lemma}[Uniqueness from the principal part]\label{lem:haagerup-character-unique}
Two character vectors for this modular representation with the same
$\gamma$ coincide in every weight.
\end{lemma}
\begin{proof}
The modular representation has finite image by the congruence theorem
\cite[Theorem~I]{DLN}. A constant vector must be fixed by both $S$ and
the \emph{character} $T$ matrix, which
includes $e^{-2\pi i/3}$. The only $T$-fixed coordinate is the object
of categorical twist $e^{2\pi i/3}$. Its $S$ image has a nonzero vacuum
component, so it cannot span an invariant line. There are therefore
no invariant constants, and Lemma~\ref{lem:pole-cancellation} applies.
\end{proof}

Evans--Gannon \cite[Sections~1.4 and~2.4, equation~(2.7)]{EG}
provide a fundamental matrix $\Xi$ whose
vacuum and $d_3$ columns have precisely the two principal parts above.
Thus, in their labeling,
\[
 \boldsymbol\chi=\Xi_{:,0}+\gamma\Xi_{:,d_3},
 \qquad\dim U_1=6+13\gamma.
\]
The actual six-current result forces $\gamma=0$. Their recursive
construction of $\Xi$ and the uniqueness lemma give the entire vector,
not just a finite coefficient match. In particular,
\begin{equation}
 \operatorname{ch}U
 =q^{-1/3}\bigl(1+6q+120q^2+956q^3+6010q^4+\cdots\bigr).
 \label{eq:vacuum-character}
\end{equation}
The lowest weights and top dimensions are
\[
\begin{array}{c|c|r@{\qquad}c|c|r}
 \text{label}&h&\dim M_{\rm top}&\text{label}&h&\dim M_{\rm top}\\\hline
 0&0&1&d_1&6/13&7\\
 b&1&80&d_2&11/13&42\\
 a&1&81&d_3&15/13&119\\
 c_0&1&81&d_4&5/13&5\\
 c_1&1/3&3&d_5&7/13&13\\
 c_2&2/3&27&d_6&8/13&14
\end{array}
\]
In particular the final maximum is $15/13$, improving the preliminary
$17/10$ bound. The proof record is
\artifact{hi_cft_20260923_next_structure_conditional_characters.md}.
Although this vector fixes every unflavoured homogeneous dimension,
it does not determine the separate branching functions for the
fourteen-module core $A_L$, and it was not used to prove properness.

\begin{corollary}[The four unmixed lattice restrictions]
\label{cor:unmixed-lattice-restrictions}
The two lattice modules labelled $(0,\pm1)\in(\mathbb Z/3)^2$
restrict to the simple $U$-module $c_1$, and the two labelled
$(\pm1,0)$ restrict to $c_2$.
\end{corollary}
\begin{proof}
The first pair has lowest weight $1/3$, top dimension three, and
weights in $1/3+\mathbb Z_{\ge0}$. The next possible summand top,
$4/3$, exceeds the established maximum $15/13$. The character list
contains a unique simple of lowest weight $1/3$, with top dimension
three, so semisimplicity gives precisely one copy of $c_1$.
The second pair has lowest weight $2/3$, top dimension $27$, and
next possible top $5/3$; the same argument gives $c_2$.
No additional finite-top irreducibility calculation is required.
\end{proof}

% Fragment: analytic comparison, unitary identification, and full theories.
% No preamble.  The master document supplies theorem-like environments,
% criticalnote/reviewnote, \artifact, \Rep, \Com, \FPdim, \Vec, and \HH.

\section{From vacuum unitarity to complete unitarity}
\label{sec:analytic-unitarity}

The exact strong-generation certificates concern the actual subVOA
$U\subset B\subset H=V_{E_8}$.  In particular, they establish
$\dim U_1=6$ and $U\ne B$; they do not establish positivity on all
modules or on fusion products.  We now separate those questions.
\emph{Vacuum unitarity} is the positive invariant form inherited by the
PCT-stable subVOA $U$.  \emph{Strong unitarity} means that all simple
ordinary $U$-modules admit compatible positive invariant forms.
\emph{Complete unitarity} additionally includes positivity of the
canonical Hermitian fusion forms, giving the VOA representation category
its unitary modular tensor structure.  An abstract ribbon equivalence to
a unitary category does not, by itself, prove this last positivity statement.

We use the preceding algebraic analysis as input: regularity of $U$,
the actual restrictions
\begin{equation}
 B\mathord\downarrow_U=U\oplus X,
 \qquad
 H\mathord\downarrow_U=U\oplus X\oplus Y^{\oplus2},
 \qquad h_X=h_Y=1,
 \label{eq:analytic-branching}
\end{equation}
the twelve simple algebraic $U$-modules, and the
$\operatorname{HI}(\mathbb Z/3)$ fusion ring of the all-module category
for the original extension.  Every use of these inputs inherits the
qualifications attached to their preceding proofs.  Agreement of modular
data alone is never used as a classification of tensor categories.

\begin{criticalnote}{Analytic proof boundary}
We distinguish two comparisons.  The linear reconstruction of modules
and the independent sector count are proved below from the stated
algebraic inputs and external reconstruction theorems.  These resolve
the earlier obligations A3 and A2, respectively.  The remaining
obligation A1 identifies the particular bounded annulus insertions used
in the positivity theorem.  It is stronger than a linear equivalence
of module categories.  Complete unitarity, the unitary ribbon
identification for the VOA, and the specifically $\HH_3$ VOA extension
retain this explicit hypothesis.  The net category itself is identified
unitarily without A1.  The arguments below are mathematical deductions
for review, not an assertion of independent expert validation.
\end{criticalnote}

\subsection{The subnet and its sector count}
\label{subsec:analytic-net-count}

There are two net constructions in use.  Write $\mathcal N_V$ for the
net generated by the closed smeared fields of a strongly local VOA $V$,
and temporarily write $\mathcal A_V^{\mathrm{bl}}$ for the net generated
by Tener's bounded localized insertions.  For $V=U,B,H$ both exist:
strong locality passes to unitary subalgebras, as do bounded localized
vertex operators.  The conformal vectors are the ambient ones, so the
rotations and central charge agree in each construction.

First work with the bounded-insertion inclusion.  Regularity of $U$
and Tener~\cite[Corollary~4.24]{TenFusion} give
\[
 [\mathcal A_H^{\mathrm{bl}}:\mathcal A_U^{\mathrm{bl}}]<\infty.
\]
The ambient net is the holomorphic lattice net, by
Tener~\cite[Proposition~7.5]{TenFields}.  Hence
Longo~\cite[Theorem~24]{Longo} makes $\mathcal A_U^{\mathrm{bl}}$ completely
rational.  Tener's Lemma~7.6 gives a vacuum net isomorphism
$\mathcal A_U^{\mathrm{bl}}\cong\mathcal N_U$, using the ambient
lattice comparison.  It follows that $\mathcal N_U$ is completely
rational as well.  Only this isomorphism-invariant property is
transferred here: we do not infer equality of the two concrete net
constructions or an identification of their embedded extensions.
In particular, this step does not assume strong unitarity of $U$.

The actual smeared-field inclusion $\mathcal N_U\subset\mathcal N_H$
will be used for the sector calculation.  Its branching and index are
established below, after the elementary domain lemma needed for
reconstruction.

\begin{proposition}[The actual net branching and sector count; A2]
\label{clm:analytic-net-count}
Assume the algebraic branching and rank-six ring deductions of
Section~\ref{sec:branching}, including their certificate hypotheses.
The restriction of the holomorphic $H$ net has precisely the three
inequivalent irreducible $\mathcal N_U$ summands corresponding to
\eqref{eq:analytic-branching}, with multiplicities $1,1,2$.
Its boundary fusion category $\mathcal F_{\mathrm{net}}$ has rank six,
untwisted invertible subcategory $\Vec_{\mathbb Z/3}$, and the
$\operatorname{HI}(\mathbb Z/3)$ fusion ring.  Consequently
\[
 \Rep(\mathcal N_U)\simeq Z(\mathcal F_{\mathrm{net}}),
 \qquad
 \#\operatorname{Irr}\Rep(\mathcal N_U)=12.
\]
\end{proposition}
\begin{proof}
Every field $v\in U$ satisfies polynomial energy bounds on $H$, since
it is a field of the strongly local lattice VOA and the conformal
Hamiltonian is unchanged.  Its restriction to an algebraic unitary
summand satisfies the same bound.  Consequently the finite-energy
spaces are cores for the smeared fields in the respective Hilbert
completions.  Apply Lemma~\ref{lem:analytic-graph-reduction} below to
the orthogonal summand projections.  This proves that the algebraic
decomposition of $H$ is a Hilbert direct-sum decomposition for the
actual represented subnet.  Conversely, a subnet intertwiner commutes
with its local diffeomorphisms and hence with $L_0$.  The equality of
closed fields in reconstruction, or Fourier coefficients of their
matrix elements, then makes its finite-energy restriction a
$U$-module map.  Thus simple summands stay irreducible and inequivalent,
and the multiplicities remain $1,1,2$.

The inclusion is irreducible.  To see the local assertion explicitly,
let $x\in\mathcal N_H(I)\cap\mathcal N_U(I)'$.  Locality makes $x$
commute with $\mathcal N_U(I')$ as well.  Strong additivity of
$\mathcal N_U$ shows that the algebras of $I$ and its complementary
interval generate the global represented subnet: remove the two
endpoints from any test interval and apply strong additivity twice.
Thus $x$ commutes with the global subnet and, in particular, its local
diffeomorphisms and common rotations.  Its action on the vacuum lies
in the one-dimensional zero-energy space of $H$, so the separating
property of the vacuum makes it scalar.  The precise finite-index
input is now
Kawahigashi--Longo~\cite[Proposition~2.3]{KL}: an irreducible local extension
of a net with finitely many irreducible DHR sectors, and in which every
sector is a direct sum of finite-statistical-dimension sectors, has
finite index.  These hypotheses follow
from the complete rationality already established for $\mathcal N_U$.
Thus the actual inclusion has finite index.  Its dual canonical sector
is the restriction of the $H$ vacuum computed above, and its dimension
is the inclusion index.  Its algebra $A_{\mathrm{net}}$ is Lagrangian because
$\mathcal N_H$ is holomorphic.  The preceding comparison also gives
the algebra identity
\begin{equation}
 \operatorname{End}(A_{\mathrm{net}})
 \cong\mathbb C\oplus\mathbb C\oplus M_2(\mathbb C).
 \label{eq:analytic-end-algebra}
\end{equation}
Here the multiplication is composition, rather than a dimension count
alone.  Center induction identifies its dimension with
$\operatorname{rank}(\mathcal F_{\mathrm{net}})$, and identifies the
algebra with the complex fusion algebra of
$\mathcal F_{\mathrm{net}}$.  Thus the latter has rank six and is
noncommutative.

We next identify the invertible objects using only the actual
$\mathbb Z/3$ lattice-extension automorphisms of $H/B$.  They act
on the net inclusion and hence on its canonical algebra.  The
automorphism--invertible-object injection of
Lemma~\ref{lem:automorphisms-invertibles} embeds this group in
$\operatorname{Inv}(\mathcal F_{\mathrm{net}})$.  Since the category
has rank six, the latter group has order three or six.  If its order
were six, the category would be pointed of global dimension six.
The intermediate lattice inclusion has $[\mathcal N_H:\mathcal N_B]=3$,
so multiplicativity of index would give
$[\mathcal N_B:\mathcal N_U]=2$.  An irreducible local index-two
extension is the fixed-point extension of its sign involution: its
canonical algebra is $\mathbf1\oplus\beta$, where $\beta^2=\mathbf1$,
and $\mathbf1\oplus(-\mathrm{id}_\beta)$ is its nontrivial algebra
automorphism.  Theorem~6.9 of \cite{CKLW} lifts this vacuum-preserving
net automorphism to a unitary automorphism $g$ of the actual $B$.
Their Proposition~7.6 gives
$\mathcal N_{B^g}=\mathcal N_B^g=\mathcal N_U$, and injectivity of
their unitary-subalgebra/subnet correspondence (Theorem~7.5) gives
$B^g=U$.  This contradicts
Lemma~\ref{lem:no-sign-extension}.  Consequently
\[
 \operatorname{Inv}(\mathcal F_{\mathrm{net}})\cong\mathbb Z/3.
\]
This argument needs no assertion about the full pointwise stabilizer
$\operatorname{Aut}(B/U)$.

The lattice net $\mathcal N_B$ has discriminant metric group
$\mathbb Z/3\times\mathbb Z/3$ with
$q(a,b)=(2a^2+b^2)/3\pmod{\mathbb Z}$.  The extension $H/B$ is its
diagonal isotropic line $L_+=\{(a,a)\}$, and
$L_-=\{(a,-a)\}$ is a transverse isotropic line.  The subcategory
$\mathcal T_-$ supported on $L_-$ is Tannakian and is tensor
equivalent to the untwisted $\Vec_{\mathbb Z/3}$.  Let $Q_+$ be the
extension algebra supported on $L_+$.  Free-module induction
$t\mapsto t\otimes Q_+$ is a tensor functor from $\mathcal T_-$ to
the category of all $Q_+$-modules.  For simple $t,t'\in\mathcal T_-$,
its adjunction gives
\[
 \operatorname{Hom}_{Q_+}(t\otimes Q_+,t'\otimes Q_+)
 \cong\operatorname{Hom}(t,t'\otimes Q_+)
 =\begin{cases}\mathbb C,&t=t',\\0,&t\ne t'.\end{cases}
\]
Here transversality $L_+\cap L_-=\{0\}$ gives the last equality.
The target has dimension $9/3=3$, so these three distinct invertible
images exhaust its simple objects.  Induction is therefore a tensor
equivalence.  Forgetting the intermediate $B$-action embeds this
category fully faithfully and tensorially in
$\mathcal F_{\mathrm{net}}$; relative tensor product over the same
$H$ algebra is unchanged.  Thus the entire pointed subcategory is
untwisted.

We may now apply the rank-six ring and indicator argument of
Section~\ref{sec:branching} to $\mathcal F_{\mathrm{net}}$ itself.
Its hypotheses are the fusion-algebra decomposition above, the three
invertibles, and their untwisted associator.  They have just been
proved for this category, and its spherical structure is positive
because it is a $C^*$ fusion category.  The argument gives
\eqref{eq:rank6-ring} with multiplicity one.  In particular, for
$d=(3+\sqrt{13})/2$ and $s=1+d^2$, the center calculation gives dimensions
\begin{equation}
 1,\ d^2,\ \underbrace{s,s,s,s}_{4},\
 \underbrace{3d,3d,3d,3d,3d,3d}_{6}.
 \label{eq:analytic-center-dimensions}
\end{equation}
The same center-induction and indicator calculation gives the last
six twists with exponents
$\{2,5,6,7,8,11\}/13$.  This is the source of the net count
twelve.  Finally the Lagrangian-algebra theorem gives
$\Rep(\mathcal N_U)\simeq Z(\mathcal F_{\mathrm{net}})$, and the
displayed center dimensions give its rank twelve.  No VOA character
transformation matrix was assigned to the net category.
\end{proof}

For comparison, the first paragraph also follows for the
bounded-insertion inclusion directly from
Tener~\cite[Proposition~5.5 and Theorem~5.6]{TenFields}.  The argument above has the
additional advantage that the later automorphism comparison refers
throughout to the original lattice fields.  The purely categorical
deductions used in the proof are recorded in
\artifact{hi_cft_20260923_next_structure_rank46.md},
\artifact{hi_cft_20260923_next_structure_rank6_indicators.md}, and
\artifact{hi_cft_20260923_next_structure_rank6_center.md}.

\subsection{Reverse reconstruction and exhaustion}
\label{subsec:analytic-reconstruction}

Henriques--Tener~\cite[Theorem~9.5 and Corollary~9.7]{HTNet} reconstruct unitary
modules and identify their closed smeared fields.  Remark~9.8 extends
the construction beyond the finite-multiplicity hypothesis.  This
extension matters: finite-dimensional energy eigenspaces must not be
silently assumed for the net representation one is trying to reconstruct.
Their Theorem~8.9 identifies the reconstructed vacuum VOA of
$\mathcal N_U$ with the original $U$, since $U$ is strongly local.

The closure step can be isolated from the VOA theory.  Its formulation
also explains why a formal mode identity alone is insufficient.

\begin{lemma}[Reduction on a common core]
\label{lem:analytic-graph-reduction}
Let $T_0:D\to\mathcal K$ be a densely defined closable operator and
$P$ an orthogonal projection.  Suppose $PD\subset D$ and
$T_0P\xi=PT_0\xi$ for $\xi\in D$.  Then $P$ reduces
$T=\overline{T_0}$ and $T^*$, and it commutes with the von Neumann
algebra generated by $T$.  In particular, this holds simultaneously
for a family of fields satisfying these hypotheses on a common core.
\end{lemma}
\begin{proof}
If $\xi_j\in D$ converges to $\xi\in D(T)$ in graph norm, then
$P\xi_j\to P\xi$ and $T_0P\xi_j=PT_0\xi_j\to PT\xi$.
Thus $P\xi\in D(T)$ and $TP\xi=PT\xi$.  The same assertion holds
for $1-P$.  It follows that $T$ is the direct sum of its closed
restrictions to $P\mathcal K$ and $(1-P)\mathcal K$; taking adjoints
gives the corresponding direct sum for $T^*$.  Its polar partial
isometry and the spectral projections of $|T|$ are therefore block
diagonal.  These generate the von Neumann algebra of $T$, proving the
last assertion.  Notice that no separate assertion that $D$ is a core
for $T^*$ is required.
\end{proof}

\begin{proposition}[Reverse module comparison; A3]
\label{clm:analytic-reconstruction}
For this regular $U$, reconstruction sends each irreducible
$\mathcal N_U$ sector to an irreducible ordinary unitary $U$-module.
It preserves $L_0$ and induces a fully faithful linear $*$-functor on
finite direct sums,
\[
 \mathcal R:\Rep(\mathcal N_U)\longrightarrow\Rep^{u}(U).
\]
The functor preserves conformal Hamiltonians.  No tensorator or
compatibility with braiding is asserted.
\end{proposition}
\begin{proof}
Let $\mathcal K$ be an irreducible net representation and put
\[
 M=\bigcup_{R>0}1_{[0,R]}(L_0)\mathcal K
\]
as in Henriques--Tener, Remark~9.8.  This is a dense admissible weak
$U$-module with a positive invariant Hermitian form.  Theorem~9.5(ii)
and its extension in that remark identify the closed reconstructed
smeared field with $\pi_I(Y_U(v,f))$ for
$\operatorname{supp}(f)\subset I$.  By Definition~9.4, $M$ is a
core for this closed operator.  The proof of Corollary~9.6 applies
unchanged on bounded-energy vectors, so the VOA zero mode of the
conformal vector is precisely the self-adjoint rotation generator on
this domain.

Regularity decomposes the weak module $M$ algebraically into ordinary
simple modules.  Choose a nonzero simple summand $S$.  Its inherited
form is positive and invariant.  Since its $L_0$ is the restriction
of the net Hamiltonian, its energy decomposition is
\[
 S=\bigoplus_{n\geq0}S_{h+n},\qquad
 \dim S_{h+n}<\infty,\quad h\geq0.
\]
The Hilbert closure $\overline S$ is the orthogonal Hilbert sum of
these finite-dimensional spaces.  In particular it reduces $L_0$,
and the projection $P$ onto $\overline S$ satisfies
\begin{equation}
 P1_{[0,R]}(L_0)\mathcal K
   \subset\bigoplus_{h+n\leq R}S_{h+n}\subset S,
 \qquad PM\subset S\subset M.
 \label{eq:analytic-projection-core}
\end{equation}
This establishes the domain invariance before any commutation with
fields is used.  It requires finite energy multiplicities only for the
chosen ordinary summand, not for $\mathcal K$.

For a homogeneous $v\in U$, the unitary invariance identity expresses
the adjoint of a mode $v_{(n)}$, on $M$, as a finite linear combination
of modes of the vectors $L_1^j\theta v$, where $\theta$ is PCT.
Every such mode preserves $S$.  Testing against $s\in S$ therefore
gives
\[
 \langle Pv_{(n)}\xi,s\rangle
 =\langle\xi,v_{(n)}^{\dagger}s\rangle
 =\langle P\xi,v_{(n)}^{\dagger}s\rangle
 =\langle v_{(n)}P\xi,s\rangle.
\]
Both outer vectors lie in $S$, by
\eqref{eq:analytic-projection-core}, whence
$Pv_{(n)}\xi=v_{(n)}P\xi$ for $\xi\in M$.
The smeared fields on $M$ are Hilbert-space vectors, as in the paragraph
preceding Henriques--Tener Definition~9.4.  Taking their energy
components, or equivalently their convergent Fourier sums, now yields
\[
 P Y^M(v,f)\xi=Y^M(v,f)P\xi,\qquad \xi\in M.
\]
There is no assertion here that a smeared field maps $M$ into itself.
Only its value in $\mathcal K$ and the boundedness of $P$ are needed.
Lemma~\ref{lem:analytic-graph-reduction} extends this identity to the
closures and their adjoints.

The local algebras are generated by the closed smeared fields.
Normal representations preserve this generation statement (equivalently,
apply them to the polar decompositions and spectral projections).
Thus $P$ is a net intertwiner.  Irreducibility and $S\ne0$ imply
$P=1$.  Equation~\eqref{eq:analytic-projection-core} then implies
$M=S$.  This proves ordinaryness and simplicity, including the removal
of the infinite multiplicities permitted in Remark~9.8.

Finally consider morphisms.  A net intertwiner commutes with the
canonical diffeomorphism action, because that action is generated by
the local diffeomorphism groups.  It therefore preserves bounded-energy
vectors.  Commutation with the reconstructed fields gives commutation
with their Fourier coefficients, hence a module map.  Conversely, an
isomorphism of simple ordinary unitary modules is a scalar multiple
of a unitary.  To verify this without presupposing boundedness, compare
the two invariant forms on each finite-dimensional energy space.  The
resulting grade-preserving positive operator is an algebraic
$U$-endomorphism, by invariance of the forms.  An eigenvector in the
finite-dimensional lowest space has an eigenspace that is a nonzero
$U$-submodule.  Simplicity makes that eigenspace the whole module, so
the comparison operator is a single positive scalar on all grades.
The isomorphism therefore has a bounded unitary rescaling.  Its
unitary extension intertwines the smeared fields on their common cores
and hence their closures, so it intertwines the nets.  Between
inequivalent simples both morphism spaces vanish.  Matrices of these
maps handle finite direct sums and respect adjoints.  This proves full
faithfulness and the asserted linear $*$-functor.
\end{proof}

The preliminary version of this deduction is recorded in
\artifact{hi_cft_20260923_next_m3_unitarity_scope.md}.

\begin{proposition}[Strong unitarity from sector exhaustion]
\label{prop:analytic-strong-unitarity}
Under the preceding algebraic inputs, including the count of twelve
simple $U$-modules, every simple $U$-module is
unitarizable, and its character is the character of a net sector.
\end{proposition}
\begin{proof}
Full faithfulness makes the twelve net simples pairwise inequivalent
simple ordinary $U$-modules.  There are exactly twelve algebraic simple
modules, so this injection is a bijection.  Each image carries its
reconstructed positive form, and $L_0$ is preserved.  Finite direct sums
give the corresponding statement for ordinary modules.
\end{proof}

This counting step supplies essential surjectivity only as a linear
comparison.  Neither fusion compatibility, equality of categorical
$S$ matrices, nor positivity of VOA fusion forms follows from it.

\subsection{The annulus compatibility claim}
\label{subsec:analytic-annuli}

The next step concerns the specific geometric annulus system of
Tener~\cite[Definition~3.18 and Corollaries~3.20--3.22]{TenFields}.
Let $(B,A)$ be a localized pair for an interval $I$, with incoming
factor $A$ and outgoing factor $B$.  Subscripts distinguish their
actions on the vacuum and on a reconstructed module.  For an insertion
point $z\in\operatorname{int}(B,A)$, put
\[
 X_v(B,z,A)=\overline{B_UY_U(v,z)A_U}
       \in\mathcal A_U^{\mathrm{bl}}(I).
\]
Here the bar denotes the bounded operator represented by the
sesquilinear insertion form.  At a point outside the incoming annulus
one must not silently treat the unsmeared formal field as an everywhere
defined Hilbert-space operator.  The earlier abstract isomorphism
$\mathcal A_U^{\mathrm{bl}}\cong\mathcal N_U$ does not identify these
operators.  The following claim includes the specific vacuum
identification needed to compare them, and then its module version.

\begin{criticalnote}{A1: new annulus compatibility lemma}
\begin{claim}[Annulus compatibility; unproved geometric comparison]
\label{clm:analytic-annulus}
For the actual geometric system above, choose matching paths and
central lifts in the two annulus constructions.  In the vacuum, the
geometric insertion of Henriques--Tener is $X_v(B,z,A)$ and identifies
$\mathcal A_U^{\mathrm{bl}}$ with $\mathcal N_U$ on the original
vacuum Hilbert space.  If $M$ is an ordinary
unitary module reconstructed from a sector $\pi_M$, then
\begin{equation}
 \pi_M\bigl(X_v(B,z,A)\bigr)
   =\overline{B_MY_M(v,z)A_M},
 \qquad v\in U,\quad z\in\operatorname{int}(B,A).
 \label{eq:analytic-annulus-identity}
\end{equation}
The right side exists as a bounded extension of the insertion form.
\end{claim}
The remaining point requiring verification is a comparison of the
specific geometric annulus systems and their central lifts, including
their localizing diffeomorphisms.  The domain-density and continuation
steps are proved in the lemmas below.  They do not require the
inverse-finite-energy domains to be graph cores.  Agreement after an
unspecified vacuum net isomorphism would not prove the displayed
operator identity.
\end{criticalnote}

The unresolved comparison has a concrete formulation.  For a geometric
incoming annulus $X=(\psi,\rho,t,\gamma)$ in Tener's Definition~3.16,
his vacuum operator is
\[
 U(\psi)e^{-tL(\rho)}U(\gamma)^*.
\]
One must select the corresponding centrally extended annulus
$\underline X$, with the same embedded inner boundary and boundary
parametrizations, and identify this expression with its integrated
action.  The time-ordered integration in
Henriques--Tener~\cite[Theorem~5.5 and formula~(5.2)]{HTAnn} gives the candidate:
constant generators integrate to exponentials, and thin factors to the
diffeomorphism action.  The comparison must retain the central scalar
and use the same path lifts.  A lift normalized only by its action on
the vacuum need not have the same action in a sector of fractional
lowest weight.

More precisely, choose the common localizing diffeomorphism
$\widehat\gamma$ of Tener's Corollary~3.20.  The required geometric
statement is that $\widehat\gamma\underline X$ is the corresponding
localized centrally extended annulus, with its specified lift, and
that its vacuum action is the local Virasoro operator
$U(\widehat\gamma)\pi(X)$.  Once this statement is known,
Henriques--Tener Proposition~4.17 transports it to every sector:
both constructions act there by
\[
 U^M(\widehat\gamma)^*
 \pi^M_I\bigl(U(\widehat\gamma)\pi(X)\bigr).
\]
The same localizing lift must be used for the incoming and outgoing
members of a pair.  This is the precise normalization problem retained
in A1; scalar equality in just one representation is insufficient.

\begin{lemma}[Dense inverse-energy domains]
\label{lem:analytic-inverse-domain}
Let $M$ be an ordinary unitary module and let $A_M,C_M$ be the actions
of geometric annuli on its Hilbert completion.  Suppose
\begin{equation}
 A_MC_M=\lambda r^{L_0},\qquad 0<r<1,\quad\lambda\ne0.
 \label{eq:analytic-round-factorization}
\end{equation}
Then $A_M^{-1}M^{\mathrm{fin}}=C_M M^{\mathrm{fin}}$ is dense.
The same conclusion holds for the outgoing adjoint factor.
\end{lemma}
\begin{proof}
The geometric annulus actions are injective and have dense range by
Henriques--Tener Lemma~4.15.  Equation
\eqref{eq:analytic-round-factorization} first gives
$C_M M^{\mathrm{fin}}\subset A_M^{-1}M^{\mathrm{fin}}$.
Conversely, if $A_M\xi\in M^{\mathrm{fin}}$, then
\[
 A_M\xi
  =A_MC_M\lambda^{-1}r^{-L_0}A_M\xi,
 \qquad
 \xi=C_M\lambda^{-1}r^{-L_0}A_M\xi.
\]
The last equality uses injectivity of $A_M$, and proves the reverse
inclusion.  Finally $M^{\mathrm{fin}}$ is dense and $C_M$ is bounded
with dense range; hence $C_M M^{\mathrm{fin}}$ is dense.  This is the
argument of Henriques--Tener Lemma~7.9 with a central scalar retained.
\end{proof}

For the geometric annuli under consideration the factorization itself
is elementary.  The incoming annulus is embedded in the unit disc,
its smooth inner Jordan domain contains zero, and its outer
parametrization is the identity.  Choose a sufficiently small round
disc $r\mathbb D$ inside that inner domain.  The intervening annulus
$C$ has its outer parametrization chosen to sew to the inner boundary
of $A$, and its inner parametrization is the standard round one.
Then $AC$ is the round annulus.  Choosing compatible lifts gives
\eqref{eq:analytic-round-factorization}.  This proves density after
the annulus-action identification specified above; density is not an
additional unknown analytic assumption.

Define
\begin{equation}
 \mathcal D_A=\{\xi\in\mathcal K_M:A_M\xi\in M^{\mathrm{fin}}\},
 \qquad
 \mathcal D_{B^*}=\{\eta\in\mathcal K_M:B_M^*\eta\in M^{\mathrm{fin}}\}.
\label{eq:analytic-annulus-domains}
\end{equation}
Both domains are dense by Lemma~\ref{lem:analytic-inverse-domain}.

\begin{lemma}[Continuation of the insertion form]
\label{lem:analytic-insertion-continuation}
Let $O_{\mathrm{in}}$ and $O_{\mathrm{out}}$ be the incoming and
outgoing insertion interiors, and put
\[
 \Omega=\operatorname{int}
  (\overline{O_{\mathrm{in}}}\cup\overline{O_{\mathrm{out}}}).
\]
Suppose a bounded geometric insertion $G_v(z)$ is holomorphic on
$\Omega$.  If its matrix elements equal the algebraic insertion form
on $O_{\mathrm{in}}\cup O_{\mathrm{out}}$, for vectors in the two
domains \eqref{eq:analytic-annulus-domains}, the equality holds
throughout $\Omega$.  It uniquely defines a bounded extension of that
form.
\end{lemma}
\begin{proof}
For $\xi\in\mathcal D_A$ and $\eta\in\mathcal D_{B^*}$, consider
\begin{equation}
 \left\langle Y_M(v,z)A_M\xi,B_M^*\eta\right\rangle
 \label{eq:analytic-laurent-pairing}
\end{equation}
as a formal matrix coefficient.  It is a finite Laurent polynomial:
if $v$ has weight $k$, and the input and output energies are $p,q$,
only the mode with $q=p+k-n-1$ can contribute.  Both vectors have
finite energy support.  Thus the form and the geometric matrix
coefficient are holomorphic on $\Omega\subset\mathbb C^\times$.

Every component of $\Omega$ meets
$O_{\mathrm{in}}\cup O_{\mathrm{out}}$.  Indeed, its complement
inside $\Omega$ is contained in the union of the two boundaries;
boundaries of open sets are closed with empty interior, and a finite
union of such sets contains no nonempty open component.  The identity
theorem therefore extends the equality across every component,
including any part of the sewing circle lying in $\Omega$.
For fixed $z$, the form is now the restriction of the bounded operator
$G_v(z)$ to two dense test domains.  This proves existence and
uniqueness of its bounded extension.  A graph-core assertion is neither
used nor needed.
\end{proof}

After the geometric normalization in A1, Henriques--Tener
Theorem~9.5(i) supplies the incoming equality.  Sewing is the
multiplicative action constructed in their Section~6.2; the
$*$-compatibility of Proposition~6.19 supplies the outgoing equality.
Lemma~\ref{lem:analytic-insertion-continuation} then proves the full
identity \eqref{eq:analytic-annulus-identity}.  Thus the operator-domain
and continuation parts of the earlier proposed proof have been
resolved.  The remaining issue is the specified comparison of the
two geometric operator constructions, not an implication from strong
unitarity alone.

\begin{theorem}[Conditional complete unitarity]
\label{thm:analytic-complete-unitarity}
Assume Claim~\ref{clm:analytic-annulus}, together with the preceding
algebraic inputs.  Then $U$ is completely
unitary.
\end{theorem}
\begin{proof}
Proposition~\ref{prop:analytic-strong-unitarity} unitarizes every simple
module.  Claim~\ref{clm:analytic-annulus} supplies the representation
required by
Tener~\cite[Definition~2.18]{TenFusion}, for each such module; finite direct sums
preserve the property.  His Corollary~4.6 then supplies positivity of
the canonical fusion forms.  Strong unitarity must be established first:
the unitarizable-fusion-product hypothesis in Theorem~4.4 cannot be used
to bypass the exhaustion argument.
\end{proof}

Recent comparison results leave this specific hypothesis in place.
Gui's extension comparison assumes Condition~II, including complete
unitarity and energy bounds and strong intertwining for a generating
family of charged fields; his inheritance theorem concerns regular
cosets \cite[Definition~4.15 and Theorems~4.20,~6.11]{GuiCompare}.
It does not give inheritance by the present conformal subalgebra:
the commutant of $U$ in $H$ is the vacuum VOA and its double
commutant is $H$.  Henriques--Tener construct bounded vacuum
multidisc operators and prove the vacuum trace-class condition
\cite[Theorems~3.6 and~4.1]{HTGenusZero}; these statements do not
identify the module insertion operators in
\eqref{eq:analytic-annulus-identity}.  A sufficient remaining target
is to match Tener's incoming and outgoing geometric annuli with the
integrated annuli, retaining their embedded boundaries, boundary
parametrizations, common localizing diffeomorphism, and central path
lifts, and then to identify their vacuum point insertions.  Sector
transport, the dense inverse-energy domains, and continuation proved
above would then give A1.  A linear equivalence of unitary module
categories alone supplies neither this comparison nor the positivity
of the canonical VOA fusion forms.

\section{Unitary categorical identifications}
\label{sec:analytic-category}

\begin{lemma}[All modules over a commutative Q-system]
\label{lem:analytic-unitary-all-modules}
Let $\mathcal C$ be a unitary braided fusion category and $A$ a
connected commutative Q-system in $\mathcal C$.  The category
$\mathcal C_A$ of all right $A$-modules, with relative tensor product,
is a unitarizable fusion category.  This assertion does not require
the modules to be local.
\end{lemma}
\begin{proof}
If $r:M\otimes A\to M$ is a right action, the braiding gives a left
action $l=r\circ c_{A,M}$.  Commutativity of $A$ and the hexagon
identity make these commuting actions, and hence an $A$-bimodule.
Naturality of the braiding makes every right-module morphism a
bimodule morphism for these choices.  This realizes $\mathcal C_A$
as the full subcategory of bimodules with $l=r\circ c_{A,M}$.
The relative tensor product is obtained by splitting the separability
idempotent on $M\otimes N$.  Its induced left and right actions again
satisfy this relation, by the same hexagon identity and the balancing
relation.  Thus the subcategory is closed under relative tensor
products, sums and summands, and its unit is $A$.

Every such bimodule is unitarizable by
Gui~\cite[Theorem~2.28]{Gui}; an isomorphism transporting the actions preserves
the displayed relation.  The unitary relative tensor product and its
associativity and unit constraints are supplied by Gui, Theorem~3.13.
The above full subcategory therefore inherits a $C^*$ tensor structure.
Separability gives semisimplicity, connectedness makes its unit simple,
and the usual module-category construction over a connected separable
algebra gives rigidity and finiteness.  It is consequently a unitary
fusion category up to tensor equivalence.

For clarity, locality is the additional monodromy condition
$r\circ c_{A,M}\circ c_{M,A}=r$.  We have not imposed it.  Only the
local subcategory inherits the ambient braiding by descent.
\end{proof}

\begin{theorem}[Unitary Haagerup-center identification, conditional on A1]
\label{thm:analytic-category}
Under the hypotheses of Theorem~\ref{thm:analytic-complete-unitarity},
and the preceding determination of the all-module fusion ring and its
untwisted pointed subcategory, there is a unitary ribbon equivalence
\[
 \Rep(U)\simeq Z(\HH_3).
\]
For the original embedded extension $H/U$, its all-module category is
tensor equivalent to one of $\HH_2$ and $\HH_3$.
\end{theorem}
\begin{proof}
Put $\mathcal C=\Rep(U)$ and let $A$ be the algebra of the actual
unitary extension $H$.  Complete unitarity allows application of
Gui~\cite[Theorems~2.21 and~2.28]{Gui}: $A$ is a standard commutative
Q-system and its modules are unitarizable.
Lemma~\ref{lem:analytic-unitary-all-modules} makes
$\mathcal F=\mathcal C_A$ a $C^*$ fusion category.  Here $\mathcal C_A$ means
\emph{all} modules.  The local modules form $\mathcal C_A^0\simeq\Vec$
because the extension is holomorphic, so substituting local modules
would lose the Haagerup category.

Its already determined fusion ring is
$\operatorname{HI}(\mathbb Z/3)$.  The actual untwisted pointed
subcategory makes the obstruction $c^{(0,3)}$ vanish, while the other
obstruction is in
\[
 H^1\!\left(\mathbb Z/2,H^2(\mathbb Z/3,\mathbb C^\times)\right)=0.
\]
Indeed, the cyclic-group calculation gives
$H^2(\mathbb Z/3,\mathbb C^\times)=
\mathbb C^\times/(\mathbb C^\times)^3=0$;
every nonzero complex number has a cube root.  Also
$(\mathbb Z/3)/2(\mathbb Z/3)=0$, so no choice of distinguished
noninvertible object introduces an additional classification orbit.
These are the hypotheses in
Izumi~\cite[Definition~2.7 and Theorems~5.9 and~9.1]{Izumi}, giving the two
unitary tensor classes $\HH_2,\HH_3$.  No algebra structure on
$1+\rho$ is inserted as an extra assumption.
Grossman--Snyder~\cite[Theorem~1.4]{GS} place these categories in the same
Morita class.  Their centers are therefore braided equivalent.
The positive spherical structures make this a ribbon equivalence, and
the Lagrangian-center theorem gives
\[
 \mathcal C\simeq Z(\mathcal F)\simeq Z(\HH_3).
\]
Reutter~\cite[Theorem~5]{Reutter} provides the unitary braided realization of the
equivalence, preserving the associated positive ribbon structures.
\end{proof}

This uses the published unitary classification as an external input;
its classification equations have not been replayed here.  It avoids
the broader pivotal census, but it does not prove that the natural
reconstruction functor $\mathcal R$ preserves Connes fusion, the VOA
tensor product, or braiding.  That particular comparison is a distinct
statement and is not needed for the theorem just proved conditionally.

\begin{corollary}[Haagerup-center identification for the conformal net]
\label{cor:analytic-net-haagerup}
Under the preceding algebraic inputs, the completely rational
central-charge-eight net $\mathcal N_U$ satisfies
\[
 \Rep(\mathcal N_U)\simeq Z(\HH_3)
\]
as a unitary ribbon category.  Its actual extension to the $E_8$
lattice net has boundary category equivalent to $\HH_2$ or $\HH_3$.
Neither conclusion requires Claim~\ref{clm:analytic-annulus}.
\end{corollary}
\begin{proof}
Proposition~\ref{clm:analytic-net-count} constructs
$\mathcal F_{\mathrm{net}}$ as a unitary fusion category with the
$\operatorname{HI}(\mathbb Z/3)$ fusion ring and an untwisted pointed
subcategory.  The Izumi classification and the vanishing obstruction
groups used in the preceding proof therefore apply directly to this
category.  Thus $\mathcal F_{\mathrm{net}}$ is tensor equivalent to
$\HH_2$ or $\HH_3$.  Their Morita equivalence, the unitary
Lagrangian-center theorem, and uniqueness of the unitary braided
structure give
$\Rep(\mathcal N_U)\simeq Z(\mathcal F_{\mathrm{net}})\simeq Z(\HH_3)$.
No positivity statement about VOA intertwining-operator spaces is
used in this argument.
\end{proof}

\section{Full net theory and a conditional \texorpdfstring{$E_8$}{E8} VOA extension}
\label{sec:analytic-full-theory}

\begin{corollary}[A chosen holomorphic net extension]
\label{cor:analytic-net-extension}
Under the preceding algebraic inputs, there is an irreducible local
conformal-net extension $\mathcal N_U\subset\mathcal E_3$ of central
charge eight.  It is holomorphic and has index
\[
 [\mathcal E_3:\mathcal N_U]
 =\operatorname{Dim}(\HH_3)=3(1+d^2),
 \qquad d=(3+\sqrt{13})/2.
\]
Its all-module boundary category is $\HH_3$, and its algebra-bimodule
category is $\HH_3\boxtimes\HH_3^{\mathrm{op}}$, with $36$ simple
objects.  This construction does not require A1.
\end{corollary}
\begin{proof}
Transport the canonical connected commutative Q-system
$I_{\HH_3}(\mathbf1)\in Z(\HH_3)$ along
Corollary~\ref{cor:analytic-net-haagerup}; call the result
$Q_3\in\Rep(\mathcal N_U)$.  The local-net reconstruction theorem
of Longo--Rehren~\cite[Theorem~4.9]{LRNets} produces the extension.
Connectedness is precisely the multiplicity-one hypothesis for the
identity sector, hence gives irreducibility.  The Q-system has finite
statistical dimension, which equals the displayed index; its
commutativity is precisely the locality condition in that theorem.
The conformal-net form of the correspondence is recalled in
\cite[Section~2.6]{BisTY}.  The represented local diffeomorphism
unitaries of the subnet give the conformal action of the extension,
so the central charge remains eight.

The finite-index theorem \cite[Theorem~24]{Longo} gives complete
rationality.  Since $Q_3$ is Lagrangian, the index relation gives
\[
 \mu(\mathcal E_3)
 =\frac{\mu(\mathcal N_U)}{\dim(Q_3)^2}=1,
\]
and the extension is holomorphic.  The canonical-algebra equivalences
of \cite[Corollary~4.1 and Proposition~4.4]{DMNO} identify the
all-module category with $\HH_3$.  The corresponding bimodule
category is $\HH_3\boxtimes\HH_3^{\mathrm{op}}$ by the
double-centralizer description \cite[Sections~4 and~8]{Mueger}; its
simple objects are the $6\cdot6$ ordered pairs.  As below, these
algebra-bimodule defects preserve the original $\mathcal N_U$ subnet.
\end{proof}

No identification of $\mathcal E_3$ with the $E_8$ lattice net is made
here.  Holomorphicity and central charge eight are the conclusions
of the cited net theorems; the VOA classification used later does not
by itself classify arbitrary holomorphic conformal nets.  Nor does
the present construction identify a reconstructed VOA extension with
the selected tensor category of VOA modules.

\begin{proposition}[Diagonal full theory]
\label{prop:analytic-diagonal}
Under the preceding algebraic inputs, the completely rational net
$\mathcal N_U$ admits its canonical unitary local diagonal full theory.
With the twelve-module character determination, its partition function is
\begin{equation}
 Z_{\mathrm{diag}}(\tau,\bar\tau)
   =\sum_{i=0}^{11}\chi_i(\tau)\overline{\chi_i(\tau)},
 \qquad c_L=c_R=8.
 \label{eq:analytic-diagonal-partition}
\end{equation}
\end{proposition}
\begin{proof}
Apply the diagonal Longo--Rehren construction, in the identity-coupling
case of
Rehren~\cite[Theorem~1.4, Proposition~1.5, and Corollary~1.6]{Rehren}.
Propositions~\ref{clm:analytic-net-count} and~\ref{clm:analytic-reconstruction}
identify all sector Hilbert spaces with the twelve ordinary
modules and preserve their Hamiltonians.  Hence their traces give
\eqref{eq:analytic-diagonal-partition}.  The annulus compatibility claim
is not needed for this net construction or this linear character matching.
\end{proof}

\begin{theorem}[Chosen $\HH_3$ extension and defects; conditional on A1]
\label{thm:analytic-e8-extension}
Assume the hypotheses of Theorem~\ref{thm:analytic-category} and the
preceding strict positivity of all nonvacuum lowest conformal weights.
There exists a unitary holomorphic conformal extension $H'\supset U$
isomorphic to $V_{E_8}$ whose all-module category is $\HH_3$.
For its extension algebra $A_3\in\Rep(U)$,
\begin{equation}
 \Rep(U)_{A_3}\simeq\HH_3,
 \qquad
 {}_{A_3}\Rep(U)_{A_3}\simeq\HH_3\boxtimes\HH_3^{\mathrm{op}}.
 \label{eq:analytic-extension-defects}
\end{equation}
The second category describes algebra-bimodule defects preserving the
$U$ chiral subalgebra in the corresponding full extension theory and
contains an $\HH_3$ fusion subcategory.  It has $36$ simple objects,
labelled by ordered pairs of simple objects of $\HH_3$.
\end{theorem}
\begin{proof}
Let $I_{\HH_3}$ be the right adjoint of the forgetful functor
$Z(\HH_3)\to\HH_3$.  Transport the canonical algebra
$I_{\HH_3}(\mathbf1)$ along the unitary ribbon equivalence.  By the
canonical Lagrangian-algebra construction, as in
Davydov--M\"uger--Nikshych--Ostrik~\cite[Corollary~4.1 and
Proposition~4.4]{DMNO}, the resulting $A_3$
is connected, commutative, separable, and twist-free, with
\[
 \FPdim(A_3)^2=\FPdim(\Rep(U)),
 \qquad \Rep(U)_{A_3}^{0}\simeq\Vec,
 \qquad \Rep(U)_{A_3}\simeq\HH_3.
\]
Trivial twist makes its constituent conformal weights integral.
Connectedness and strict positivity away from the vacuum give a
one-dimensional weight-zero space.  The extension correspondence of
Huang--Kirillov--Lepowsky~\cite[Theorems~3.2--3.6]{HKL}, together with Gui's
unitary Q-system theorem, produces a unitary conformal VOA extension
$H'$.  Here the transported algebra is a Q-system: the canonical
algebra in the unitary center has its standard positive Frobenius
structure, and the equivalence is unitary.  Thus Gui's theorem is
applied to a specified positive algebra structure, rather than
to an arbitrary separable algebra without a chosen $*$-structure.

For completeness, the finiteness condition on the extension can be
checked at the level of $C_2$ spaces.  Write
$H'=\bigoplus_i M_i^{\oplus m_i}$ as a finite sum of ordinary
$U$-modules and set
\[
 C_2^U(H')=\operatorname{span}\{u_{(-2)}w:u\in U,\ w\in H'\}.
\]
The module $C_2$-cofiniteness theorem of
Abe--Buhl--Dong~\cite[Proposition~5.2]{ABD} implies
$\dim H'/C_2^U(H')<\infty$.  Since
$C_2^U(H')\subset C_2(H')$, the extension is $C_2$-cofinite.
Complete reducibility for the extension follows from the separable
algebra/module correspondence in Huang--Kirillov--Lepowsky,
Theorems~3.5--3.6.  The unitary invariant form makes the simple
extension self-contragredient.  Together with the already verified
CFT-type grading, these are the strong-rationality hypotheses.
The Lagrangian condition gives
$\Rep(H')\simeq\Rep(U)_{A_3}^0\simeq\Vec$, so it is holomorphic.
Its central charge is eight, so
Dong--Mason~\cite[Theorem~1]{DM} identifies $H'$ with $V_{E_8}$.

For the canonical algebra, the algebra-bimodule duality gives the
second equivalence in \eqref{eq:analytic-extension-defects}; see
M\"uger~\cite[Sections~4 and~8]{Mueger}.  One way to keep track of the opposite
is to view $\HH_3$ as a module category over its center.  The two
commuting actions on it are
\[
 (x\boxtimes y)\cdot m=x\otimes m\otimes y.
\]
Composing two such functors reverses the order of the right factors,
which accounts for ``op'' in the formula.  The double-centralizer
description identifies these endofunctors with the bimodule category
of the canonical algebra.  The equivalence transports with $A_3$.
The first-factor embedding $x\mapsto x\boxtimes\mathbf1$ is fully
faithful and tensor, and supplies the stated $\HH_3$ fusion subcategory.
There are $6\cdot6=36$ simple pairs, with dimensions
$\dim(x)\dim(y)$.  As a consistency check, their total squared
dimension is $\operatorname{Dim}(\HH_3)^2
=\operatorname{Dim}(\Rep(U))$, as required by Morita duality.
\end{proof}

\begin{remark}[Extension choice and preserved symmetry]
\label{rem:analytic-extension-scope}
The selection of the canonical $\HH_3$ algebra proves existence of
$H'$ with the specified all-module category.  It does not identify
$A_3$ with the original embedded $H/U$ algebra; the original boundary
category may be $\HH_2$ or $\HH_3$.  Both holomorphic VOAs being
isomorphic to $V_{E_8}$ does not identify their embeddings of $U$.
Moreover, the defects in \eqref{eq:analytic-extension-defects} preserve
$U$ and need not preserve the enlarged $E_8$ chiral algebra.
The local chiral $H'$-module category is $\Vec$.
\end{remark}

The finite certificates and algebraic arguments supply the stated
representation-theory inputs.  The proved comparisons A2--A3 then give
sector exhaustion, strong unitarity, and a unitary Haagerup-center
identification for the conformal net.  The extra annulus normalization
in A1 is still needed for complete unitarity of the VOA, its specified
unitary ribbon structure, and the chosen $\HH_3$ VOA extension.
In particular, the diagonal net theory and its character interpretation
have a weaker dependency than the chosen unitary VOA extension.
The earlier synthesis
sources are
\artifact{hi_cft_20260923_c8two_final_theorem.md} and
\artifact{hi_cft_20260923_c8two_certificate_audit.md}; their unqualified
analytic conclusion is replaced here by the remaining specific annulus
hypothesis and the proved intermediate reductions.


\section{Relation to the proposed Haagerup theories}\label{sec:comparison}
Albert--Honda--Kaidi--Zheng explicitly assume the existence of a
\(c=8\) Haagerup-center theory in Section II, footnote 4 of
\cite{AHKZ}.
Their equation (12) displays the \(\gamma\)-dependent characters,
and their extension discussion relates the center theory to the
\(E_8\) theory with Haagerup defects.

Under the stated construction inputs, Theorem~\ref{thm:main-net}
supplies a $\gamma=0$ realization as a conformal net and a diagonal
full net theory. Under the additional annulus comparison,
Theorem~\ref{claim:main} supplies the completely unitary VOA
realization and a compatible $E_8$ VOA extension. The distinction
between $U$-preserving and full-$E_8$-preserving defects is maintained.
Neither this construction
nor that finite gauging identifies the continuum limit of the
\(\rho\)-projector chain. Section III of the cited paper proposes
further asymmetric quotients towards charges two and six; those require
additional conformal and defect data. No uniqueness theorem for all
central-charge-eight CFTs is asserted here.

The ambient lattice is also compatible with the conditional
characterization in \cite[Section 4.3]{GR}: a rational chiral algebra
with the Haagerup center at central charge eight must arise from a
rank-two hypergroup action on $E_{6,1}\otimes A_{2,1}$. That
characterization constrains a realization; it does not supply its
generating fields. Conversely, the concrete inclusion constructed here
does not identify the original extension's fusion category as $\HH_3$
without the separate distinction between $\HH_2$ and $\HH_3$ discussed
above. Keeping these two statements separate avoids making an
existence argument depend on a desired choice of Morita representative.

\appendix
\section{Certificate ledger and reproducibility}\label{sec:ledger}
The following frozen records supply the finite construction inputs.
The revision of 27 September 2026 additionally reran the original
provenance check, minimal generation, grade-two exact OPE and PCT
checks, the exact grade-three construction, and the final Gate~4,
charged-span, and Gate~5 verifiers. Fresh outputs were written under
\artifact{c8_revision_20260927/}; the frozen inputs were not overwritten.
All paths are relative to the project directory. The filenames preserve
an unambiguous connection to the actual records.

\begin{longtable}{@{}p{0.25\linewidth}p{0.67\linewidth}@{}}
\toprule
Result & Record and scope\\
\midrule\endfirsthead
\toprule Result & Record and scope\\\midrule\endhead
\bottomrule\endfoot
Minimal generation &
\artifact{hi_cft_20260923_c8_minimal_generators.json}:
the short generating set in \eqref{eq:U} generates all $120$ vectors
of $S_2$. This necessary provenance step was absent from the original
640-file bundle and is explicitly included in the supplemented ledger.\\
Initial spaces &
\artifact{hi_cft_20260923_c8_exact_unitary_certificate.json}
and \artifact{hi_cft_20260923_c8_grade3_exact_v2.json}:
PCT, stress and singular OPE identities, and exact dimension $956$.\\
Lower identities &
\artifact{hi_cft_20260923_next_structure_lower_closure.json}:
3,970 exact relations.\\
Actual triplet &
\artifact{hi_cft_20260923_next_structure_p4_controls.json}:
18 current-zero and six direct Virasoro controls.\\
Raising saturation &
\artifact{hi_cft_20260923_c8two_structure_saturation.json}:
644 exact images in a quotient of dimension 322.\\
Additional direct seeds &
\artifact{hi_cft_20260923_c8two_targeted_h4_trimmed_v2/summary.json}:
118 completed exact identities. The unfinished final block in this
historical directory is not counted as a completed identity.\\
Raised seeds &
\artifact{hi_cft_20260923_c8two_structure_h4_raised_exact/summary.json}:
twelve additional ordinary exact identities.\\
Singlet premise &
\artifact{hi_cft_20260923_c8two_h4_singlet_premise_verified.json}:
all 24 raising conditions replayed.\\
Projected residuals &
\artifact{hi_cft_20260923_c8two_h4_singlet_exact.json}:
twelve full exact residuals, using 21 selected columns.\\
Full Gate 4 &
\artifact{hi_cft_20260923_c8two_gate4_verified.json}:
all 2,898 identities, using singlet recovery and raising saturation.\\
Charged spanning &
\artifact{hi_cft_20260923_c8two_charged_span_verified.json}:
28 reconstructed seeds, rank 322, and current commutators.\\
Full Gate 5 &
\artifact{hi_cft_20260923_c8two_gate5_verified.json}:
all 24 channels by the stated Jacobi argument.\\
Combined provenance &
\artifact{hi_cft_20260923_c8two_proof_bundle.json}:
640 authenticated files, including earlier exact inputs and proof notes.
It checks provenance and completion, not formal correctness of every
mathematical argument.\\
\end{longtable}

The final Gate 4 verifier replays the twelve projected identities using
integer-cleared coefficients, separately from the rational producer.
A field element is represented in the basis
\(1,s,t,st\), with \(s^2=13\), \(t^2=-3\); all residuals are
exact four-component rational identities. Prime calculations are used
to select candidate columns and to certify positive spanning against
an independently known upper dimension. A vanishing prime residual
alone is never treated as a characteristic-zero containment proof.

The provenance-ledger replay can be written to a fresh output without
changing the original certificate:
\begin{quote}\small
\begin{verbatim}
python hi_cft_safe_run.py --limit-mib 1024 -- \
  python hi_cft_20260923_c8two_proof_bundle.py \
  --output /tmp/haagerup-c8-proof-review.json
\end{verbatim}
\end{quote}
These lines represent a single shell command. The output path must not
already exist. This rechecks the frozen source hashes and completion
flags; it does not rerun all original mode calculations. The source
files of the individual certificate producers and verifiers are bound
into the same ledger.

\subsection{Fresh replay coverage}
The supplemented ledger
\artifact{c8_revision_20260927/finite_input_ledger_final.json}
retains the original provenance graph and explicitly binds the
minimal-generation and exact grade-three records. This repairs an
omission in the original bundle's declared roots; it does not modify
the old bundle or infer a mathematical identity from a hash.
The principal fresh outputs are listed below. A ``replay'' means the
specific operation in the right column, rather than a full independent
reconstruction of every earlier lattice calculation.

\begin{longtable}{@{}p{0.38\linewidth}p{0.54\linewidth}@{}}
\toprule
Output under \artifact{c8_revision_20260927/} & Recomputed assertion\\
\midrule\endfirsthead
\toprule Output & Recomputed assertion\\\midrule\endhead
\bottomrule\endfoot
\artifact{provenance_replay.json} & All 640 original source hashes and
required completion flags agree.\\
\artifact{minimal_generators_replay.json} & The specified short
generating set gives the whole exact $120$-dimensional $S_2$.\\
\artifact{unitary_replay.json} & The PCT, stress, primarity and
grade-two singular-product tests, including all eight primary
multiplets, pass in the exact coefficient field.\\
\artifact{grade3_replay.json} & The characteristic-zero generated
span has dimension $956$, with normal subspace of dimension $634$.\\
\artifact{gate4_replay.json} & Integer-cleared projected residuals,
raising premises and saturation establish all $2,898$ required
Gate~4 inputs.\\
\artifact{charged_span_replay.json} & The reconstructed charged
seeds and lowering images span dimension $322$; the saved current
matrices satisfy both axis and cross-axis commutators.\\
\artifact{gate5_replay.json} & Exact quotient relations, nonzero
bottom coefficients and authenticated inputs support all $24$
channels of the printed Jacobi deduction.\\
\artifact{top_bound_replay.json} & Independent rational enumeration
reconstructs all fourteen core tops, trace coefficients, the two
current-character inequalities and the $17/10$ certificate.\\
\artifact{core_extension_replay.json} & The zero-complementary-energy
space in $S_2$ has dimension $25$ and contains all six required
order-thirteen spectral witnesses with their exact eigenvalues.\\
\artifact{top81_replay.json} & In each of four mixed lattice sectors,
exact zero-mode actions and weighted adjoints generate the entire
$81$-dimensional top from a unique Cartan eigenline.\\
\artifact{potts_replay.json} & The explicit Potts primary has the
required stress action and lies in the actual generated $S_3$.\\
\artifact{pf13_membership_replay.json} & The parafermion candidate
lies in the generated Cartan-zero $S_3$, its PCT line is stable,
and the positive stress identity has the stated exact coefficients.\\
\end{longtable}

The last two producers originally had fixed output paths. The adapter
\artifact{c8_revision_20260927/replay_structural_input.py} redirects
only that output construction to a fresh file, requires exactly one
write, and compares the complete mathematical output with the frozen
record, excluding timing and memory statistics. It authenticates the
unchanged producer and checks that the original record remains unchanged.
The parafermion polynomial OPE and decoupling records are authenticated
inputs, not claimed as newly reconstructed by this membership replay.

These calculations ran serially with at most two compute threads and
CPU affinity enforced. The permanent shared-desktop guard imposed
a $6$~GiB address ceiling, an $8$~GiB process-group RSS watchdog and
a $32$~GiB host-available-memory reserve; the exact lattice replays
used smaller address caps of at most $2$~GiB. The grade-three replay
took about $493$ seconds and peaked at $0.190$~GiB sampled group RSS.
No swap capacity was included in the allocation budget. Per-job
estimates and logs are recorded in
\artifact{c8_revision_20260927/resource_review.md}.

\section{Limited formalization}\label{sec:formalization}
Two new Lean files isolate implications that are useful independently
of the VOA implementation. They use the existing Lean~4.33.1 and
Mathlib~4.33.1 environment and were checked without \texttt{sorry} or
new axioms. The axiom reports contain only Lean's standard
propositional extensionality, quotient soundness and, for some
statements, classical choice.

\begin{enumerate}
\item \artifact{lean/HaagerupIzumi/CFT/FiniteClosure.lean} proves
raising saturation for a stable submodule, under an explicit joint
nilpotence hypothesis. It also proves recovery from a projected
residual and the abstract linear Jacobi-bottom-witness implication.
The hypotheses identify exactly where the current-representation
bounds and finite identities enter the manuscript's use of these
lemmas.
\item \artifact{lean/HaagerupIzumi/CFT/TopBound.lean} proves the
displayed fourteen-variable deficit identity and derives the
cleared-denominator $17/10$ inequality from nonnegative
multiplicities and the two stated $\mathfrak{sl}_2$ inequalities.
The independent rational checker separately reconstructs the
table to which this identity is applied.
\end{enumerate}

The safe compiler commands, issued serially from the project root, are
\begin{quote}\small
\begin{verbatim}
python lean/safe_lean.py \
  lean/HaagerupIzumi/CFT/FiniteClosure.lean
python lean/safe_lean.py \
  lean/HaagerupIzumi/CFT/TopBound.lean
\end{verbatim}
\end{quote}
The driver invokes the same resource guard. These formalizations do
not define a VOA, validate the lattice engine, prove filtered Wick
reduction, construct the embedded parafermion factor, or resolve A1.
Their value is to check the stated algebraic implications with their
premises exposed, rather than to certify the entire manuscript.

\section{Compilation and review scope}\label{sec:compilation}
The main source is \artifact{hi_cft_c8_proof.tex}; its companion sources
are \artifact{hi_cft_c8_manuscript_structure.tex} and
\artifact{hi_cft_c8_manuscript_analytic.tex}, with references in
\artifact{hi_cft_c8_bibliography.tex}. Run \texttt{pdflatex} twice
on the main source for cross-references and the table of contents.
The document uses standard TeX packages and requires no shell escape.

The frozen computational records remain unchanged. Fresh replays,
the supplemented provenance ledger and the two Lean modules are
additional checks with the scopes stated above. The original ledger
does not authenticate this revised exposition. The remaining analytic
comparison and the exact embedded-field inputs should be examined
directly when assessing the conditional VOA conclusion; neither a
successful PDF build nor a modular-character match resolves them.
\begin{thebibliography}{99}
\raggedright

\bibitem{DSK}
A.~De Sole and V.~G.~Kac,
\emph{Freely generated vertex algebras and non-linear Lie conformal algebras},
Comm. Math. Phys. \textbf{254} (2005), 659--694.
\href{https://arxiv.org/abs/math-ph/0312042}{arXiv:math-ph/0312042}.

\bibitem{EG}
D.~E.~Evans and T.~Gannon,
\emph{The exoticness and realisability of twisted Haagerup--Izumi modular data},
Comm. Math. Phys. \textbf{307} (2011), 463--512.
\href{https://arxiv.org/abs/1006.1326}{arXiv:1006.1326}.

\bibitem{CKLW}
S.~Carpi, Y.~Kawahigashi, R.~Longo, and M.~Weiner,
\emph{From vertex operator algebras to conformal nets and back},
Mem. Amer. Math. Soc. \textbf{254} (2018), no.~1213.
\href{https://arxiv.org/abs/1503.01260v4}{arXiv:1503.01260v4}.

\bibitem{AHKZ}
J.~Albert, Y.~Honda, J.~Kaidi, and Y.~Zheng,
\emph{Haagerup symmetry in $(E_8)_1$?}, preprint (2025), revised 2026.
\href{https://arxiv.org/abs/2512.08225v2}{arXiv:2512.08225v2}.

\bibitem{GR}
T.~Gannon and B.~C.~Rayhaun,
\emph{Hypergroup symmetry in relative quantum field theories and chiral algebras},
preprint (2026).
\href{https://arxiv.org/abs/2606.05279}{arXiv:2606.05279}.

\bibitem{ABD}
T.~Abe, G.~Buhl, and C.~Dong,
\emph{Rationality, regularity, and $C_2$-cofiniteness},
Trans. Amer. Math. Soc. \textbf{356} (2004), 3391--3402.
\href{https://doi.org/10.1090/S0002-9947-03-03413-5}{doi:10.1090/S0002-9947-03-03413-5};
\href{https://arxiv.org/abs/math/0204021}{arXiv:math/0204021}.

\bibitem{EP}
P.~Etingof and D.~Penneys,
\emph{Rigidity of non-negligible objects of moderate growth in braided categories},
\href{https://arxiv.org/abs/2412.17681v3}{arXiv:2412.17681v3}.

\bibitem{McRae}
R.~McRae,
\emph{On rationality for $C_2$-cofinite vertex operator algebras},
\href{https://arxiv.org/abs/2108.01898}{arXiv:2108.01898}.

\bibitem{ALY}
T.~Arakawa, C.~H.~Lam, and H.~Yamada,
\emph{Parafermion vertex operator algebras and $W$-algebras},
\href{https://arxiv.org/abs/1701.06229}{arXiv:1701.06229}.

\bibitem{BJ}
M.~Bischoff and C.~Jones,
\emph{Computing fusion rules for spherical $G$-extensions of fusion categories},
\href{https://arxiv.org/abs/1909.02816}{arXiv:1909.02816}.

\bibitem{RLYG}
M.~Reeder, P.~Levy, J.-K.~Yu, and B.~H.~Gross,
\emph{Gradings of positive rank on simple Lie algebras},
\href{https://arxiv.org/abs/1307.5765}{arXiv:1307.5765}.

\bibitem{Larson}
H.~Larson,
\emph{Pseudo-unitary non-self-dual fusion categories of rank $4$},
\href{https://arxiv.org/abs/1401.1879}{arXiv:1401.1879}.

\bibitem{EGO}
P.~Etingof, S.~Gelaki, and V.~Ostrik,
\emph{Classification of fusion categories of dimension $pq$},
\href{https://arxiv.org/abs/math/0304194}{arXiv:math/0304194}.

\bibitem{KO}
A.~Kirillov, Jr., and V.~Ostrik,
\emph{On a $q$-analogue of the McKay correspondence and the ADE classification
of $\widehat{\mathfrak{sl}}_2$ conformal field theories},
\href{https://arxiv.org/abs/math/0101219}{arXiv:math/0101219}.

\bibitem{AGJS}
Md.~Abhishek, S.~Grover, D.~P.~Jatkar, and K.~Singh,
\emph{Looking for the $G_2$ Higgs branch of 4D rank 1 SCFTs},
\href{https://arxiv.org/abs/2312.00275}{arXiv:2312.00275}.

\bibitem{DLN}
C.~Dong, X.~Lin, and S.-H.~Ng,
\emph{Congruence property in conformal field theory},
\href{https://msp.org/ant/2015/9-9/ant-v9-n9-p03-s.pdf}{Algebra Number Theory
\textbf{9} (2015), no.~9}.

\bibitem{Yu}
Z.~Yu,
\emph{On the center of near-group fusion category of type $\mathbb Z_3+6$},
\href{https://arxiv.org/abs/2312.12195v4}{arXiv:2312.12195v4}.

\bibitem{Ostrik13}
V.~Ostrik,
\emph{Pivotal fusion categories of rank $3$},
with an appendix jointly with D.~Nikshych,
\href{https://arxiv.org/abs/1309.4822}{arXiv:1309.4822}.

\bibitem{NS}
S.-H.~Ng and P.~Schauenburg,
\emph{Frobenius--Schur indicators and exponents of spherical categories},
\href{https://arxiv.org/abs/math/0601012}{arXiv:math/0601012}.

\bibitem{GSY}
T.~Gannon, A.~Schopieray, and H.~Yadav,
\emph{On Haagerup--Izumi fusion categories},
\href{https://arxiv.org/abs/2609.25185v1}{arXiv:2609.25185v1}.

\bibitem{TenFields}
J.~E.~Tener,
\emph{Representation theory in chiral conformal field theory:
from fields to observables},
Selecta Math. (N.S.) \textbf{25} (2019), article 76.
\href{https://doi.org/10.1007/s00029-019-0526-3}{doi:10.1007/s00029-019-0526-3};
\href{https://arxiv.org/abs/1810.08168v2}{arXiv:1810.08168v2}.

\bibitem{TenFusion}
J.~E.~Tener,
\emph{Fusion and positivity in chiral conformal field theory},
Geom. Funct. Anal. \textbf{34} (2024), 1226--1296.
\href{https://doi.org/10.1007/s00039-024-00685-8}{doi:10.1007/s00039-024-00685-8};
\href{https://arxiv.org/abs/1910.08257v3}{arXiv:1910.08257v3}.

\bibitem{HTNet}
A.~G.~Henriques and J.~E.~Tener,
\emph{Every conformal net has an associated unitary VOA},
preprint (2025),
\href{https://arxiv.org/abs/2507.20735v3}{arXiv:2507.20735v3}.

\bibitem{HTAnn}
A.~G.~Henriques and J.~E.~Tener,
\emph{Integrating positive energy representations of the Virasoro algebra},
preprint (2025),
\href{https://arxiv.org/abs/2506.08684}{arXiv:2506.08684}.

\bibitem{Longo}
R.~Longo,
\emph{Conformal subnets and intermediate subfactors},
Comm. Math. Phys. \textbf{237} (2003), 7--30.
\href{https://doi.org/10.1007/s00220-003-0814-8}{doi:10.1007/s00220-003-0814-8};
\href{https://arxiv.org/abs/math/0102196}{arXiv:math/0102196}.

\bibitem{KL}
Y.~Kawahigashi and R.~Longo,
\emph{Classification of local conformal nets. Case $c<1$},
Ann. of Math. (2) \textbf{160} (2004), 493--522.
\href{https://annals.math.princeton.edu/wp-content/uploads/annals-v160-n2-p04.pdf}{Published article}.

\bibitem{Gui}
B.~Gui,
\emph{Q-systems and extensions of completely unitary vertex operator algebras},
Int. Math. Res. Not. IMRN \textbf{2022}, no.~10, 7550--7614.
\href{https://doi.org/10.1093/imrn/rnaa300}{doi:10.1093/imrn/rnaa300};
\href{https://arxiv.org/abs/1908.03484v3}{arXiv:1908.03484v3}.

\bibitem{Izumi}
M.~Izumi,
\emph{The classification of $3^n$ subfactors and related fusion categories},
Quantum Topol. \textbf{9} (2018), 473--562.
\href{https://doi.org/10.4171/QT/113}{doi:10.4171/QT/113};
\href{https://arxiv.org/abs/1609.07604v2}{arXiv:1609.07604v2}.

\bibitem{GS}
P.~Grossman and N.~Snyder,
\emph{Quantum subgroups of the Haagerup fusion categories},
Comm. Math. Phys. \textbf{311} (2012), 617--643.
\href{https://doi.org/10.1007/s00220-012-1427-x}{doi:10.1007/s00220-012-1427-x};
\href{https://arxiv.org/abs/1102.2631v2}{arXiv:1102.2631v2}.

\bibitem{Reutter}
D.~Reutter,
\emph{Uniqueness of unitary structure for unitarizable fusion categories},
Comm. Math. Phys. \textbf{397} (2023), 37--52.
\href{https://doi.org/10.1007/s00220-022-04425-7}{doi:10.1007/s00220-022-04425-7};
\href{https://arxiv.org/abs/1906.09710}{arXiv:1906.09710}.

\bibitem{Rehren}
K.-H.~Rehren,
\emph{Canonical tensor product subfactors},
Comm. Math. Phys. \textbf{211} (2000), 395--406.
\href{https://doi.org/10.1007/s002200050818}{doi:10.1007/s002200050818};
\href{https://arxiv.org/abs/math/9911149v2}{arXiv:math/9911149v2}.

\bibitem{DMNO}
A.~Davydov, M.~M\"uger, D.~Nikshych, and V.~Ostrik,
\emph{The Witt group of non-degenerate braided fusion categories},
J. Reine Angew. Math. \textbf{677} (2013), 135--177.
\href{https://arxiv.org/abs/1009.2117v2}{arXiv:1009.2117v2}.

\bibitem{HKL}
Y.-Z.~Huang, A.~Kirillov, Jr., and J.~Lepowsky,
\emph{Braided tensor categories and extensions of vertex operator algebras},
Comm. Math. Phys. \textbf{337} (2015), 1143--1159.
\href{https://doi.org/10.1007/s00220-015-2292-1}{doi:10.1007/s00220-015-2292-1};
\href{https://arxiv.org/abs/1406.3420}{arXiv:1406.3420}.

\bibitem{DM}
C.~Dong and G.~Mason,
\emph{Holomorphic vertex operator algebras of small central charge},
Pacific J. Math. \textbf{213} (2004), 253--266.
\href{https://doi.org/10.2140/pjm.2004.213.253}{doi:10.2140/pjm.2004.213.253};
\href{https://arxiv.org/abs/math/0203005}{arXiv:math/0203005}.

\bibitem{Mueger}
M.~M\"uger,
\emph{From subfactors to categories and topology II.
The quantum double of tensor categories and subfactors},
J. Pure Appl. Algebra \textbf{180} (2003), 159--219.
\href{https://arxiv.org/abs/math/0111205}{arXiv:math/0111205}.

\bibitem{Vercleyen}
G.~Vercleyen,
\emph{On Low-Rank Multiplicity-Free Fusion Categories},
Ph.D. thesis, Maynooth University, 2024.
\href{https://mural.maynoothuniversity.ie/id/eprint/18958/1/phd.pdf}{Institutional thesis PDF}.

\bibitem{EG15}
D.~E.~Evans and T.~Gannon,
\emph{Non-unitary fusion categories and their doubles via endomorphisms},
\href{https://arxiv.org/abs/1506.03546}{arXiv:1506.03546}.

\bibitem{LRNets}
R.~Longo and K.-H.~Rehren,
\emph{Nets of subfactors},
Rev. Math. Phys. \textbf{7} (1995), 567--597.
\href{https://doi.org/10.1142/S0129055X95000232}{doi:10.1142/S0129055X95000232};
\href{https://arxiv.org/abs/hep-th/9411077v2}{arXiv:hep-th/9411077v2}.

\bibitem{BisTY}
M.~Bischoff,
\emph{Conformal net realizability of Tambara--Yamagami categories
and generalized metaplectic modular categories},
preprint (2018),
\href{https://arxiv.org/abs/1803.04949}{arXiv:1803.04949}.

\bibitem{Linshaw}
A.~R.~Linshaw,
\emph{Universal two-parameter $\mathcal W_\infty$-algebra and vertex
algebras of type $\mathcal W(2,3,\ldots,N)$},
Compos. Math. \textbf{157} (2021), no.~1, 12--82.
\href{https://doi.org/10.1112/S0010437X20007514}{doi:10.1112/S0010437X20007514};
\href{https://arxiv.org/abs/1710.02275}{arXiv:1710.02275}.

\bibitem{DLY}
C.~Dong, C.~H.~Lam, and H.~Yamada,
\emph{$W$-algebras related to parafermion algebras},
J. Algebra \textbf{322} (2009), no.~7, 2366--2403.
\href{https://arxiv.org/abs/0809.3630}{arXiv:0809.3630}.

\bibitem{GuiCompare}
B.~Gui,
\emph{Comparison of extensions of unitary vertex operator algebras
and conformal nets},
preprint (2025), version~3, 18 February 2026,
\href{https://arxiv.org/abs/2505.03235v3}{arXiv:2505.03235v3}.

\bibitem{HTGenusZero}
A.~G.~Henriques and J.~E.~Tener,
\emph{A genus zero functorial CFT associated to the vacuum sector
of a conformal net},
preprint (2026), version~1, 25 August 2026,
\href{https://arxiv.org/abs/2608.24437v1}{arXiv:2608.24437v1}.

\end{thebibliography}

\end{document}
